\chapter{Problems Related to Volume I}
\label{ch:companion-i}

\paragraph{Main-text correspondence.}
This Part accompanies \emph{Volume I: Linear Algebra}. Problems are grouped by
mathematical theme rather than by reproducing the chapter sequence of the main text.
Each section gives the corresponding main-text chapter range for further reading.

\section{Vector Spaces}

\paragraph{Related material.} Volume I, Chapters \texttt{I/01--I/06}.

\begin{problem}[CP-I-0002]
\label{prob:cp-i-0002}
\textbf{Level Sets of a Cubic Surface}

\hspace*{1.5em}Consider the subset
	\[
	\Sigma \subset \mathbb{R}^3
	\qquad\text{defined by}\qquad
	\Sigma = \{(x,y,z) : z = x^3 + y^3 - 3xy\}.
	\]

		\par\noindent\textbullet\quad
		Show that $\Sigma$ is a smooth surface in $\mathbb{R}^3$.

		\par\noindent\textbullet\quad
		Show that the level set
		\[
		\Sigma \cap \{z=c\}
		\]
		is a (not necessarily connected) smooth curve in the $(x,y)$–plane
		unless
		\[
		c \in \{-1,0\}.
		\]

		\par\noindent\textbullet\quad
		Sketch the level sets for the special values $c=-1$ and $c=0$.

		(For $c=-1$, you may use the factorization)
		\[
		x^3 + y^3 - 3xy + 1
		= (x+y+1)\bigl(x^2 + y^2 - xy - x - y + 1\bigr).
		\]

	\textbf{Worked solution.}

	Define
	\[
	f(x,y) = x^3 + y^3 - 3xy, \qquad
	F(x,y,z) = z - f(x,y).
	\]
	Then
	\[
	\Sigma = F^{-1}(0).
	\]
	We have
	\[
	\frac{\partial F}{\partial z}(x,y,z) = 1 \neq 0
	\]
	for all $(x,y,z)$.
	By the implicit function theorem, the level set $F^{-1}(0)$ is a smooth
	$2$--dimensional submanifold of $\mathbb{R}^3$; in fact it is the graph of the
	smooth function $f$.  Hence $\Sigma$ is a smooth surface.

	\paragraph{(b) Level sets are smooth curves for $c\notin\{-1,0\}$.}
	Fix $c\in\mathbb{R}$ and intersect $\Sigma$ with the plane $z=c$.
	On this plane, points have the form $(x,y,c)$ and satisfy
	\[
	c = x^3 + y^3 - 3xy,
	\]
	or equivalently
	\[
	g_c(x,y) := x^3 + y^3 - 3xy - c = 0.
	\]
	Thus
	\[
	\Sigma \cap \{z=c\}
	= \{(x,y,c) : g_c(x,y)=0\}
	\]
	is naturally identified with the level set $g_c^{-1}(0)\subset\mathbb{R}^2$.

	We compute
	\[
	\nabla g_c(x,y) = \bigl(3x^2 - 3y,\; 3y^2 - 3x\bigr),
	\]
	which is independent of $c$.
	The gradient vanishes exactly when
	\[
	3x^2 - 3y = 0, \qquad 3y^2 - 3x = 0
	\quad\Longleftrightarrow\quad
	x^2 = y, \quad y^2 = x.
	\]
	Substituting $y=x^2$ into $y^2=x$ gives
	\[
	x^4 = x \quad\Longleftrightarrow\quad x(x^3-1)=0,
	\]
	so $x=0$ or $x=1$.  Thus the only critical points of $g_c$ are
	\[
	(0,0) \quad\text{and}\quad (1,1).
	\]

	We now determine for which $c$ these points lie on the level set $g_c^{-1}(0)$:

		\par\noindent\textbullet\quad At $(0,0)$ we have $x^3 + y^3 - 3xy = 0$, so
		\[
		g_c(0,0) = -c,
		\]
		which vanishes iff $c=0$.
		\par\noindent\textbullet\quad At $(1,1)$ we have $1^3+1^3-3\cdot1\cdot1 = -1$, so
		\[
		g_c(1,1) = -1 - c,
		\]
		which vanishes iff $c=-1$.

	Hence, if $c\notin\{-1,0\}$, the level set $g_c^{-1}(0)$ contains no critical
	points of $g_c$, so $\nabla g_c\neq 0$ at every point on $g_c^{-1}(0)$.
	By the implicit function theorem in $\mathbb{R}^2$, each such level set is a
	(possibly disconnected) smooth $1$--dimensional submanifold of
	$\mathbb{R}^2$, i.e.\ a smooth plane curve.  Therefore
	$\Sigma\cap\{z=c\}$ is a smooth curve in the plane for all
	$c\notin\{-1,0\}$.

	For $c=-1$ the equation becomes
	\[
	x^3 + y^3 - 3xy + 1 = 0.
	\]
	The given factorisation
	\[
	x^3 + y^3 - 3xy + 1
	= (x+y+1)\bigl(x^2 + y^2 - xy - x - y + 1\bigr)
	\]
	shows that the level set is the union of
	\[
	x + y + 1 = 0
	\]
	(a straight line) and
	\[
	x^2 + y^2 - xy - x - y + 1 = 0,
	\]
	which is an ellipse (the associated quadratic form is positive definite).
	Thus the intersection $\Sigma\cap\{z=-1\}$ consists of a line and a smooth
	oval.

	For $c=0$ the equation is
	\[
	x^3 + y^3 - 3xy = 0.
	\]
	The origin $(0,0)$ lies on this level set and is a critical point of $g_0$
	since $\nabla g_0(0,0)=(0,0)$.  Thus the level set $g_0^{-1}(0)$ has a
	singularity at the origin, where several branches of the cubic curve meet.
	Away from $(0,0)$, $\nabla g_0\neq 0$, so the curve is smooth at all other
	points.  Hence $\Sigma\cap\{z=0\}$ is a singular cubic curve passing through
	$(0,0)$, smooth everywhere except at the origin.

	For $c=-1$ we have
	\[
	x^3 + y^3 - 3xy + 1 = 0,
	\]
	and the given factorisation
	\[
	x^3 + y^3 - 3xy + 1
	= (x+y+1)\bigl(x^2 + y^2 - xy - x - y + 1\bigr)
	\]
	shows that the level set consists of two components:
	the straight line $x+y+1=0$ and the conic
	\[
	x^2 + y^2 - xy - x - y + 1 = 0,
	\]
	which is an ellipse (since the associated quadratic form is positive
	definite).  Thus $\Sigma\cap\{z=-1\}$ is the union of a line and a smooth
	oval lying in the $(x,y)$–plane.

For $c=0$ the level set is given by
\[
x^3 + y^3 - 3xy = 0.
\]
The origin $(0,0)$ lies on this curve and is a critical point, since
\[
\nabla(x^3 + y^3 - 3xy)(0,0) = (0,0).
\]
Thus the origin is a singular point, where several smooth branches of the
cubic meet.  Away from $(0,0)$ the gradient is non-zero, so the curve is a
smooth $1$--dimensional submanifold of the plane except at this single
singular point.
\end{problem}

\begin{problem}[CP-I-0014]
\label{prob:cp-i-0014}
\textbf{Images and kernels of two maps on $\mathbb R^3$}

\textbf{Exact statement.}
\hspace*{1.5em}Give a geometrical description of the images and kernels of each of the following linear maps on $\mathbb R^3$:
\[
\text{(a)}\quad T:(x,y,z)\mapsto (x+2y+z,\ x+2y+z,\ 2x+4y+2z),
\]
\[
\text{(b)}\quad S:(x,y,z)\mapsto (x+2y+3z,\ x-y+z,\ x+5y+5z).
\]

\end{problem}

\begin{solution}
(a) Put $s=x+2y+z$. Then
\[
T(x,y,z)=(s,s,2s)=s(1,1,2).
\]
Hence $\mathrm{Im}\,T=\mathrm{span}\{(1,1,2)\}$ (a line through the origin) and
\[
\ker T=\{(x,y,z):x+2y+z=0\}
\]
(a plane through the origin).

(b) The matrix of $S$ is
\[
A=\begin{pmatrix}1&2&3\\ 1&-1&1\\ 1&5&5\end{pmatrix},
\]
which has rank $2$, so $\dim\ker S=1$ and $\dim\mathrm{Im}\,S=2$.
Solving $A(x,y,z)^T=0$ gives a kernel direction $(-5,-2,3)$, hence
\[
\ker S=\mathrm{span}\{(-5,-2,3)\}.
\]
A basis for the image is given by independent columns, e.g.
\[
(1,1,1),\quad (2,-1,5),
\]
so $\mathrm{Im}\,S=\mathrm{span}\{(1,1,1),(2,-1,5)\}$, a plane through the origin.
A normal vector is $(1,1,1)\times(2,-1,5)=(2,-1,-1)$, so equivalently
\[
\mathrm{Im}\,S=\{(X,Y,Z):2X-Y-Z=0\}.
\]
\end{solution}

\begin{problem}[CP-I-0016]
\label{prob:cp-i-0016}
\textbf{Dirichlet Green Function for \(-u''+u\) on \([0,1]\)}

\textbf{Canonical authored reconstruction from the retained Green-kernel fragment.}

\hspace*{1.5em}Consider the boundary-value operator
\[
L u=-u''+u,\qquad u(0)=u(1)=0.
\]
Let
\[
u_{<}(x)=\sinh x,\qquad u_{>}(x)=\sinh(1-x).
\]
Show that the Green function for \(L\) is
\[
G(x,\xi)=\frac{1}{\sinh 1}
\begin{cases}
\sinh x\,\sinh(1-\xi),&x\le \xi,\\[2mm]
\sinh \xi\,\sinh(1-x),&x\ge \xi.
\end{cases}
\]
Verify the boundary conditions, continuity at \(x=\xi\), and the derivative jump required by
\[
(-\partial_x^2+1)G(\,\cdot\,,\xi)=\delta_\xi .
\]
\end{problem}

\begin{solution}
For \(x\neq \xi\), each branch is a product of a constant in \(x\) with a solution of
\[
-u''+u=0,
\]
so \(LG=0\) away from \(x=\xi\).

The boundary conditions are immediate:
\[
G(0,\xi)=0,\qquad G(1,\xi)=0.
\]
At \(x=\xi\), both branches have the same value
\[
G(\xi,\xi)=\frac{\sinh \xi\,\sinh(1-\xi)}{\sinh 1},
\]
so \(G\) is continuous.

For the one-sided derivatives,
\[
\partial_xG(\xi^{-},\xi)
=\frac{\cosh \xi\,\sinh(1-\xi)}{\sinh 1},
\]
while
\[
\partial_xG(\xi^{+},\xi)
=-\frac{\sinh \xi\,\cosh(1-\xi)}{\sinh 1}.
\]
Hence
\[
\partial_xG(\xi^{+},\xi)-\partial_xG(\xi^{-},\xi)
=
-\frac{\sinh \xi\,\cosh(1-\xi)+\cosh \xi\,\sinh(1-\xi)}
{\sinh 1}
=-1,
\]
because
\[
\sinh \xi\,\cosh(1-\xi)+\cosh \xi\,\sinh(1-\xi)=\sinh 1.
\]
Thus the jump in the first derivative is \(-1\), exactly the distributional jump required by
\[
-\partial_x^2G=\delta_\xi+\text{regular terms}.
\]
Since the regular terms are canceled by \(+G\), we obtain
\[
(-\partial_x^2+1)G(\,\cdot\,,\xi)=\delta_\xi.
\]
\end{solution}

\begin{problem}[CP-I-0018]
\label{prob:cp-i-0018}
\textbf{Sum of linear maps; images/kernels; idempotents}

\textbf{Exact statement.}

	\par\hspace*{1.5em} (i)\quad If $\alpha$ and $\beta$ are linear maps from $U$ to $V$ show that $\alpha+\beta$ is linear.
	Give explicit counter-examples to:
	\[
	\text{(a) }\operatorname{im}(\alpha+\beta)=\operatorname{im}(\alpha)+\operatorname{im}(\beta),\qquad
	\text{(b) }\ker(\alpha+\beta)=\ker(\alpha)\cap\ker(\beta).
	\]
	Show that in general each equality can be replaced by a valid inclusion of one side in the other.
	\par\noindent (ii)\quad Let $\alpha:V\to V$ be linear. Show that if $\alpha^2=\alpha$ then
	$V=\ker(\alpha)\oplus \operatorname{im}(\alpha)$.
	Does your proof still work if $V$ is infinite dimensional? Is the result still true?
\end{problem}

\begin{solution}
\textbf{(i)} For $u,u'\in U$ and $\lambda\in\mathbb R$,
\[
(\alpha+\beta)(u+u')=\alpha(u+u')+\beta(u+u')
=(\alpha u+\alpha u')+(\beta u+\beta u')
=(\alpha+\beta)u+(\alpha+\beta)u',
\]
\[
(\alpha+\beta)(\lambda u)=\alpha(\lambda u)+\beta(\lambda u)=\lambda\alpha(u)+\lambda\beta(u)
=\lambda(\alpha+\beta)(u),
\]
so $\alpha+\beta$ is linear.

Counterexample to (a): take $U=V=\mathbb R$, $\alpha=\mathrm{id}$, $\beta=-\mathrm{id}$.
Then $\alpha+\beta=0$, hence $\operatorname{im}(\alpha+\beta)=\{0\}$ but
$\operatorname{im}(\alpha)+\operatorname{im}(\beta)=\mathbb R$.
In general,
\[
\operatorname{im}(\alpha+\beta)\subseteq \operatorname{im}(\alpha)+\operatorname{im}(\beta),
\]
since $(\alpha+\beta)(u)=\alpha(u)+\beta(u)$.

Counterexample to (b): with the same $\alpha,\beta$, we have $\ker(\alpha+\beta)=\mathbb R$ but
$\ker(\alpha)\cap\ker(\beta)=\{0\}$.
In general,
\[
\ker(\alpha)\cap\ker(\beta)\subseteq \ker(\alpha+\beta),
\]
since $\alpha(u)=\beta(u)=0\Rightarrow (\alpha+\beta)(u)=0$.

\medskip
\textbf{(ii)} Assume $\alpha^2=\alpha$.
For any $x\in V$,
\[
x=(x-\alpha x)+\alpha x.
\]
Moreover,
\[
\alpha(x-\alpha x)=\alpha x-\alpha^2 x=\alpha x-\alpha x=0,
\]
so $x-\alpha x\in\ker(\alpha)$ and $\alpha x\in\operatorname{im}(\alpha)$, hence
$V=\ker(\alpha)+\operatorname{im}(\alpha)$.

If $y\in\ker(\alpha)\cap\operatorname{im}(\alpha)$, then $y=\alpha(x)$ for some $x$ and
$0=\alpha(y)=\alpha(\alpha x)=\alpha^2 x=\alpha x=y$, so $y=0$. Thus the sum is direct:
\[
V=\ker(\alpha)\oplus\operatorname{im}(\alpha).
\]
No finite-dimensional hypothesis was used, so the proof works for infinite-dimensional $V$ and the result remains true.

\bigskip
\end{solution}

\begin{problem}[CP-I-0019]
\label{prob:cp-i-0019}
\textbf{Distance bounds from a disk (triangle inequalities)}

\textbf{Exact statement.}
\hspace*{1.5em}Let \(D\) be the interior of the circle \(|z-1-i|=1\). Show, by using suitable inequalities for \(|z_1\pm z_2|\), that
if \(z\in D\) then
\[
\sqrt5-1<|z-3|<\sqrt5+1.
\]
Obtain the same result geometrically \([\,\)start by considering the line through the centre of the circle and the point \(3\,]\).

\end{problem}

\begin{solution}
Let \(c=1+i\). Then \(|z-c|<1\) for \(z\in D\). Write
\[
z-3=(z-c)+(c-3).
\]
By the triangle inequality and reverse triangle inequality,
\[
|z-3|\le |z-c|+|c-3|,\qquad |z-3|\ge \bigl||c-3|-|z-c|\bigr|.
\]
Now
\[
|c-3|=|(1+i)-3|=|-2+i|=\sqrt{4+1}=\sqrt5.
\]
Hence, since \(|z-c|<1\),
\[
|z-3|<1+\sqrt5,\qquad |z-3|>\sqrt5-1,
\]
which gives \(\sqrt5-1<|z-3|<\sqrt5+1\).

Geometrically: \(D\) is the open disk of radius \(1\) centered at \(c\). The distance from the fixed point \(3\) to points in that disk ranges strictly between \(|3-c|-1=\sqrt5-1\) and \(|3-c|+1=\sqrt5+1\).
\end{solution}

\begin{problem}[CP-I-0026]
\label{prob:cp-i-0026}
\textbf{Common complementary subspace}

\textbf{Exact statement.}
\hspace*{1.5em}Show that any two subspaces of the same dimension in a finite dimensional vector space have a common complementary subspace.
\end{problem}

\begin{solution}
Let $U,W\le V$ with $\dim U=\dim W=k$ and $\dim V=n$. Set $I=U\cap W$ and $\dim I=r$.
Choose bases
\[
I=\mathrm{span}(a_1,\dots,a_r),
\]
extend to bases
\[
U=\mathrm{span}(a_1,\dots,a_r,b_1,\dots,b_{k-r}),\qquad
W=\mathrm{span}(a_1,\dots,a_r,c_1,\dots,c_{k-r}).
\]
Define
\[
D:=\mathrm{span}\{c_i-b_i:\ i=1,\dots,k-r\}\subseteq U+W.
\]
If $\sum t_i(c_i-b_i)\in U$, then $\sum t_i c_i\in U\cap W=I$, which forces $t_i=0$ since
$c_1,\dots,c_{k-r}$ are independent modulo $I$. Hence $D\cap U=\{0\}$. Similarly $D\cap W=\{0\}$.
Therefore
\[
U\oplus D=U+W,\qquad W\oplus D=U+W.
\]
Choose a complement $Z_0$ of $U+W$ in $V$ (possible since $V$ is finite dimensional):
\[
V=(U+W)\oplus Z_0.
\]
Set $Z:=D\oplus Z_0$. Then $\dim Z=n-k$ and
\[
V=U\oplus Z\qquad\text{and}\qquad V=W\oplus Z,
\]
so $Z$ is a common complementary subspace.
\end{solution}

\begin{problem}[CP-I-0033]
\label{prob:cp-i-0033}
\textbf{Fifth roots of unity factorization and a trig identity}

\textbf{Exact statement.}
\hspace*{1.5em}Express
\[
I=\frac{z^5-1}{z-1}
\]
as a polynomial in \(z\). By considering the complex fifth root of unity \(\omega\), obtain the four factors of \(I\) linear in \(z\).
Hence write \(I\) as the product of two real quadratic factors. By considering the term in \(z^2\) in the identity so obtained for \(I\),
show that
\[
4\cos\frac{\pi}{5}\,\sin\frac{\pi}{10}=1.
\]

\end{problem}

\begin{solution}
For \(z\neq1\),
\[
\frac{z^5-1}{z-1}=z^4+z^3+z^2+z+1,
\]
so as a polynomial,
\[
I=z^4+z^3+z^2+z+1.
\]
Let \(\omega=e^{2\pi i/5}\). Then
\[
z^5-1=(z-1)(z-\omega)(z-\omega^2)(z-\omega^3)(z-\omega^4),
\]
hence
\[
I=(z-\omega)(z-\omega^2)(z-\omega^3)(z-\omega^4).
\]
Pair conjugates \(\omega^4=\overline{\omega}\) and \(\omega^3=\overline{\omega^2}\):
\[
(z-\omega)(z-\omega^4)=z^2-(\omega+\omega^4)z+1
= z^2-2\cos\!\left(\frac{2\pi}{5}\right)z+1,
\]
\[
(z-\omega^2)(z-\omega^3)=z^2-(\omega^2+\omega^3)z+1
= z^2-2\cos\!\left(\frac{4\pi}{5}\right)z+1.
\]
Therefore
\[
I=\Bigl(z^2-2\cos\!\left(\frac{2\pi}{5}\right)z+1\Bigr)
\Bigl(z^2-2\cos\!\left(\frac{4\pi}{5}\right)z+1\Bigr).
\]
Expanding \((z^2-az+1)(z^2-bz+1)\) gives the \(z^2\)-coefficient \(ab+2\). Comparing with
\(I=z^4+z^3+z^2+z+1\) yields
\[
ab+2=1\quad\Longrightarrow\quad ab=-1,
\]
with \(a=2\cos(2\pi/5)\), \(b=2\cos(4\pi/5)\). Hence
\[
4\cos\!\left(\frac{2\pi}{5}\right)\cos\!\left(\frac{4\pi}{5}\right)=-1.
\]
Using \(\cos(4\pi/5)=-\cos(\pi/5)\) and \(\cos(2\pi/5)=\sin(\pi/10)\), we get
\[
4\cos\!\left(\frac{\pi}{5}\right)\sin\!\left(\frac{\pi}{10}\right)=1.
\]
\end{solution}

\begin{problem}[CP-I-0035]
\label{prob:cp-i-0035}
\textbf{A Subspace Defined by a Four-Term Recurrence}

\textbf{Canonical authored reconstruction from the retained recurrence fragment.}

\hspace*{1.5em}For \(n\ge 4\), let
\[
V=\left\{x=(x_1,\dots,x_n)\in\mathbb R^n:
x_i+x_{i+1}+x_{i+2}+x_{i+3}=0
\text{ for }1\le i\le n-3\right\}.
\]
Show that \(V\) is a subspace and that every vector in \(V\) is uniquely determined by
\((x_1,x_2,x_3)\). Deduce that
\[
\dim V=3.
\]
Show moreover that every solution satisfies \(x_{i+4}=x_i\) whenever both sides are defined.
\end{problem}

\begin{solution}
The defining equations are homogeneous linear equations, so their common solution set is a
subspace of \(\mathbb R^n\).

Each equation gives
\[
x_{i+3}=-(x_i+x_{i+1}+x_{i+2}).
\]
Thus after choosing \(x_1,x_2,x_3\), all later coordinates are determined recursively. Conversely,
every initial triple produces a unique vector satisfying all the recurrence equations. Hence the
map
\[
V\longrightarrow\mathbb R^3,\qquad x\longmapsto(x_1,x_2,x_3)
\]
is a linear isomorphism, so \(\dim V=3\).

Applying the recurrence twice,
\[
x_{i+4}=-(x_{i+1}+x_{i+2}+x_{i+3}).
\]
Since
\[
x_i+x_{i+1}+x_{i+2}+x_{i+3}=0,
\]
we obtain \(x_{i+4}=x_i\). Thus every solution is \(4\)-periodic on the available index range.
Equivalently, the characteristic polynomial
\[
r^3+r^2+r+1=(r+1)(r^2+1)
\]
has roots \(-1,\pm i\), all fourth roots of unity other than \(1\).
\end{solution}

\begin{problem}[CP-I-0036]
\label{prob:cp-i-0036}
\textbf{Parameter-Dependent Rank via Minors}

\textbf{Source-recovered canonical completion of the parameter-rank method note.}

\hspace*{1.5em}Let \(M(a,b)\in M_4(\mathbb R)\) be a matrix whose entries depend polynomially
on parameters \(a,b\). Prove that
\[
\dim(\operatorname{Im}M)=\operatorname{rank}(M),
\qquad
\dim(\ker M)=4-\operatorname{rank}(M).
\]
Explain how the rank of \(M(a,b)\) can be determined by its minors, and prove that the set of
parameter values where the rank is at most \(r\) is cut out by the vanishing of all
\((r+1)\times(r+1)\) minors. Conclude that rank drops can occur only on special algebraic
parameter loci.
\end{problem}

\begin{solution}
The first identity is the definition of matrix rank:
\[
\operatorname{rank}(M)=\dim(\operatorname{Im}M).
\]
Rank--nullity for \(M:\mathbb R^4\to\mathbb R^4\) gives
\[
\dim(\ker M)+\dim(\operatorname{Im}M)=4,
\]
hence
\[
\dim(\ker M)=4-\operatorname{rank}(M).
\]

A matrix has rank at least \(k\) if and only if at least one \(k\times k\) minor is nonzero.
Equivalently,
\[
\operatorname{rank}M(a,b)\le r
\]
if and only if every \((r+1)\times(r+1)\) minor vanishes.

Because the entries of \(M(a,b)\) are polynomial in \(a,b\), each such minor is also a polynomial
in \(a,b\). Therefore the parameter set on which the rank is at most \(r\) is the common zero
set of finitely many polynomials. Thus rank can drop only on special algebraic parameter loci.
\end{solution}

\begin{problem}[CP-I-0048]
\label{prob:cp-i-0048}
\textbf{Which modified lists are bases?}

\textbf{Exact statement.}
\hspace*{1.5em}Suppose that the vectors $e_1,\dots,e_n$ form a basis for a real vector space $V$.
Which of the following are also bases?

	\par\noindent (a)\quad $e_1+e_2,\ e_2+e_3,\ \dots,\ e_{n-1}+e_n,\ e_n$;
	\par\noindent (b)\quad $e_1+e_2,\ e_2+e_3,\ \dots,\ e_{n-1}+e_n,\ e_n+e_1$;
	\par\noindent (c)\quad $e_1-e_n,\ e_2+e_{n-1},\ \dots,\ e_n+(-1)^n e_1$.
\end{problem}

\begin{solution}
Write the proposed lists as $v_1,\dots,v_n$.

\medskip
\noindent\textbf{(a)} Let $v_k=e_k+e_{k+1}$ for $1\le k\le n-1$ and $v_n=e_n$.
Then we can recover the original basis:
\[
e_n=v_n,\qquad
e_{n-1}=v_{n-1}-e_n,\qquad
e_{n-2}=v_{n-2}-e_{n-1},\ \dots,\
e_1=v_1-e_2.
\]
Hence $v_1,\dots,v_n$ span $V$ and are linearly independent; they form a basis.

\medskip
\noindent\textbf{(b)} Let $v_k=e_k+e_{k+1}$ for $1\le k\le n-1$ and $v_n=e_n+e_1$.
Consider the alternating sum
\[
S:=v_1-v_2+v_3-\cdots+(-1)^{n-1}v_n.
\]
A direct cancellation shows:
\[
S=
\begin{cases}
	0, & n\ \text{even},\\[4pt]
	2e_1, & n\ \text{odd}.
\end{cases}
\]
If $n$ is even, this gives a nontrivial linear dependence, so the list is not a basis.
If $n$ is odd, then $e_1=\tfrac12 S$ lies in $\mathrm{span}(v_1,\dots,v_n)$, and then
\[
e_2=v_1-e_1,\quad e_3=v_2-e_2,\quad \dots,\quad e_n=v_{n-1}-e_{n-1},
\]
so all $e_i$ lie in the span; hence $v_1,\dots,v_n$ span $V$. Since there are $n$ vectors spanning an
$n$-dimensional space, they form a basis. Therefore (b) is a basis iff $n$ is odd.

\medskip
\noindent\textbf{(c)} Interpret the pattern as
\[
v_k:=e_k+(-1)^k e_{n+1-k}\qquad (1\le k\le n).
\]
Pair indices $k$ and $m:=n+1-k$. Then
\[
v_k=e_k+(-1)^k e_m,\qquad
v_m=e_m+(-1)^m e_k.
\]
If $n$ is odd, then $(-1)^m=(-1)^{n+1-k}=(-1)^{n+1}(-1)^k=(+1)(-1)^k=(-1)^k$, so
\[
v_k-(-1)^k v_m = 0
\]
for each pair, giving a nontrivial dependence; hence not a basis.

If $n$ is even, then $(-1)^m=(-1)^{n+1-k}=(-1)^{n+1}(-1)^k=-( -1)^k$, so
\[
v_k=e_k+s e_m,\qquad v_m=e_m-s e_k,\qquad (s:=(-1)^k).
\]
Solving yields
\[
e_k=\frac{v_k-sv_m}{2},\qquad e_m=\frac{sv_k+v_m}{2}.
\]
Thus every $e_i$ lies in $\mathrm{span}(v_1,\dots,v_n)$, so the $v_i$ span $V$ and hence form a basis.
Therefore (c) is a basis iff $n$ is even.

\bigskip
\end{solution}

\begin{problem}[CP-I-0051]
\label{prob:cp-i-0051}
\textbf{Generalised Eigenspace Decomposition Theorem}

\textbf{Exact statement.}
\hspace*{1.5em}Prove the Generalised Eigenspace Decomposition Theorem.

\end{problem}

\begin{solution}
Let $T\in\mathrm{End}(V)$ with $V$ finite-dimensional over $\mathbb C$ and
\[
m_T(t)=\prod_{i=1}^r (t-\lambda_i)^{k_i}
\]
with distinct $\lambda_i$.
Set $G_{\lambda_i}=\ker((T-\lambda_i I)^{k_i})$.
The factors $(t-\lambda_i)^{k_i}$ are pairwise coprime, so there exist polynomials $a_i(t)$ such that
\[
\sum_{i=1}^r a_i(t)\,\frac{m_T(t)}{(t-\lambda_i)^{k_i}}=1.
\]
Write $m_i(t)=m_T(t)/(t-\lambda_i)^{k_i}$ and define $P_i=a_i(T)m_i(T)$.
Then $\sum_i P_i=I$, $P_iP_j=0$ for $i\neq j$, and $P_i^2=P_i$.
Also $(T-\lambda_i I)^{k_i}P_i=0$, hence $\mathrm{im}(P_i)\subseteq G_{\lambda_i}$.
If $v\in G_{\lambda_i}$ then $m_j(T)v=0$ for $j\neq i$, so from $\sum_j P_j=I$ we get $v=P_iv$, hence $G_{\lambda_i}\subseteq\mathrm{im}(P_i)$.
Thus $\mathrm{im}(P_i)=G_{\lambda_i}$ and
\[
V=\sum_i \mathrm{im}(P_i)=\sum_i G_{\lambda_i}.
\]
The sum is direct because $P_iP_j=0$ implies $G_{\lambda_i}\cap\sum_{j\neq i}G_{\lambda_j}=\{0\}$.
Therefore,
\[
V=\bigoplus_{i=1}^r G_{\lambda_i}.
\]
\end{solution}

\begin{problem}[CP-I-0052]
\label{prob:cp-i-0052}
\textbf{Cauchy--Schwarz and two applications}

\textbf{Exact statement.}
\hspace*{1.5em}State the Cauchy--Schwarz inequality for vectors $\mathbf u$ and $\mathbf v$ in $\mathbb{R}^n$ and give a necessary and sufficient condition for equality to hold.

	\par\noindent\textbullet\quad By considering suitable vectors in $\mathbb{R}^3$, or otherwise, show that
	\[
	x^2+y^2+z^2 \ge yz+zx+xy
	\]
	for any real numbers $x,y,z$.
	\par\noindent\textbullet\quad By considering suitable vectors in $\mathbb{R}^4$, or otherwise, show that
	\[
	3(x^2+y^2+z^2+4)-2(yz+zx+xy)-4(x+y+z)=0
	\]
	holds for unique real values of $x,y,z$, to be determined.
\end{problem}

\begin{solution}

Cauchy--Schwarz: for $\mathbf u,\mathbf v\in\mathbb{R}^n$,
\[
|\mathbf u\cdot \mathbf v|\le \|\mathbf u\|\,\|\mathbf v\|.
\]
Equality holds if and only if $\mathbf u$ and $\mathbf v$ are linearly dependent (including the case that one is zero).

	\par\noindent\textbullet\quad Let $\mathbf u=(x,y,z)$ and $\mathbf v=(y,z,x)$ in $\mathbb{R}^3$. Then
	\[
	\mathbf u\cdot\mathbf v=xy+yz+zx,\qquad
	\|\mathbf u\|^2=\|\mathbf v\|^2=x^2+y^2+z^2.
	\]
	By Cauchy--Schwarz,
	\[
	|xy+yz+zx|\le \|\mathbf u\|\,\|\mathbf v\|=x^2+y^2+z^2,
	\]
	hence $xy+yz+zx\le x^2+y^2+z^2$.
	Equality forces $(x,y,z)=\lambda(y,z,x)$, which over $\mathbb{R}$ implies $x=y=z$ (or all zero).
	\par\noindent\textbullet\quad Write $s=x+y+z$ and $Q=x^2+y^2+z^2$. Since
	\[
	xy+yz+zx=\frac{s^2-Q}{2},
	\]
	the equation becomes
	\[
	3(Q+4)-2\cdot\frac{s^2-Q}{2}-4s=0
	\quad\Longleftrightarrow\quad
	4Q-s^2-4s+12=0.
	\]
	Let $m=s/3$ and set $p=x-m$, $q=y-m$, $r=z-m$ so that $p+q+r=0$ and
	\[
	Q=3m^2+(p^2+q^2+r^2).
	\]
	Substitution gives
	\[
	4Q-s^2-4s+12
	=3(m-2)^2+4(p^2+q^2+r^2).
	\]
	Hence the equation holds iff
	\[
	3(m-2)^2+4(p^2+q^2+r^2)=0,
	\]
	so $m=2$ and $p=q=r=0$, i.e.
	\[
	x=y=z=2.
	\]
	Therefore the unique real solution is $(x,y,z)=(2,2,2)$.
\end{solution}

\begin{problem}[CP-I-0056]
\label{prob:cp-i-0056}
\textbf{Descending images, ascending kernels, and rank constraints}

\textbf{Exact statement.}

	\par\hspace*{1.5em} (i)\quad Let $\alpha:V\to V$ be an endomorphism of a finite dimensional vector space $V$.
	Show that
	\[
	V\supseteq \operatorname{im}(\alpha)\supseteq \operatorname{im}(\alpha^2)\supseteq \cdots,\qquad
	\{0\}\subseteq \ker(\alpha)\subseteq \ker(\alpha^2)\subseteq \cdots.
	\]
	If $r_k=\mathrm{rk}(\alpha^k)$, deduce that $r_k\ge r_{k+1}$ and
	$r_k-r_{k+1}\ge r_{k+1}-r_{k+2}$.
	Conclude that if for some $k\ge 0$ we have $r_k=r_{k+1}$, then $r_k=r_{k+\ell}$ for all $\ell\ge 0$.
	\par\noindent (ii)\quad Suppose $\dim(V)=5$, $\alpha^3=0$, but $\alpha^2\ne 0$.
	What possibilities are there for $\mathrm{rk}(\alpha)$ and $\mathrm{rk}(\alpha^2)$?
\end{problem}

\begin{solution}
\textbf{(i)} Since
\[
\operatorname{im}(\alpha^{k+1})=\alpha(\operatorname{im}(\alpha^k))\subseteq \operatorname{im}(\alpha^k),
\]
the images form a descending chain. Also $\alpha^k(v)=0\Rightarrow \alpha^{k+1}(v)=\alpha(0)=0$, hence
$\ker(\alpha^k)\subseteq \ker(\alpha^{k+1})$.

Let $r_k=\dim\operatorname{im}(\alpha^k)$. The inclusion
$\operatorname{im}(\alpha^{k+1})\subseteq \operatorname{im}(\alpha^k)$ implies $r_k\ge r_{k+1}$.

Consider the restriction $\alpha:\operatorname{im}(\alpha^k)\to \operatorname{im}(\alpha^{k+1})$, which is surjective.
Rank--nullity for this restricted map gives
\[
\dim(\ker(\alpha)\cap \operatorname{im}(\alpha^k))
=\dim\operatorname{im}(\alpha^k)-\dim\operatorname{im}(\alpha^{k+1})
=r_k-r_{k+1}.
\]
Similarly,
\[
r_{k+1}-r_{k+2}=\dim(\ker(\alpha)\cap \operatorname{im}(\alpha^{k+1})).
\]
Since $\operatorname{im}(\alpha^{k+1})\subseteq \operatorname{im}(\alpha^k)$, we have
\[
\ker(\alpha)\cap \operatorname{im}(\alpha^{k+1})\subseteq \ker(\alpha)\cap \operatorname{im}(\alpha^k),
\]
so $r_{k+1}-r_{k+2}\le r_k-r_{k+1}$, i.e.
\[
r_k-r_{k+1}\ge r_{k+1}-r_{k+2}.
\]

If $r_k=r_{k+1}$, then $r_k-r_{k+1}=0$, hence $0\ge r_{k+1}-r_{k+2}$.
But also $r_{k+1}\ge r_{k+2}$, so $r_{k+1}=r_{k+2}$. Inductively, $r_k=r_{k+\ell}$ for all $\ell\ge 0$.

\medskip
\textbf{(ii)} Here $r_3=0$ and $r_2\ge 1$.
Also $\operatorname{im}(\alpha^2)\subseteq \ker(\alpha)$ because $\alpha(\alpha^2(v))=\alpha^3(v)=0$.
Thus
\[
r_2\le \dim\ker(\alpha)=5-r_1 \quad\Rightarrow\quad r_1+r_2\le 5.
\]
From (i), the successive drops satisfy
\[
(5-r_1)\ge (r_1-r_2)\ge (r_2-0)=r_2,
\]
so $r_1\ge 2r_2$ and $r_2\ge 1$.
The only integer solutions are
\[
(\mathrm{rk}(\alpha),\mathrm{rk}(\alpha^2))=(2,1)\quad\text{or}\quad(3,1).
\]
Both occur (Jordan types $3+1+1$ and $3+2$ respectively).
\end{solution}

\begin{problem}[CP-I-0057]
\label{prob:cp-i-0057}
\textbf{combinatorial picture for the six vectors in $\mathbb R^4$}

\hspace*{1.5em}The six vectors are $e_i+e_j$ for $1\le i<j\le 4$, which can be viewed as the edges of the complete graph $K_4$ on vertices $\{1,2,3,4\}$.
A useful rule of thumb:

	\par\noindent\textbullet\quad A set of such vectors is linearly independent exactly when the corresponding edges contain no cycle (a forest).
	\par\noindent\textbullet\quad A maximal independent set of size $3$ corresponds to a spanning tree on $4$ vertices.

Typical dependencies come from cycles, e.g.
\[
(e_1+e_2)+(e_3+e_4)=(e_1+e_3)+(e_2+e_4).
\]
\end{problem}

\begin{problem}[CP-I-0060]
\label{prob:cp-i-0060}
\textbf{A real $\lambda$ with $\det(A+\lambda B)\neq 0$; real similarity}

\textbf{Exact statement.}
\hspace*{1.5em}Let $C\in M_n(\mathbb C)$ and write $C=A+iB$ with $A,B\in M_n(\mathbb R)$.
By considering $\det(A+\lambda B)$ as a function of $\lambda$, show that if $C$ is invertible then there exists $\lambda\in\mathbb R$ such that $A+\lambda B$ is invertible.
Deduce that if $P,Q\in M_n(\mathbb R)$ are similar over $\mathbb C$, then they are similar over $\mathbb R$.

\end{problem}

\begin{solution}
Set $p(\lambda)=\det(A+\lambda B)\in\mathbb R[\lambda]$. If $A+\lambda B$ were singular for all real $\lambda$, then $p$ would have infinitely many real roots, hence $p\equiv 0$, so $0=p(i)=\det(A+iB)=\det(C)$, contradicting invertibility of $C$.
Thus $\exists \lambda\in\mathbb R$ with $p(\lambda)\neq 0$, i.e.\ $A+\lambda B$ invertible.

Now suppose $Q=S^{-1}PS$ for some invertible $S\in M_n(\mathbb C)$, with $P,Q\in M_n(\mathbb R)$.
Write $S=A+iB$ with $A,B$ real. The relation $PS=SQ$ gives
\[
P(A+iB)=(A+iB)Q \;\Rightarrow\; PA=AQ,\quad PB=BQ.
\]
Hence for any real $\lambda$,
\[
P(A+\lambda B)=(A+\lambda B)Q.
\]
Since $S=A+iB$ is invertible, the first part gives $\lambda\in\mathbb R$ such that $T:=A+\lambda B$ is invertible.
Then $PT=TQ$ implies $Q=T^{-1}PT$, so $P$ and $Q$ are similar over $\mathbb R$.
\end{solution}

\begin{problem}[CP-I-0061]
\label{prob:cp-i-0061}
\textbf{Concurrency of the altitudes (orthocenter) via vectors}

\textbf{Exact statement.}
\hspace*{1.5em}Show by vector methods that the altitudes of a triangle are concurrent.
\([\,\)Hint: let the altitudes \(AD,BE\) of \(\triangle ABC\) meet at \(H\), and show that \(CH\) is perpendicular to \(AB\,]\).

\end{problem}

\begin{solution}
Let the position vectors of \(A,B,C,H\) be \(\mathbf a,\mathbf b,\mathbf c,\mathbf h\).
Since \(H\) lies on the altitude from \(A\), we have \(AH\perp BC\), i.e.
\[
(\mathbf h-\mathbf a)\cdot(\mathbf c-\mathbf b)=0
\quad\Longrightarrow\quad
\mathbf h\cdot(\mathbf c-\mathbf b)=\mathbf a\cdot(\mathbf c-\mathbf b).
\]
Since \(H\) lies on the altitude from \(B\), we have \(BH\perp AC\), i.e.
\[
(\mathbf h-\mathbf b)\cdot(\mathbf c-\mathbf a)=0
\quad\Longrightarrow\quad
\mathbf h\cdot(\mathbf c-\mathbf a)=\mathbf b\cdot(\mathbf c-\mathbf a).
\]
Now write \(\mathbf b-\mathbf a=(\mathbf b-\mathbf c)+(\mathbf c-\mathbf a)\). Then
\[
\mathbf h\cdot(\mathbf b-\mathbf a)
=\mathbf h\cdot(\mathbf b-\mathbf c)+\mathbf h\cdot(\mathbf c-\mathbf a)
=-\mathbf h\cdot(\mathbf c-\mathbf b)+\mathbf h\cdot(\mathbf c-\mathbf a).
\]
Substitute the relations above:
\[
\mathbf h\cdot(\mathbf b-\mathbf a)
=-\mathbf a\cdot(\mathbf c-\mathbf b)+\mathbf b\cdot(\mathbf c-\mathbf a)
=\mathbf c\cdot(\mathbf b-\mathbf a).
\]
Therefore
\[
(\mathbf h-\mathbf c)\cdot(\mathbf b-\mathbf a)=0,
\]
so \(CH\perp AB\). Hence \(H\) also lies on the altitude from \(C\), and all three altitudes are concurrent.
\end{solution}

\begin{problem}[CP-I-0062]
\label{prob:cp-i-0062}
\textbf{Cross products in a triangle; sine rule; a converse}

\textbf{Exact statement.}
\hspace*{1.5em}(a) In \(\triangle ABC\), let \(\overrightarrow{AB}=\mathbf u\), \(\overrightarrow{BC}=\mathbf v\) and \(\overrightarrow{CA}=\mathbf w\).
Show that
\[
\mathbf u\times \mathbf v=\mathbf v\times \mathbf w=\mathbf w\times \mathbf u,
\]
and hence obtain the sine rule for \(\triangle ABC\).

(b) Given any three vectors \(\mathbf p,\mathbf q,\mathbf r\) such that
\[
\mathbf p\times \mathbf q=\mathbf q\times \mathbf r=\mathbf r\times \mathbf p,
\]
and \(|\mathbf p\times \mathbf q|\ne 0\), show that
\[
\mathbf p+\mathbf q+\mathbf r=\mathbf 0.
\]

\end{problem}

\begin{solution}
\textbf{(a)} In a triangle,
\[
\mathbf u+\mathbf v+\mathbf w=\mathbf 0
\quad\Longrightarrow\quad
\mathbf w=-(\mathbf u+\mathbf v).
\]
Then
\[
\mathbf v\times \mathbf w=\mathbf v\times\bigl(-(\mathbf u+\mathbf v)\bigr)=-(\mathbf v\times\mathbf u)=\mathbf u\times\mathbf v,
\]
and similarly \(\mathbf w\times\mathbf u=\mathbf u\times\mathbf v\). Hence
\(\mathbf u\times\mathbf v=\mathbf v\times\mathbf w=\mathbf w\times\mathbf u\).

Taking magnitudes and using \(|\mathbf a\times\mathbf b|=|\mathbf a||\mathbf b|\sin(\angle(\mathbf a,\mathbf b))\),
we obtain
\[
|\mathbf u||\mathbf v|\sin B=|\mathbf v||\mathbf w|\sin C=|\mathbf w||\mathbf u|\sin A,
\]
which gives the sine rule
\[
\frac{|BC|}{\sin A}=\frac{|CA|}{\sin B}=\frac{|AB|}{\sin C}.
\]

\medskip
\noindent\textbf{(b)} From \(\mathbf p\times\mathbf q=\mathbf q\times\mathbf r\),
\[
\mathbf p\times\mathbf q-\mathbf q\times\mathbf r=\mathbf 0
\ \Longrightarrow\
\mathbf p\times\mathbf q+\mathbf r\times\mathbf q=\mathbf 0
\ \Longrightarrow\
(\mathbf p+\mathbf r)\times\mathbf q=\mathbf 0,
\]
so \(\mathbf p+\mathbf r\parallel \mathbf q\).
Also \(\mathbf r\times\mathbf p=\mathbf p\times\mathbf q\) gives \((\mathbf q+\mathbf r)\times\mathbf p=\mathbf 0\),
so \(\mathbf q+\mathbf r\parallel \mathbf p\).
Since \(|\mathbf p\times\mathbf q|\ne 0\), \(\mathbf p\) and \(\mathbf q\) are linearly independent, forcing
\(\mathbf r=-\mathbf p-\mathbf q\). Hence \(\mathbf p+\mathbf q+\mathbf r=\mathbf 0\).
\end{solution}

\begin{problem}[CP-I-0063]
\label{prob:cp-i-0063}
\textbf{Image, kernel, and solvability of \(Mx=d\) (Vandermonde-type matrix)}

\textbf{Exact statement.}
\hspace*{1.5em}Let \(M\) be a real \(3\times 3\) matrix, and let \(d\) be a 3 component column vector. Explain briefly how the
general solution of the matrix equation \(Mx=d\), where \(x\) is a 3 component column vector, depends on
the image and kernel of the linear map \(x\mapsto Mx\).

Consider the case
\[
M=\begin{pmatrix}
	1&1&1\\
	1&a&b\\
	1&a^2&b^2
\end{pmatrix},\qquad
x=\begin{pmatrix}x\\y\\z\end{pmatrix},\qquad
d=\begin{pmatrix}1\\-1\\1\end{pmatrix}.
\]
Find the image and kernel of the corresponding map, noting the different possibilities according to different values of \(a\) and \(b\).
For which values of \(a\) and \(b\) do the equations \(Mx=d\) have (i) a unique solution, (ii) more than one solution, (iii) no solution?
For each pair \((a,b)\) satisfying (ii), give the solutions as the sum of a fixed solution and the general solution of the corresponding homogeneous equations.

\end{problem}

\begin{solution}
Let \(T:\mathbb R^3\to\mathbb R^3\) be \(T(x)=Mx\). Then:
\[
Mx=d \text{ is solvable } \Longleftrightarrow d\in \operatorname{Im}(M),
\]
and if \(x_0\) is one solution, then all solutions are
\[
x=x_0+v,\qquad v\in\ker(M).
\]
The solution is unique iff \(\ker(M)=\{0\}\), equivalently \(\det(M)\neq 0\).

Write the columns of \(M\) as
\[
c_1=\begin{pmatrix}1\\1\\1\end{pmatrix},\quad
c_2=\begin{pmatrix}1\\a\\a^2\end{pmatrix},\quad
c_3=\begin{pmatrix}1\\b\\b^2\end{pmatrix}.
\]
Then \(\operatorname{Im}(M)=\operatorname{span}\{c_1,c_2,c_3\}\) and a direct computation gives
\[
\det M=(a-1)(b-1)(b-a).
\]

\textbf{Image and kernel (by cases).}

\emph{(A) If \((a-1)(b-1)(a-b)\neq 0\).}
Then \(M\) is invertible, so
\[
\ker(M)=\{0\},\qquad \operatorname{Im}(M)=\mathbb R^3.
\]

\emph{(B) Rank \(2\) cases (exactly two of \(1,a,b\) coincide).}

	\par\noindent\textbullet\quad If \(a=1\) and \(b\neq 1\), then \(c_2=c_1\), hence
	\[
	\ker(M)=\operatorname{span}\{(1,-1,0)\},\quad
	\operatorname{Im}(M)=\operatorname{span}\left\{\begin{pmatrix}1\\1\\1\end{pmatrix},\begin{pmatrix}1\\b\\b^2\end{pmatrix}\right\}.
	\]
	\par\noindent\textbullet\quad If \(b=1\) and \(a\neq 1\), then \(c_3=c_1\), hence
	\[
	\ker(M)=\operatorname{span}\{(1,0,-1)\},\quad
	\operatorname{Im}(M)=\operatorname{span}\left\{\begin{pmatrix}1\\1\\1\end{pmatrix},\begin{pmatrix}1\\a\\a^2\end{pmatrix}\right\}.
	\]
	\par\noindent\textbullet\quad If \(a=b\neq 1\), then \(c_2=c_3\), hence
	\[
	\ker(M)=\operatorname{span}\{(0,1,-1)\},\quad
	\operatorname{Im}(M)=\operatorname{span}\left\{\begin{pmatrix}1\\1\\1\end{pmatrix},\begin{pmatrix}1\\a\\a^2\end{pmatrix}\right\}.
	\]

\emph{(C) Rank \(1\) case.}
If \(a=b=1\), then \(c_1=c_2=c_3\), hence
\[
\operatorname{Im}(M)=\operatorname{span}\left\{\begin{pmatrix}1\\1\\1\end{pmatrix}\right\},\qquad
\ker(M)=\{(u,v,w)\in\mathbb R^3:\ u+v+w=0\}.
\]

\textbf{Solving \(Mx=d\) for \(d=(1,-1,1)^T\).}

\emph{(i) Unique solution.}
This occurs exactly when \(\det M\neq 0\), i.e.\ when \((a-1)(b-1)(a-b)\neq 0\).
Then
\[
x=\frac{(a+1)(b+1)}{(a-1)(b-1)},\quad
y=\frac{2(b+1)}{(a-1)(a-b)},\quad
z=-\frac{2(a+1)}{(b-1)(a-b)}.
\]

\emph{(ii) More than one solution.}
This requires \(\det M=0\) and \(d\in\operatorname{Im}(M)\). This happens only for
\[
(a,b)=(1,-1),\quad (a,b)=(-1,1),\quad (a,b)=(-1,-1).
\]
For these pairs:

	\par\noindent\textbullet\quad If \((a,b)=(1,-1)\), one fixed solution is \((0,0,1)\), and \(\ker(M)=\operatorname{span}\{(1,-1,0)\}\), so
	\[
	(x,y,z)=(0,0,1)+s(1,-1,0),\qquad s\in\mathbb R.
	\]
	\par\noindent\textbullet\quad If \((a,b)=(-1,1)\), one fixed solution is \((0,1,0)\), and \(\ker(M)=\operatorname{span}\{(1,0,-1)\}\), so
	\[
	(x,y,z)=(0,1,0)+s(1,0,-1),\qquad s\in\mathbb R.
	\]
	\par\noindent\textbullet\quad If \((a,b)=(-1,-1)\), one fixed solution is \((0,1,0)\), and \(\ker(M)=\operatorname{span}\{(0,1,-1)\}\), so
	\[
	(x,y,z)=(0,1,0)+s(0,1,-1),\qquad s\in\mathbb R.
	\]

\emph{(iii) No solution.}
For all other \((a,b)\) with \(\det M=0\), one has \(d\notin \operatorname{Im}(M)\), hence no solutions exist.
\end{solution}

\begin{problem}[CP-I-0066]
\label{prob:cp-i-0066}
\textbf{Restriction and induced quotient map; block matrices}

\textbf{Exact statement.}
\hspace*{1.5em}Let $Y\le V$ and $Z\le W$ be subspaces of finite dimensional vector spaces.
Suppose $\alpha:V\to W$ is linear and $\alpha(Y)\subseteq Z$.
Show that $\alpha$ induces linear maps $\alpha|_Y:Y\to Z$ and
$\overline{\alpha}:V/Y\to W/Z$ defined by
\[
\alpha|_Y(y)=\alpha(y),\qquad \overline{\alpha}(v+Y)=\alpha(v)+Z.
\]
Let $(v_1,\dots,v_n)$ be a basis of $V$ containing a basis $(v_1,\dots,v_k)$ for $Y$ and
$(w_1,\dots,w_m)$ a basis of $W$ containing a basis $(w_1,\dots,w_\ell)$ for $Z$.
Show the matrix of $\alpha$ has the block form
\[
\begin{pmatrix}
	A & C\\
	0 & B
\end{pmatrix},
\]
and explain how to read off the matrices of $\alpha|_Y$ and $\overline{\alpha}$ from this block matrix.
\end{problem}

\begin{solution}
Since $\alpha(Y)\subseteq Z$, the restriction $\alpha|_Y:Y\to Z$ is well-defined and linear.

Define $\overline{\alpha}:V/Y\to W/Z$ by $\overline{\alpha}(v+Y)=\alpha(v)+Z$.
If $v+Y=v'+Y$, then $v-v'\in Y$, so $\alpha(v-v')\in Z$, hence $\alpha(v)+Z=\alpha(v')+Z$.
Thus $\overline{\alpha}$ is well-defined, and it is linear.

With the chosen bases, for $j\le k$ we have $v_j\in Y$ and hence $\alpha(v_j)\in Z$.
Therefore the coordinates of $\alpha(v_j)$ along $w_{\ell+1},\dots,w_m$ vanish, which forces the lower-left block to be $0$.
Thus the matrix of $\alpha$ has the form
\[
[\alpha]=\begin{pmatrix}
	A & C\\
	0 & B
\end{pmatrix},
\]
where $A$ is the $\ell\times k$ block describing $\alpha(v_1),\dots,\alpha(v_k)$ in the basis of $Z$,
and $B$ is the $(m-\ell)\times(n-k)$ block describing the components of $\alpha(v_{k+1}),\dots,\alpha(v_n)$
modulo $Z$.

Hence the matrix of $\alpha|_Y:Y\to Z$ with respect to bases $(v_1,\dots,v_k)$ and $(w_1,\dots,w_\ell)$ is exactly $A$.

Moreover, the induced map $\overline{\alpha}:V/Y\to W/Z$ is represented (with respect to bases
$(v_{k+1}+Y,\dots,v_n+Y)$ and $(w_{\ell+1}+Z,\dots,w_m+Z)$) by the block $B$, since modding out by $Z$
kills the top $\ell$ coordinates.

\bigskip
\end{solution}

\begin{problem}[CP-I-0069]
\label{prob:cp-i-0069}
\textbf{Dual of \(P\) as sequences; dual maps of \(D\) and \(S\)}

\textbf{Exact statement.}
\hspace*{1.5em}Show \(P^*\cong\mathbb R^{\mathbb N}\) via \(\xi\mapsto (\xi(1),\xi(t),\xi(t^2),\dots)\).
Under this identification, describe the dual maps of
\(D(p)(t)=p'(t)\), \(S(p)(t)=p(t^2)\), and of \(DS\), \(SD\).
Verify \((DS)^*=S^*D^*\) and \((SD)^*=D^*S^*\).

\end{problem}

\begin{solution}
Let \(a_k=\xi(t^k)\). For \(T:P\to P\), \(b_k=(T^*\xi)(t^k)=\xi(T(t^k))\).

For \(D\): \(D(t^k)=k t^{k-1}\), so
\[
b_0=0,\qquad b_k=k a_{k-1}\ (k\ge1).
\]
For \(S\): \(S(t^k)=t^{2k}\), so
\[
b_k=a_{2k}.
\]
For \(DS\): \(DS(t^k)=D(t^{2k})=2k\,t^{2k-1}\), so
\[
b_0=0,\qquad b_k=2k\,a_{2k-1}\ (k\ge1).
\]
For \(SD\): \(SD(t^k)=S(k t^{k-1})=k\,t^{2k-2}\), so
\[
b_0=0,\qquad b_k=k\,a_{2k-2}\ (k\ge1).
\]
Then \((S^*D^*a)_k=(D^*a)_{2k}=2k a_{2k-1}\) and \((D^*S^*a)_k=k(S^*a)_{k-1}=k a_{2k-2}\), giving the identities.
\end{solution}

\begin{problem}[CP-I-0072]
\label{prob:cp-i-0072}
\textbf{Nondegenerate skew-symmetric forms and maximal isotropic dimension}

\textbf{Exact statement.}
\hspace*{1.5em}A bilinear form $\varphi:\mathbb R^n\times\mathbb R^n\to\mathbb R$ is called \emph{skew-symmetric} if $\varphi(u,v)=-\varphi(v,u)$ for all $u,v$.
If $\varphi$ is non-degenerate, show that $n$ is even, and that there is a basis with respect to which the matrix of $\varphi$ is block-diagonal with blocks
\[
J=\begin{pmatrix}0&1\\-1&0\end{pmatrix}.
\]
What is the maximum dimension of a subspace on which $\varphi$ vanishes?

\end{problem}

\begin{solution}
Let $A$ be the matrix of $\varphi$ in some basis. Skew-symmetry gives $A^T=-A$, hence
\[
\det(A)=\det(A^T)=\det(-A)=(-1)^n\det(A).
\]
If $n$ is odd then $\det(A)=0$, contradicting nondegeneracy. Thus $n=2m$ is even.

To obtain the canonical form, choose $v_1\neq 0$. Nondegeneracy gives $w_1$ with $\varphi(v_1,w_1)\neq 0$; scale so $\varphi(v_1,w_1)=1$.
Let $U_1=\mathrm{span}\{v_1,w_1\}$ and $U_1^\perp=\{x:\varphi(x,U_1)=0\}$.
Then $\mathbb R^{2m}=U_1\oplus U_1^\perp$ and $\varphi|_{U_1^\perp}$ is again nondegenerate.
Iterating yields a basis $(v_1,w_1,\dots,v_m,w_m)$ in which the matrix is $\mathrm{diag}(J,\dots,J)$.

If $W$ is a subspace on which $\varphi$ vanishes (isotropic), then $W\subseteq W^\perp$.
Since $\varphi$ is nondegenerate, $\dim W+\dim W^\perp=2m$, so $2\dim W\le 2m$, hence $\dim W\le m=\frac n2$.
This bound is sharp (take $W=\mathrm{span}\{v_1,\dots,v_m\}$), so the maximum isotropic dimension is $n/2$.
\end{solution}

\begin{problem}[CP-I-0074]
\label{prob:cp-i-0074}
\textbf{Medians of a complex triangle are concurrent}

\textbf{Exact statement.}
\hspace*{1.5em}Consider a triangle in the complex plane with vertices at \(0\), \(z_1\) and \(z_2\). Write down an expression for the
general point on the median through \(z_1\), and a similar expression for the general point on the median through \(z_2\).
Show that the three medians of the triangle are concurrent.

\end{problem}

\begin{solution}
The midpoint of the side from \(0\) to \(z_2\) is \(z_2/2\). Hence the median through \(z_1\) is the line segment joining
\(z_1\) to \(z_2/2\), so a general point on it is
\[
m_1(t)=(1-t)z_1+t\frac{z_2}{2},\qquad t\in\mathbb{R}.
\]
Similarly, the midpoint of the side from \(0\) to \(z_1\) is \(z_1/2\), so the median through \(z_2\) is
\[
m_2(s)=(1-s)z_2+s\frac{z_1}{2},\qquad s\in\mathbb{R}.
\]
Let
\[
g=\frac{z_1+z_2}{3}.
\]
Then
\[
g=\frac13 z_1+\frac13 z_2=(1-\tfrac23)z_1+\tfrac23\cdot\frac{z_2}{2}=m_1\!\left(\frac23\right),
\]
so \(g\) lies on the median through \(z_1\). Similarly,
\[
g=(1-\tfrac23)z_2+\tfrac23\cdot\frac{z_1}{2}=m_2\!\left(\frac23\right),
\]
so \(g\) lies on the median through \(z_2\). The median through \(0\) goes to the midpoint \((z_1+z_2)/2\), and
\[
g=\frac{2}{3}\cdot \frac{z_1+z_2}{2},
\]
so \(g\) lies on that median too. Hence all three medians meet at \(g\), so they are concurrent.
\end{solution}

\begin{problem}[CP-I-0076]
\label{prob:cp-i-0076}
\textbf{Commuting Operators Preserve Eigenspaces}

\textbf{Canonical authored completion of the retained invariance fragment.}

\hspace*{1.5em}Let \(A,B\in\operatorname{End}(V)\) satisfy \(AB=BA\).
For an eigenvalue \(\lambda\) of \(A\), prove that
\[
B\bigl(E_\lambda(A)\bigr)\subseteq E_\lambda(A),
\qquad
E_\lambda(A)=\ker(A-\lambda I).
\]
Deduce that if \(A\) is diagonalizable, then
\[
V=\bigoplus_{\lambda}E_\lambda(A)
\]
is a decomposition of \(V\) into \(B\)-invariant subspaces.
\end{problem}

\begin{solution}
Let \(x\in E_\lambda(A)\). Then \(Ax=\lambda x\). Using \(AB=BA\),
\[
A(Bx)=B(Ax)=B(\lambda x)=\lambda Bx.
\]
Hence \(Bx\in E_\lambda(A)\), proving
\[
B(E_\lambda(A))\subseteq E_\lambda(A).
\]

If \(A\) is diagonalizable, then its eigenspaces span \(V\) and their sum is direct:
\[
V=\bigoplus_{\lambda}E_\lambda(A).
\]
The first part shows that every summand is preserved by \(B\), so this is a decomposition into
\(B\)-invariant subspaces.
\end{solution}

\begin{problem}[CP-I-0083]
\label{prob:cp-i-0083}
\textbf{Matrices of a Linear Map in General Bases}

\textbf{Canonical authored completion of the retained change-of-basis fragment.}

\hspace*{1.5em}Let \(T:V\to W\) be linear. Let \(E,E'\) be bases of \(V\) and \(F,F'\)
bases of \(W\). Suppose
\[
[T(v)]_F=A_{F,E}[v]_E.
\]
If
\[
[v]_E=C_E[v]_{E'},\qquad [w]_F=C_F[w]_{F'},
\]
prove that
\[
A_{F',E'}=C_F^{-1}A_{F,E}C_E.
\]
Deduce that when \(V=W\) and the same basis change is used in domain and codomain, matrices
of the same endomorphism are similar.
\end{problem}

\begin{solution}
Starting from
\[
[T(v)]_F=A_{F,E}[v]_E
\]
and using the coordinate changes,
\[
C_F[T(v)]_{F'}=A_{F,E}C_E[v]_{E'}.
\]
Multiplying by \(C_F^{-1}\),
\[
[T(v)]_{F'}=C_F^{-1}A_{F,E}C_E[v]_{E'}.
\]
Since this holds for every \(v\),
\[
A_{F',E'}=C_F^{-1}A_{F,E}C_E.
\]

If \(V=W\) and the same change-of-basis matrix \(R\) is used on both sides, then
\[
A'=R^{-1}AR.
\]
Thus matrices representing the same endomorphism in two bases are similar.
\end{solution}

\begin{problem}[CP-I-0084]
\label{prob:cp-i-0084}
\textbf{Solving linear vector equations with \(\times\) and \(\cdot\)}

\textbf{Exact statement.}
\hspace*{1.5em}Let \(\mathbf a,\mathbf b,\mathbf c,\mathbf d\) be fixed vectors in three dimensions. For each of the following equations,
find all solutions for \(\mathbf r\):
\[
\text{(i)}\ \mathbf r+\mathbf r\times\mathbf d=\mathbf c;\qquad
\text{(ii)}\ \mathbf r+(\mathbf r\cdot\mathbf a)\mathbf b=\mathbf c.
\]
\([\,\)In (ii), consider separately the cases \(\mathbf a\cdot\mathbf b\ne -1\) and \(\mathbf a\cdot\mathbf b=-1\,]\).

\end{problem}

\begin{solution}
\textbf{(i)} A closed-form solution is
\[
\mathbf r=\frac{\mathbf c-\mathbf c\times\mathbf d+(\mathbf c\cdot\mathbf d)\mathbf d}{1+|\mathbf d|^2}.
\]
Indeed, using \((\mathbf c\times\mathbf d)\times\mathbf d=\mathbf d(\mathbf c\cdot\mathbf d)-\mathbf c|\mathbf d|^2\),
one checks that \(\mathbf r+\mathbf r\times\mathbf d=\mathbf c\).

\medskip
\noindent\textbf{(ii)} Let \(s=\mathbf r\cdot\mathbf a\). Then \(\mathbf r=\mathbf c-s\mathbf b\), and dotting with \(\mathbf a\) gives
\[
s=\mathbf c\cdot\mathbf a-s(\mathbf b\cdot\mathbf a)
\quad\Longrightarrow\quad
s(1+\mathbf a\cdot\mathbf b)=\mathbf c\cdot\mathbf a.
\]
If \(\mathbf a\cdot\mathbf b\ne -1\), then
\[
s=\frac{\mathbf c\cdot\mathbf a}{1+\mathbf a\cdot\mathbf b},
\qquad
\mathbf r=\mathbf c-\frac{\mathbf c\cdot\mathbf a}{1+\mathbf a\cdot\mathbf b}\,\mathbf b.
\]
If \(\mathbf a\cdot\mathbf b=-1\), then solutions exist iff \(\mathbf c\cdot\mathbf a=0\). In that case \(s\) is arbitrary and
\[
\mathbf r=\mathbf c-t\,\mathbf b,\qquad t\in\mathbb R.
\]
If \(\mathbf c\cdot\mathbf a\ne 0\), there are no solutions.
\end{solution}

\begin{problem}[CP-I-0085]
\label{prob:cp-i-0085}
\textbf{Cyclic basis for a nilpotent map; commuting endomorphisms}

\textbf{Exact statement.}
\hspace*{1.5em}Let $\dim V=n$ and $\alpha\in\mathrm{End}(V)$ satisfy $\alpha^n=0$ but $\alpha^{n-1}\neq 0$.
Show there exists $y$ such that $(y,\alpha y,\dots,\alpha^{n-1}y)$ is a basis.
Find the matrix of $\alpha$ and of $\alpha^k$ in this basis.
Show that if $\beta$ commutes with $\alpha$, then $\beta=p(\alpha)$ for some polynomial $p$.
\end{problem}

\begin{solution}
Since $\alpha^{n-1}\neq 0$, choose $y$ with $\alpha^{n-1}y\neq 0$.
Suppose $\sum_{i=0}^{n-1} c_i\alpha^i y=0$ and let $r$ be the smallest index with $c_r\neq 0$.
Applying $\alpha^{n-1-r}$ gives
\[
0=\sum_{i=r}^{n-1} c_i\alpha^{i+n-1-r}y=c_r\alpha^{n-1}y,
\]
so $c_r=0$, a contradiction. Hence $y,\alpha y,\dots,\alpha^{n-1}y$ are linearly independent and form a basis.

Let $b_i=\alpha^i y$ for $i=0,\dots,n-1$. Then $\alpha(b_i)=b_{i+1}$ for $i\le n-2$ and $\alpha(b_{n-1})=0$.
Thus
\[
[\alpha]_{(b_0,\dots,b_{n-1})}=
\begin{pmatrix}
	0&0&0&\cdots&0\\
	1&0&0&\cdots&0\\
	0&1&0&\cdots&0\\
	\vdots&\ddots&\ddots&\ddots&\vdots\\
	0&\cdots&0&1&0
\end{pmatrix}.
\]
Moreover, $\alpha^k(b_i)=b_{i+k}$ if $i+k\le n-1$ and $0$ otherwise, so $[\alpha^k]$ has $1$'s on the $k$-th subdiagonal.

If $\beta\alpha=\alpha\beta$, write
\[
\beta(y)=\sum_{j=0}^{n-1} c_j \alpha^j y.
\]
Then for each $i$,
\[
\beta(\alpha^i y)=\alpha^i\beta(y)=\sum_{j=0}^{n-1} c_j \alpha^{i+j}y.
\]
But the operator $p(\alpha)=\sum_{j=0}^{n-1} c_j \alpha^j$ acts identically on the basis vectors $\alpha^i y$.
Hence $\beta=p(\alpha)$.
\end{solution}

\begin{problem}[CP-I-0087]
\label{prob:cp-i-0087}
\textbf{Functionals and intersections of kernels}

\textbf{Exact statement.}
\hspace*{1.5em}Let \(V\) be a finite-dimensional vector space and \(f_1,\dots,f_n,g\in V^*\).
Show
\[
g\in \operatorname{span}(f_1,\dots,f_n)
\iff
\bigcap_{i=1}^n \ker(f_i)\subseteq \ker(g).
\]
What if \(V\) is infinite dimensional?



The trace satisfies the cyclic identity \(\operatorname{tr}(AB)=\operatorname{tr}(BA)\).
The commutator subspace
\[
[M_n,M_n]=\operatorname{span}\{AB-BA:A,B\in M_n(\mathbb F)\}
\]
equals \(\{X\in M_n(\mathbb F):\operatorname{tr}(X)=0\}\). Hence any linear functional \(f\) with \(f(AB)=f(BA)\) vanishes on
commutators and factors through the 1-dimensional quotient \(M_n/[M_n,M_n]\); therefore \(f\) is a scalar multiple of \(\operatorname{tr}\),
and the normalization \(f(I)=n\) forces \(f=\operatorname{tr}\).

For endomorphisms on a finite-dimensional \(V\),
\[
\operatorname{tr}(\alpha\beta)=\operatorname{tr}(\beta\alpha)\quad\Rightarrow\quad \operatorname{tr}(\alpha\beta-\beta\alpha)=0,
\]
so a commutator cannot equal \(id_V\) since \(\operatorname{tr}(id_V)=\dim V\neq0\).
On \(C^\infty(\mathbb R)\), differentiation \(D\) and multiplication \(M_t\) satisfy the Weyl relation
\([D,M_t]=I\).

For distinct \(a_0,\dots,a_n\), the Vandermonde matrix \(V=(a_j^i)_{0\le i,j\le n}\) has
\[
\det(V)=\prod_{0\le i<j\le n}(a_j-a_i)\neq 0,
\]
equivalently: a polynomial \(p\in P_n\) with \(p(a_0)=\cdots=p(a_n)=0\) must be \(0\).
Thus the evaluation map
\[
\mathrm{ev}:P_n\to \mathbb F^{n+1},\quad p\mapsto (p(a_0),\dots,p(a_n))
\]
is an isomorphism, and the evaluation functionals \(e_{a_i}(p)=p(a_i)\) form a basis of \(P_n^*\).
The dual basis in \(P_n\) is the Lagrange basis \(\ell_i\) defined by \(\ell_i(a_j)=\delta_{ij}\), explicitly
\[
\ell_i(t)=\prod_{j\ne i}\frac{t-a_j}{a_i-a_j}.
\]

Monomials \((1,t,t^2,\dots)\) form a basis of \(P\), so every \(\xi\in P^*\) is determined by \(a_k=\xi(t^k)\).
Hence \(P^*\cong \mathbb R^{\mathbb N}\) via \(\xi\mapsto (a_0,a_1,\dots)\).
For \(T:P\to P\), the dual map is \(T^*\xi=\xi\circ T\). In sequence form:
\[
b_k=(T^*\xi)(t^k)=\xi(T(t^k)).
\]
Composition reverses: \((ST)^*=T^*S^*\).

For a finite family \(f_1,\dots,f_n\in V^*\) and \(g\in V^*\),
\[
g\in \operatorname{span}(f_1,\dots,f_n)
\iff
\bigcap_{i=1}^n \ker(f_i)\subseteq \ker(g).
\]
Proof idea: define \(F:V\to\mathbb F^n\) by \(F(v)=(f_1(v),\dots,f_n(v))\). Then \(\ker(F)=\bigcap\ker(f_i)\).
If \(\ker(F)\subseteq\ker(g)\), then \(g\) factors through \(V/\ker(F)\cong \operatorname{im}(F)\subseteq\mathbb F^n\), hence
\(g=\sum_{i=1}^n c_i f_i\).
This equivalence holds for any \(V\) (finite or infinite dimensional) provided the family \(f_1,\dots,f_n\) is finite.
For infinite families \(\{f_i\}\), the kernel condition may hold while \(g\) is not a finite linear combination of the \(f_i\).

For distinct scalars \(a_0,\dots,a_n\), the Vandermonde matrix is
\[
V(a_0,\dots,a_n)=
\begin{pmatrix}
	1 & 1 & \cdots & 1\\
	a_0 & a_1 & \cdots & a_n\\
	a_0^2 & a_1^2 & \cdots & a_n^2\\
	\vdots & \vdots & \ddots & \vdots\\
	a_0^n & a_1^n & \cdots & a_n^n
\end{pmatrix}.
\]
Its determinant is
\[
\det V(a_0,\dots,a_n)=\prod_{0\le i<j\le n}(a_j-a_i),
\]
so \(\det V\neq 0\) iff the \(a_i\) are pairwise distinct.

For distinct nodes \(a_0,\dots,a_n\), define \(\ell_0,\dots,\ell_n\in P_n\) by
\(\ell_i(a_j)=\delta_{ij}\). Explicitly,
\[
\ell_i(t)=\prod_{\substack{0\le j\le n\\ j\ne i}}\frac{t-a_j}{a_i-a_j}.
\]
Then every \(p\in P_n\) satisfies the interpolation formula
\[
p(t)=\sum_{i=0}^n p(a_i)\,\ell_i(t).
\]
If \(e_{a_j}(p)=p(a_j)\), then \(e_{a_j}(\ell_i)=\delta_{ij}\), so \((\ell_i)\) is the basis of \(P_n\) dual to
\((e_{a_0},\dots,e_{a_n})\).\end{problem}

\begin{solution}
(\(\Rightarrow\)) If \(g=\sum_{i=1}^n c_i f_i\) and \(v\in\bigcap \ker(f_i)\), then \(g(v)=\sum c_i f_i(v)=0\), so \(v\in\ker(g)\).

(\(\Leftarrow\)) Define \(F:V\to\mathbb F^n\) by \(F(v)=(f_1(v),\dots,f_n(v))\). Then \(\ker(F)=\bigcap_{i=1}^n\ker(f_i)\).
If \(\ker(F)\subseteq\ker(g)\), then \(g\) induces \(\bar g\) on \(V/\ker(F)\). Since \(V/\ker(F)\cong \operatorname{im}(F)\subseteq\mathbb F^n\),
there is \(\varphi\in(\operatorname{im}F)^*\) with \(g=\varphi\circ F\). Extend \(\varphi\) to \(\mathbb F^n\), so \(\varphi(y)=c\cdot y\).
Hence \(g(v)=\sum_{i=1}^n c_i f_i(v)\), i.e. \(g\in\operatorname{span}(f_1,\dots,f_n)\).

The same argument works for infinite-dimensional \(V\) (no finiteness of \(V\) is used), so the equivalence remains true for any \(V\) when
\(f_1,\dots,f_n\) is a finite family.
\end{solution}

\begin{problem}[CP-I-0090]
\label{prob:cp-i-0090}
\textbf{Bases for $U$, $W$, $U\cap W$, and a description of $U+W$ in $\mathbb R^5$}

\textbf{Exact statement.}
\hspace*{1.5em}Let
\begin{center}
\resizebox{\linewidth}{!}{$\displaystyle
U=\{x\in\mathbb R^5:\ x_1+x_3+x_4=0,\ \ 2x_1+2x_2+x_5=0\},\qquad
W=\{x\in\mathbb R^5:\ x_1+x_5=0,\ \ x_2=x_3=x_4\}.
$}
\end{center}
Find bases for $U$ and $W$ containing a basis for $U\cap W$ as a subset.
Give a basis for $U+W$ and show that
\[
U+W=\{x\in\mathbb R^5:\ x_1+2x_2+x_5=x_3+x_4\}.
\]
\end{problem}

\begin{solution}
For $U$, solve
\[
x_1=-x_3-x_4,\qquad x_5=-2x_1-2x_2=2x_3+2x_4-2x_2.
\]
With free variables $x_2,x_3,x_4$,
\[
x=x_2(0,1,0,0,-2)+x_3(-1,0,1,0,2)+x_4(-1,0,0,1,2).
\]
Thus a basis of $U$ is
\[
u_1=(0,1,0,0,-2),\quad u_2=(-1,0,1,0,2),\quad u_3=(-1,0,0,1,2).
\]

For $W$, write $x_2=x_3=x_4=t$ and $x_1=s$, so $x_5=-s$. Then
\[
x=(s,t,t,t,-s)=s(1,0,0,0,-1)+t(0,1,1,1,0),
\]
hence a basis of $W$ is
\[
w_1=(1,0,0,0,-1),\quad w_2=(0,1,1,1,0).
\]

For $U\cap W$, take $x=(s,t,t,t,-s)\in W$ and impose $x_1+x_3+x_4=0$:
\[
s+2t=0\Rightarrow s=-2t.
\]
Therefore
\[
U\cap W=\mathrm{span}\{v\},\qquad v=(-2,1,1,1,2).
\]
Note that $v=u_1+u_2+u_3$ and $v=-2w_1+w_2$, so bases containing $v$ are
\[
\mathcal B_U=\{v,u_2,u_3\}\ \text{ for }U,\qquad \mathcal B_W=\{v,w_1\}\ \text{ for }W.
\]
Since $\dim(U+W)=3+2-1=4$, a basis for $U+W$ is
\[
\mathcal B_{U+W}=\{v,u_2,u_3,w_1\}.
\]

Let $S:=\{x\in\mathbb R^5:\ x_1+2x_2+x_5=x_3+x_4\}$.
If $x\in U+W$ and $x=a v+b u_2+c u_3+d w_1$, then a coordinate check gives
$x_1+2x_2+x_5=x_3+x_4$, hence $U+W\subseteq S$.

Conversely, given $x\in S$, set
\[
a=x_2,\qquad b=x_3-x_2,\qquad c=x_4-x_2,\qquad d=x_1+x_3+x_4.
\]
Then one verifies $x=a v+b u_2+c u_3+d w_1$, hence $x\in U+W$, so $S\subseteq U+W$.
Therefore $U+W=S$.

\bigskip
\end{solution}

\begin{problem}[CP-I-0093]
\label{prob:cp-i-0093}
\textbf{A vector equation with dot–product constraints and equality cases}

\textbf{Exact statement.}
\hspace*{1.5em}In three dimensions, the vectors \(\mathbf x,\mathbf y,\mathbf a\) satisfy
\[
\mathbf x+\mathbf y(\mathbf x\cdot \mathbf y)=\mathbf a.
\]
Show that
\[
(\mathbf x\cdot \mathbf y)^2=\frac{|\mathbf a|^2-|\mathbf x|^2}{2+|\mathbf y|^2}.
\]
Use an inequality involving \(\mathbf x\cdot \mathbf y\) and the lengths of \(\mathbf x\) and \(\mathbf y\) to deduce that
\[
|\mathbf x|\,(1+|\mathbf y|^2)\ge |\mathbf a|\ge |\mathbf x|.
\]
Explain the circumstances in which either of the inequalities above become equalities, and describe the relation between
\(\mathbf x,\mathbf y,\mathbf a\) in these circumstances.

\end{problem}

\begin{solution}
Let \(t=\mathbf x\cdot \mathbf y\). Then \(\mathbf a=\mathbf x+t\mathbf y\). Taking norms,
\[
|\mathbf a|^2=|\mathbf x+t\mathbf y|^2
=|\mathbf x|^2+2t(\mathbf x\cdot \mathbf y)+t^2|\mathbf y|^2
=|\mathbf x|^2+t^2(2+|\mathbf y|^2).
\]
Hence
\[
t^2=\frac{|\mathbf a|^2-|\mathbf x|^2}{2+|\mathbf y|^2}.
\]
From this, \(|\mathbf a|^2\ge |\mathbf x|^2\), so \(|\mathbf a|\ge |\mathbf x|\).
Also by Cauchy--Schwarz, \(|t|=|\mathbf x\cdot \mathbf y|\le |\mathbf x||\mathbf y|\), so
\[
\frac{|\mathbf a|^2-|\mathbf x|^2}{2+|\mathbf y|^2}=t^2\le |\mathbf x|^2|\mathbf y|^2.
\]
Therefore
\[
|\mathbf a|^2\le |\mathbf x|^2\bigl(1+(2+|\mathbf y|^2)|\mathbf y|^2\bigr)
=|\mathbf x|^2(1+|\mathbf y|^2)^2,
\]
hence \(|\mathbf a|\le |\mathbf x|(1+|\mathbf y|^2)\).

Equality \(|\mathbf a|=|\mathbf x|\) holds iff \(t=0\), i.e. \(\mathbf x\cdot \mathbf y=0\), in which case \(\mathbf a=\mathbf x\).
Equality \(|\mathbf a|=|\mathbf x|(1+|\mathbf y|^2)\) holds iff \(|\mathbf x\cdot\mathbf y|=|\mathbf x||\mathbf y|\), i.e. \(\mathbf x\parallel \mathbf y\),
and then \(\mathbf a=\mathbf x+t\mathbf y\) is also parallel to them.
\end{solution}

\begin{problem}[CP-I-0094]
\label{prob:cp-i-0094}
\textbf{Lagrange’s identity, Jacobi identity, and scalar triple products}

\textbf{Exact statement.}
\hspace*{1.5em}(a) Using the identity \(\mathbf a\times(\mathbf b\times\mathbf c)=(\mathbf a\cdot\mathbf c)\mathbf b-(\mathbf a\cdot\mathbf b)\mathbf c\), show that
\[
\text{(i)}\quad
(\mathbf a\times\mathbf b)\cdot(\mathbf c\times\mathbf d)
=(\mathbf a\cdot\mathbf c)(\mathbf b\cdot\mathbf d)-(\mathbf a\cdot\mathbf d)(\mathbf b\cdot\mathbf c),
\]
\[
\text{(ii)}\quad
\mathbf a\times(\mathbf b\times\mathbf c)+\mathbf b\times(\mathbf c\times\mathbf a)+\mathbf c\times(\mathbf a\times\mathbf b)=\mathbf 0.
\]
Relate the case \(\mathbf c=\mathbf a\) and \(\mathbf d=\mathbf b\) in (i) to a well-known trigonometric identity.
Evaluate \((\mathbf a\times\mathbf b)\times(\mathbf c\times\mathbf d)\) in two ways and compare to find an explicit linear combination of
\(\mathbf a,\mathbf b,\mathbf c,\mathbf d\) that is zero.

(b) Given that \([\mathbf a,\mathbf b,\mathbf c]=\mathbf a\cdot(\mathbf b\times\mathbf c)\), show that
\[
[\mathbf a\times\mathbf b,\ \mathbf b\times\mathbf c,\ \mathbf c\times\mathbf a]=[\mathbf a,\mathbf b,\mathbf c]^2.
\]

\end{problem}

\begin{solution}
\textbf{(a)(i)} Using \((\mathbf a\times\mathbf b)\cdot(\mathbf c\times\mathbf d)=\mathbf c\cdot(\mathbf d\times(\mathbf a\times\mathbf b))\) and
\(\mathbf d\times(\mathbf a\times\mathbf b)=(\mathbf d\cdot\mathbf b)\mathbf a-(\mathbf d\cdot\mathbf a)\mathbf b\), we obtain
\[
(\mathbf a\times\mathbf b)\cdot(\mathbf c\times\mathbf d)
=(\mathbf a\cdot\mathbf c)(\mathbf b\cdot\mathbf d)-(\mathbf a\cdot\mathbf d)(\mathbf b\cdot\mathbf c).
\]
Putting \(\mathbf c=\mathbf a\), \(\mathbf d=\mathbf b\) gives
\[
|\mathbf a\times\mathbf b|^2=|\mathbf a|^2|\mathbf b|^2-(\mathbf a\cdot\mathbf b)^2,
\]
equivalent to \(\sin^2\theta+\cos^2\theta=1\).

\textbf{(a)(ii)} Expand each term:
\[
\mathbf a\times(\mathbf b\times\mathbf c)=(\mathbf a\cdot\mathbf c)\mathbf b-(\mathbf a\cdot\mathbf b)\mathbf c,
\]
\[
\mathbf b\times(\mathbf c\times\mathbf a)=(\mathbf b\cdot\mathbf a)\mathbf c-(\mathbf b\cdot\mathbf c)\mathbf a,
\]
\[
\mathbf c\times(\mathbf a\times\mathbf b)=(\mathbf c\cdot\mathbf b)\mathbf a-(\mathbf c\cdot\mathbf a)\mathbf b,
\]
and they cancel to give \(\mathbf 0\).

\medskip
\noindent\textbf{Linear combination.}
Compute
\[
(\mathbf a\times\mathbf b)\times(\mathbf c\times\mathbf d)
=\mathbf c\,[(\mathbf a\times\mathbf b)\cdot\mathbf d]-\mathbf d\,[(\mathbf a\times\mathbf b)\cdot\mathbf c]
=\mathbf c\,[\mathbf a,\mathbf b,\mathbf d]-\mathbf d\,[\mathbf a,\mathbf b,\mathbf c].
\]
Also
\[
(\mathbf a\times\mathbf b)\times(\mathbf c\times\mathbf d)
=-(\mathbf c\times\mathbf d)\times(\mathbf a\times\mathbf b)
=\mathbf b\,[\mathbf c,\mathbf d,\mathbf a]-\mathbf a\,[\mathbf c,\mathbf d,\mathbf b].
\]
Using cyclic symmetry of \([\cdot,\cdot,\cdot]\) and equating gives
\[
[\mathbf b,\mathbf c,\mathbf d]\mathbf a-[\mathbf a,\mathbf c,\mathbf d]\mathbf b+[\mathbf a,\mathbf b,\mathbf d]\mathbf c-[\mathbf a,\mathbf b,\mathbf c]\mathbf d=\mathbf 0.
\]

\medskip
\noindent\textbf{(b)}
\[
[\mathbf a\times\mathbf b,\ \mathbf b\times\mathbf c,\ \mathbf c\times\mathbf a]
=(\mathbf a\times\mathbf b)\cdot\bigl((\mathbf b\times\mathbf c)\times(\mathbf c\times\mathbf a)\bigr).
\]
But
\[
(\mathbf b\times\mathbf c)\times(\mathbf c\times\mathbf a)
=\mathbf c\,[(\mathbf b\times\mathbf c)\cdot\mathbf a]-\mathbf a\,[(\mathbf b\times\mathbf c)\cdot\mathbf c]
=\mathbf c\,[\mathbf a,\mathbf b,\mathbf c],
\]
so
\[
[\mathbf a\times\mathbf b,\ \mathbf b\times\mathbf c,\ \mathbf c\times\mathbf a]
=[\mathbf a,\mathbf b,\mathbf c]\;(\mathbf a\times\mathbf b)\cdot\mathbf c
=[\mathbf a,\mathbf b,\mathbf c]^2.
\]
\end{solution}

\begin{problem}[CP-I-0097]
\label{prob:cp-i-0097}
\textbf{Geometric meaning of affine combinations (line and plane)}

\textbf{Exact statement.}
\hspace*{1.5em}Interpret geometrically the equations
\[
\mathbf r=(1-\lambda)\mathbf a+\lambda \mathbf b,
\qquad
\mathbf r=(1-\lambda-\mu)\mathbf a+\lambda \mathbf b+\mu \mathbf c,
\]
where \(\mathbf a,\mathbf b,\mathbf c\) are the position vectors of points in three dimensions and \(\lambda,\mu\) are scalar parameters
that take any real values. For each equation, obtain an equivalent form that does not involve parameters \(\lambda\) or \(\mu\).
How do your equations behave when the points \(\mathbf a,\mathbf b,\mathbf c\) are interchanged?

\end{problem}

\begin{solution}
\[
\mathbf r=(1-\lambda)\mathbf a+\lambda \mathbf b=\mathbf a+\lambda(\mathbf b-\mathbf a)
\]
describes the line through the points with position vectors \(\mathbf a\) and \(\mathbf b\).
A parameter-free form is
\[
(\mathbf r-\mathbf a)\times(\mathbf b-\mathbf a)=\mathbf 0.
\]
Similarly,
\[
\mathbf r=(1-\lambda-\mu)\mathbf a+\lambda \mathbf b+\mu \mathbf c
=\mathbf a+\lambda(\mathbf b-\mathbf a)+\mu(\mathbf c-\mathbf a)
\]
describes the plane through the points \(\mathbf a,\mathbf b,\mathbf c\).
A parameter-free form is
\[
(\mathbf r-\mathbf a)\cdot\bigl((\mathbf b-\mathbf a)\times(\mathbf c-\mathbf a)\bigr)=0.
\]
Interchanging \(\mathbf a,\mathbf b\) (or permuting \(\mathbf a,\mathbf b,\mathbf c\)) changes \(\lambda,\mu\) but not the underlying
line/plane; the parameter-free equations remain the same (up to an overall sign in the normal vector, which does not affect ``\(=0\)'').
\end{solution}

\begin{problem}[CP-I-0099]
\label{prob:cp-i-0099}
\textbf{Characterizing trace; (non)existence of \([\,\alpha,\beta\,]=id\)}

\textbf{Exact statement.}

	\par\hspace*{1.5em} (a)\quad Suppose \(f\in M_n(\mathbb R)^*\) satisfies \(f(AB)=f(BA)\) for all \(A,B\in M_n(\mathbb R)\) and \(f(I)=n\).
	Show that \(f(A)=\operatorname{tr}(A)\) for all \(A\).
	\par\noindent (b)\quad Let \(V\neq0\) be a finite-dimensional real vector space. Show there are no endomorphisms \(\alpha,\beta\) with
	\(\alpha\beta-\beta\alpha=id_V\).
	\par\noindent (c)\quad Let \(V=C^\infty(\mathbb R,\mathbb R)\). Find \(\alpha,\beta\in\operatorname{End}(V)\) such that
	\(\alpha\beta-\beta\alpha=id_V\).

\end{problem}

\begin{solution}
(a) Let \(E_{ij}\) be the matrix units. If \(i\neq j\), take \(A=E_{ii}\), \(B=E_{ij}\). Then \(AB=E_{ij}\), \(BA=0\), hence
\(f(E_{ij})=f(AB)=f(BA)=0\).
If \(i\neq j\), take \(A=E_{ij}\), \(B=E_{ji}\). Then \(AB=E_{ii}\), \(BA=E_{jj}\), hence \(f(E_{ii})=f(E_{jj})\).
Let \(c=f(E_{11})\). Then \(f(E_{ii})=c\) for all \(i\), so
\[
f(I)=f\Big(\sum_{i=1}^n E_{ii}\Big)=\sum_{i=1}^n f(E_{ii})=nc.
\]
Since \(f(I)=n\), we get \(c=1\).
For any \(A=\sum_{i,j} a_{ij}E_{ij}\),
\[
f(A)=\sum_{i,j} a_{ij}f(E_{ij})=\sum_i a_{ii}=\operatorname{tr}(A).
\]

(b) In finite dimension, \(\operatorname{tr}(\alpha\beta)=\operatorname{tr}(\beta\alpha)\), so
\(\operatorname{tr}(\alpha\beta-\beta\alpha)=0\).
If \(\alpha\beta-\beta\alpha=id_V\), then
\[
0=\operatorname{tr}(\alpha\beta-\beta\alpha)=\operatorname{tr}(id_V)=\dim V,
\]
contradiction.

(c) On \(V=C^\infty(\mathbb R)\), define \((\alpha f)(t)=f'(t)\) and \((\beta f)(t)=t f(t)\). Then
\[
(\alpha\beta f)(t)=(t f(t))'=f(t)+t f'(t),\qquad (\beta\alpha f)(t)=t f'(t),
\]
so \((\alpha\beta-\beta\alpha)f=f\), i.e. \(\alpha\beta-\beta\alpha=id_V\).
\end{solution}

\begin{problem}[CP-I-0100]
\label{prob:cp-i-0100}
\textbf{The subspace of maps sending $Y$ into $Z$}

\textbf{Exact statement.}
\hspace*{1.5em}Let $Y$ and $Z$ be subspaces of the finite dimensional vector spaces $V$ and $W$, respectively.
Show that
\[
R=\{\alpha\in \mathcal L(V,W):\ \alpha(Y)\subseteq Z\}
\]
is a subspace of $\mathcal L(V,W)$. What is the dimension of $R$?
\end{problem}

\begin{solution}
If $\alpha,\beta\in R$ and $c\in\mathbb F$, then for any $y\in Y$,
\[
(\alpha+\beta)(y)=\alpha(y)+\beta(y)\in Z,\qquad (c\alpha)(y)=c\,\alpha(y)\in Z,
\]
so $\alpha+\beta\in R$ and $c\alpha\in R$. Also $0\in R$. Hence $R\le \mathcal L(V,W)$.

Let $\dim V=n$, $\dim W=m$, $\dim Y=k$, $\dim Z=\ell$.
Choose bases $v_1,\dots,v_n$ of $V$ with $v_1,\dots,v_k$ a basis of $Y$, and $w_1,\dots,w_m$ of $W$ with
$w_1,\dots,w_\ell$ a basis of $Z$.
A map $\alpha$ is determined by the images of the $v_i$.
The condition $\alpha(Y)\subseteq Z$ forces $\alpha(v_i)\in Z$ for $i\le k$ (each gives $\ell$ degrees of freedom),
while $\alpha(v_i)\in W$ is arbitrary for $i>k$ (each gives $m$ degrees of freedom).
Therefore
\[
\dim R=k\ell+(n-k)m=nm-k(m-\ell).
\]

\bigskip
\end{solution}

\begin{problem}[CP-I-0104]
\label{prob:cp-i-0104}
\textbf{Equivalent characterizations of a direct sum}

\textbf{Exact statement.}
\hspace*{1.5em}Let $U_1,\dots,U_k$ be subspaces of a vector space $V$ and let $B_i$ be a basis for $U_i$.
Show that the following are equivalent:

	\par\noindent (i)\quad $U=\sum_{i=1}^k U_i$ is a direct sum, i.e.\ every element of $U$ can be written uniquely as $\sum_i u_i$ with $u_i\in U_i$.
	\par\noindent (ii)\quad $U_j\cap \sum_{i\ne j}U_i=\{0\}$ for all $j$.
	\par\noindent (iii)\quad The $B_i$ are pairwise disjoint and their union is a basis for $\sum_i U_i$.

Give an example where $U_i\cap U_j=\{0\}$ for all $i\ne j$, yet $U_1+\cdots+U_k$ is not a direct sum.
\end{problem}

\begin{solution}
(i)$\Rightarrow$(ii): If $x\in U_j\cap\sum_{i\ne j}U_i$, then $x=x+0+\cdots+0$ and also $x$ has a decomposition with
$j$-th component $0$. Uniqueness forces $x=0$.

(ii)$\Rightarrow$(i): If $\sum_i u_i=\sum_i u_i'$, then for each $j$,
\[
u_j-u_j'=\sum_{i\ne j}(u_i'-u_i)\in U_j\cap\sum_{i\ne j}U_i=\{0\},
\]
hence $u_j=u_j'$ for all $j$.

(iii)$\Rightarrow$(i): If $\bigcup_i B_i$ is a basis of $\sum_i U_i$, every vector has a unique linear combination of basis
vectors; grouping terms by $B_i$ yields a unique decomposition $\sum_i u_i$ with $u_i\in U_i$.

(ii)$\Rightarrow$(iii): The union $\bigcup_i B_i$ spans $\sum_i U_i$. For linear independence, suppose
$\sum_i \sum_{b\in B_i} c_b b=0$ and set $u_i=\sum_{b\in B_i} c_b b\in U_i$. Then $\sum_i u_i=0$, so
\[
u_j=-\sum_{i\ne j}u_i\in U_j\cap\sum_{i\ne j}U_i=\{0\},
\]
hence $u_j=0$ for all $j$ and therefore all $c_b=0$.

Counterexample: in $\mathbb R^2$ let
\[
U_1=\mathrm{span}(e_1),\quad U_2=\mathrm{span}(e_2),\quad U_3=\mathrm{span}(e_1+e_2).
\]
Then $U_i\cap U_j=\{0\}$ for $i\ne j$, but
\[
e_1+e_2=(e_1)+(e_2)+0=0+0+(e_1+e_2),
\]
so the sum is not direct.

\bigskip
\end{solution}

\begin{problem}[CP-I-0105]
\label{prob:cp-i-0105}
\textbf{Computing an Eigenspace as a Kernel}

\textbf{Canonical authored completion of the retained eigenspace fragment.}

\hspace*{1.5em}Let \(A\in M_n(\mathbb F)\) and let \(\lambda\in\mathbb F\).
Prove that
\[
E_\lambda(A)=\ker(A-\lambda I).
\]
Explain why \(\lambda\) is an eigenvalue exactly when this kernel is nonzero, and describe how
row reduction of \(A-\lambda I\) produces a basis of the eigenspace.
\end{problem}

\begin{solution}
By definition, \(x\neq0\) is a \(\lambda\)-eigenvector exactly when
\[
Ax=\lambda x.
\]
This is equivalent to
\[
(A-\lambda I)x=0.
\]
Including the zero vector, the full eigenspace is therefore
\[
E_\lambda(A)=\ker(A-\lambda I).
\]
Hence \(\lambda\) is an eigenvalue if and only if this kernel contains a nonzero vector, equivalently
if and only if \(A-\lambda I\) is singular.

To compute a basis, row-reduce \(A-\lambda I\), solve the resulting homogeneous system in
terms of its free variables, and take the coefficient vectors of those free parameters. These
vectors form a basis of \(\ker(A-\lambda I)\), hence of \(E_\lambda(A)\).
\end{solution}

\begin{problem}[CP-I-0119]
\label{prob:cp-i-0119}
\textbf{Projection and cross-product maps (unit vector $\mathbf n\in\mathbb R^3$)}

\textbf{Exact statement.}
\hspace*{1.5em}Let $\mathbf n$ be a unit vector in $\mathbb R^3$. Identify the image and kernel (null space) of each of the following linear maps $\mathbb R^3\to\mathbb R^3$:
\[
\text{(a)}\quad T:\ \mathbf x\mapsto \mathbf x'=\mathbf x-(\mathbf x\cdot\mathbf n)\mathbf n,
\qquad
\text{(b)}\quad Q:\ \mathbf x\mapsto \mathbf x'=\mathbf n\times \mathbf x.
\]
Show that $T^2=T$ and interpret the map $T$ geometrically. Interpret the maps $Q^2$ and $Q^3+Q$, and show that $Q^4=T$.

\end{problem}

\begin{solution}
For $T$, write $\mathbf x=\mathbf x_\perp+(\mathbf x\cdot\mathbf n)\mathbf n$ with $\mathbf x_\perp\cdot\mathbf n=0$. Then
\[
T(\mathbf x)=\mathbf x-(\mathbf x\cdot\mathbf n)\mathbf n=\mathbf x_\perp.
\]
Hence
\[
\ker T=\mathrm{span}\{\mathbf n\},\qquad \mathrm{Im}\,T=\mathbf n^\perp=\{\mathbf y:\mathbf y\cdot\mathbf n=0\}.
\]
Also $T(\mathbf x)\cdot\mathbf n=0$, so
\[
T^2(\mathbf x)=T(T(\mathbf x))=T(\mathbf x),
\]
i.e. $T^2=T$. Geometrically, $T$ is the orthogonal projection onto the plane $\mathbf n^\perp$.

For $Q(\mathbf x)=\mathbf n\times\mathbf x$, we have $\mathbf n\times\mathbf x=\mathbf 0$ iff $\mathbf x$ is parallel to $\mathbf n$, so
\[
\ker Q=\mathrm{span}\{\mathbf n\}.
\]
Also $\mathbf n\cdot(\mathbf n\times\mathbf x)=0$, so $\mathrm{Im}\,Q\subseteq \mathbf n^\perp$. Conversely, if $\mathbf y\in\mathbf n^\perp$, take $\mathbf x=-\mathbf n\times\mathbf y$ and use
\[
\mathbf n\times(\mathbf n\times\mathbf y)=\mathbf n(\mathbf n\cdot\mathbf y)-\mathbf y(\mathbf n\cdot\mathbf n)=-\mathbf y
\]
to get $Q(\mathbf x)=\mathbf y$. Thus $\mathrm{Im}\,Q=\mathbf n^\perp$.

Moreover,
\[
Q^2(\mathbf x)=\mathbf n\times(\mathbf n\times\mathbf x)=\mathbf n(\mathbf n\cdot\mathbf x)-\mathbf x
=-(\mathbf x-(\mathbf x\cdot\mathbf n)\mathbf n)=-T(\mathbf x),
\]
so $Q^2=-T$. Then
\[
Q^3=Q(Q^2)=Q(-T)=-QT.
\]
But $Q(\mathbf x)=Q(T\mathbf x)$ since $Q$ kills the component parallel to $\mathbf n$, hence $QT=Q$ and $Q^3=-Q$, i.e. $Q^3+Q=0$.
Finally,
\[
Q^4=(Q^2)^2=(-T)^2=T^2=T.
\]
\end{solution}

\begin{problem}[CP-I-0121]
\label{prob:cp-i-0121}
\textbf{Realizing any \(P\in GL_n(\mathbb F)\) as a product of change-of-basis matrices}

\textbf{Exact statement.}
\hspace*{1.5em}Let \(V\) be an \(n\)-dimensional vector space over a field \(\mathbb F\). Let \(B\) be a basis of \(V\), and let \(P\in GL_n(\mathbb F)\).
Prove there exist bases \(C\) and \(D\) of \(V\) such that
\[
[id_V]^{C}_{B}\,[id_V]^{B}_{D}=P.
\]
\end{problem}

\begin{solution}
Take \(C=B\). Then \([id_V]^C_B=[id_V]^B_B=I\), so it suffices to find a basis \(D\) with \([id_V]^B_D=P\).
Write \(B=(b_1,\dots,b_n)\) and let the columns of \(P\) be \(p^{(1)},\dots,p^{(n)}\in\mathbb F^n\).
Define
\[
d_j := \sum_{i=1}^n p_{ij}\,b_i \quad (j=1,\dots,n).
\]
Then \([d_j]_B=p^{(j)}\), hence \([id_V]^B_D=P\).
Since \(P\) is invertible, its columns are linearly independent, so \(d_1,\dots,d_n\) are linearly independent and \(D=(d_1,\dots,d_n)\) is a basis.
Therefore
\[
[id_V]^C_B[id_V]^B_D=I\cdot P=P.
\]
\end{solution}

\begin{problem}[CP-I-0122]
\label{prob:cp-i-0122}
\textbf{The composition operator $\Phi(\theta)=\beta\circ\theta\circ\alpha$}

\textbf{Exact statement.}
\hspace*{1.5em}Let $T,U,V,W$ be vector spaces over $\mathbb F$ and let $\alpha:T\to U$, $\beta:V\to W$ be fixed linear maps.
Show that
\[
\Phi:\mathcal L(U,V)\to \mathcal L(T,W),\qquad \Phi(\theta)=\beta\circ\theta\circ\alpha
\]
is linear. If the spaces are finite-dimensional and $\alpha$ and $\beta$ have ranks $r$ and $s$ respectively,
find the rank of $\Phi$.
\end{problem}

\begin{solution}
For $\theta_1,\theta_2\in\mathcal L(U,V)$ and $c\in\mathbb F$,
\[
\Phi(\theta_1+c\theta_2)
=\beta\circ(\theta_1+c\theta_2)\circ\alpha
=\beta\circ\theta_1\circ\alpha + c\,\beta\circ\theta_2\circ\alpha
=\Phi(\theta_1)+c\,\Phi(\theta_2),
\]
so $\Phi$ is linear.

Assume $\mathrm{rk}(\alpha)=r$ and $\mathrm{rk}(\beta)=s$.
Choose bases so that $\alpha$ and $\beta$ are represented by matrices
\[
[\alpha]=\begin{pmatrix}I_r&0\\0&0\end{pmatrix},
\qquad
[\beta]=\begin{pmatrix}I_s&0\\0&0\end{pmatrix}.
\]
For any $\theta\in\mathcal L(U,V)$ with matrix $[\theta]$, the matrix of $\Phi(\theta)=\beta\theta\alpha$
is obtained by taking the $s\times r$ block of $[\theta]$ (and placing it in the corresponding position),
and all other entries are forced to be $0$.
Thus the image of $\Phi$ is isomorphic to the space of all $s\times r$ matrices, so
\[
\mathrm{rk}(\Phi)=rs.
\]
\end{solution}

\begin{problem}[CP-I-0123]
\label{prob:cp-i-0123}
\textbf{Image and kernel in $\mathbb R^4$ depending on parameters $a,b$}

\textbf{Exact statement.}
\hspace*{1.5em}A linear map $\mathbb R^4\to\mathbb R^4$ is defined by $\mathbf x\mapsto M\mathbf x$, where
\[
M=\begin{pmatrix}
	a&a&b&a\\
	a&a&b&0\\
	a&b&a&b\\
	a&b&a&0
\end{pmatrix}.
\]
Find the image and kernel of this map for all real values of $a$ and $b$.

\end{problem}

\begin{solution}
\emph{Case 1: $a=b=0$.} Then $M=0$, so $\mathrm{Im}\,M=\{0\}$ and $\ker M=\mathbb R^4$.

\emph{Case 2: $a=b\neq 0$.} Then $\mathrm{rank}(M)=2$ and the equations $M\mathbf x=0$ reduce to
\[
x_4=0,\qquad x_1+x_2+x_3=0.
\]
Hence
\[
\ker M=\{(x_1,x_2,x_3,0):x_1+x_2+x_3=0\}
=\mathrm{span}\{(-1,1,0,0),(-1,0,1,0)\}.
\]
Also the columns satisfy $c_1=c_2=c_3=a(1,1,1,1)^T$ and $c_4=a(1,0,1,0)^T$, so
\[
\mathrm{Im}\,M=\mathrm{span}\{(1,1,1,1),(1,0,1,0)\}
=\{(u,v,u,v):u,v\in\mathbb R\}.
\]

\emph{Case 3: $a\neq b$.} Then $\mathrm{rank}(M)=3$ and $\dim\ker M=1$. One generator of the kernel is
\[
\mathbf k=(-(a+b),\,a,\,a,\,0),
\qquad\text{so}\qquad
\ker M=\mathrm{span}\{\mathbf k\}.
\]
Since $\dim(\mathrm{Im}\,M)=3$, the image is a hyperplane in $\mathbb R^4$. A vector orthogonal to all columns is
\[
\mathbf w=(b,-b,-a,a),
\]
so
\[
\mathrm{Im}\,M=\{\mathbf y\in\mathbb R^4:\mathbf w\cdot\mathbf y=0\}
=\{(y_1,y_2,y_3,y_4):\ b(y_1-y_2)+a(y_4-y_3)=0\}.
\]
\end{solution}

\section{Matrices and Operators}

\paragraph{Related material.} Volume I, Chapters \texttt{I/07--I/12}.

\begin{problem}[CP-I-0003]
\label{prob:cp-i-0003}
\textbf{Trace zero endomorphisms: zero diagonal form and commutators}

\textbf{Exact statement.}
\hspace*{1.5em}Let \(\alpha:V\to V\) be a linear map on a real finite-dimensional vector space \(V\) with \(\operatorname{tr}(\alpha)=0\).

	\par\noindent (i)\quad Show that if \(\alpha\neq 0\), there exists \(v\in V\) such that \(v,\alpha(v)\) are linearly independent.
	Deduce that there is a basis of \(V\) in which the matrix of \(\alpha\) has all diagonal entries equal to \(0\).
	\par\noindent (ii)\quad Show that there exist endomorphisms \(\beta,\gamma\in\operatorname{End}(V)\) such that \(\alpha=\beta\gamma-\gamma\beta\).

\end{problem}

\begin{solution}
(i) Suppose \(\alpha(v)\in \langle v\rangle\) for all \(v\neq0\). Then \(\alpha(v)=\lambda(v)v\).
For independent \(u,w\),
\[
\alpha(u+w)=\alpha(u)+\alpha(w)=\lambda(u)u+\lambda(w)w
\]
but also \(\alpha(u+w)=\lambda(u+w)(u+w)\). Comparing coefficients in the independent pair \(\{u,w\}\) yields
\(\lambda(u)=\lambda(w)=\lambda(u+w)\). Hence \(\lambda(v)\) is constant on \(V\setminus\{0\}\), so \(\alpha=\lambda I\).
Then \(0=\operatorname{tr}(\alpha)=\lambda\dim V\), hence \(\lambda=0\), contradicting \(\alpha\neq0\). Therefore some \(v\) has
\(v,\alpha(v)\) independent.

To get a zero-diagonal matrix, proceed by induction on \(n=\dim V\). If \(\alpha=0\) we are done. Otherwise choose \(v_1\) with
\(v_1,\alpha(v_1)\) independent. Choose \(f\in V^*\) with \(f(v_1)=1\) and \(f(\alpha(v_1))=0\), and set \(H=\ker f\).
Then \(\alpha(v_1)\in H\) and \(V=\langle v_1\rangle\oplus H\). Choose a basis \(v_2,\dots,v_n\) of \(H\).
In the basis \((v_1,v_2,\dots,v_n)\), the first diagonal entry is \(0\). Let \(\pi_H\) be the projection onto \(H\) along \(\langle v_1\rangle\),
and define \(\alpha_H:=\pi_H\circ \alpha|_H:H\to H\). The remaining diagonal entries are exactly those of \(\alpha_H\), and the block form
\[
\begin{pmatrix}
	0 & *\\
	* & [\alpha_H]
\end{pmatrix}
\]
shows \(\operatorname{tr}(\alpha_H)=\operatorname{tr}(\alpha)=0\). By induction, choose a basis of \(H\) for which \([\alpha_H]\) has zero diagonal.
Together with \(v_1\), this yields a basis of \(V\) in which \([\alpha]\) has all diagonal entries \(0\).

(ii) In a basis where \([\alpha]=A=(a_{ij})\) has \(a_{ii}=0\), pick distinct \(\lambda_1,\dots,\lambda_n\in\mathbb R\) and set
\(B=\mathrm{diag}(\lambda_1,\dots,\lambda_n)\). Define \(C=(c_{ij})\) by \(c_{ii}=0\) and
\[
c_{ij}=\frac{a_{ij}}{\lambda_i-\lambda_j}\quad (i\neq j).
\]
Then \((BC-CB)_{ij}=(\lambda_i-\lambda_j)c_{ij}=a_{ij}\) for \(i\neq j\), and the diagonal is \(0\). Hence \(BC-CB=A\),
so \(\alpha=\beta\gamma-\gamma\beta\) for the corresponding endomorphisms \(\beta,\gamma\).
\end{solution}

\begin{problem}[CP-I-0004]
\label{prob:cp-i-0004}
\textbf{Torus endomorphism induced by an integer matrix}

\hspace*{1.5em}Let \(A\) be a \(2\times 2\) matrix with entries in \(\mathbb Z\). Show that the linear map
\[
A:\mathbb R^2\to \mathbb R^2
\]
preserves the equivalence relation
\[
(a,b)\sim (a',b')
\quad\Longleftrightarrow\quad
(a-a',\,b-b')\in \mathbb Z^2,
\]
and therefore induces a continuous map
\[
f_A:T\to T
\]
on the torus
\[
T=\mathbb R^2/\mathbb Z^2.
\]
Compute the induced homomorphisms
\[
(f_A)_*:H_k(T)\to H_k(T)
\]
on the homology of \(T\).
\end{problem}

\begin{solution}

Recall that the torus is the quotient
\[
T=\mathbb R^2/\mathbb Z^2,
\]
where two points of \(\mathbb R^2\) are identified if they differ by an element of the lattice \(\mathbb Z^2\).

We first show that \(A\) respects this equivalence relation. Suppose
\[
(a,b)\sim (a',b').
\]
Then
\[
(a-a',\,b-b')\in \mathbb Z^2.
\]
Writing \(x=(a,b)\) and \(y=(a',b')\), this means \(x-y\in \mathbb Z^2\). Since \(A\) has integer entries, it sends integer vectors to integer vectors:
\[
A(\mathbb Z^2)\subseteq \mathbb Z^2.
\]
Hence
\[
A(x-y)=Ax-Ay\in \mathbb Z^2,
\]
so \(Ax\sim Ay\). Therefore \(A\) preserves the quotient relation and induces a well-defined map
\[
f_A:T\to T,
\qquad
f_A([x])=[Ax].
\]
Because the quotient map \(\pi:\mathbb R^2\to T\) is continuous and \(A\) is continuous, the induced map \(f_A\) is continuous.

We now compute the induced map on homology.

For the torus,
\[
H_0(T)\cong \mathbb Z,\qquad
H_1(T)\cong \mathbb Z^2,\qquad
H_2(T)\cong \mathbb Z,
\]
and \(H_k(T)=0\) for all \(k\ge 3\).

Choose the standard generators of \(H_1(T)\cong \mathbb Z^2\) represented by the two coordinate circles:
\[
\alpha(t)=[(t,0)],
\qquad
\beta(t)=[(0,t)],
\qquad t\in \mathbb R/\mathbb Z.
\]
Write
\[
A=
\begin{pmatrix}
	a & b\\
	c & d
\end{pmatrix}.
\]
Then
\[
A(t,0)=(at,ct),
\qquad
A(0,t)=(bt,dt).
\]
Therefore the induced loops satisfy
\[
f_A\circ \alpha \sim a\alpha+c\beta,
\qquad
f_A\circ \beta \sim b\alpha+d\beta.
\]
Hence, with respect to the basis \([\alpha],[\beta]\) of \(H_1(T)\),
\[
(f_A)_1=
\begin{pmatrix}
	a & b\\
	c & d
\end{pmatrix}
=A.
\]

Since \(H_2(T)\cong \mathbb Z\), the induced map on \(H_2\) must be multiplication by the degree of \(f_A\). For a torus map induced by the linear map \(A\), the degree is
\[
\deg(f_A)=\det(A).
\]
Thus
\[
(f_A)_2=\times \det(A).
\]

So the induced maps on homology are:
\[
(f_A)_*:H_0(T)\to H_0(T)
\quad\text{is}\quad
\mathrm{id}_{\mathbb Z},
\]
\[
(f_A)_*:H_1(T)\to H_1(T)
\quad\text{is}\quad
A,
\]
\[
(f_A)_*:H_2(T)\to H_2(T)
\quad\text{is multiplication by}\quad
\det(A),
\]
and for \(k\ge 3\),
\[
H_k(T)=0.
\]

Thus the complete answer is
\[
(f_A)_0=\mathrm{id},\qquad
(f_A)_1=A,\qquad
(f_A)_2=\times \det(A).
\]
\end{solution}

\begin{problem}[CP-I-0005]
\label{prob:cp-i-0005}
\textbf{Classifying Quadrics from the Eigenvalues of a Symmetric Matrix}

\textbf{Canonical authored completion of the retained quadric-classification fragment.}

\hspace*{1.5em}Let \(A=A^T\in M_3(\mathbb R)\), and consider the quadric
\[
Q=\{x\in\mathbb R^3:x^TAx=1\}.
\]
Using an orthonormal eigenbasis of \(A\), show that \(Q\) can be written
\[
\lambda_1u_1^2+\lambda_2u_2^2+\lambda_3u_3^2=1.
\]
Explain how the signs and zeros of the \(\lambda_i\) distinguish the basic cases
ellipsoid, cylinder, and hyperboloid. When all \(\lambda_i>0\), determine the semi-axis lengths.
\end{problem}

\begin{solution}
By the spectral theorem there is an orthogonal matrix \(P\) such that
\[
P^TAP=\operatorname{diag}(\lambda_1,\lambda_2,\lambda_3).
\]
Writing \(x=Pu\), we obtain
\[
x^TAx
=
u^TP^TAPu
=
\lambda_1u_1^2+\lambda_2u_2^2+\lambda_3u_3^2.
\]
Thus the quadric is diagonal in principal-axis coordinates.

If all three eigenvalues are positive, then
\[
\frac{u_1^2}{1/\lambda_1}
+\frac{u_2^2}{1/\lambda_2}
+\frac{u_3^2}{1/\lambda_3}
=1,
\]
so \(Q\) is an ellipsoid with semi-axis lengths
\[
\frac1{\sqrt{\lambda_1}},\qquad
\frac1{\sqrt{\lambda_2}},\qquad
\frac1{\sqrt{\lambda_3}}.
\]

If one eigenvalue is zero while the other two are positive, the corresponding coordinate is
unconstrained and the quadric is a cylinder. If the eigenvalues have mixed signs, the diagonal
equation is of hyperboloid type. Thus the signature and nullity of \(A\) control the basic
classification.
\end{solution}

\begin{problem}[CP-I-0007]
\label{prob:cp-i-0007}
\textbf{(Anosov/Lefschetz torus automorphism with explicit fixed points)}

\hspace*{1.5em}Let $T^2=\mathbb{R}^2/\mathbb{Z}^2$ and take
\[
A=\begin{pmatrix}3&2\\1&1\end{pmatrix}\in SL(2,\mathbb{Z}).
\]
Define the induced torus map
\[
f_A:T^2\to T^2,\qquad f_A([x])=[Ax].
\]

\medskip
\noindent\textbf{Hyperbolic (Anosov) property.}
We have $\mathrm{tr}(A)=4$ and $\det(A)=1$, hence the eigenvalues are
\[
\lambda_{\pm}=\frac{4\pm\sqrt{4^2-4}}{2}=2\pm\sqrt{3},
\]
so $A$ has one eigenvalue $>1$ and one eigenvalue $<1$.
Therefore $f_A$ is a hyperbolic (Anosov) torus automorphism.

\medskip
\noindent\textbf{Lefschetz condition.}
Since $Df_A\equiv A$, a fixed point is nondegenerate iff $\det(I-A)\neq 0$.
Compute
\[
A-I=\begin{pmatrix}2&2\\1&0\end{pmatrix},
\qquad
\det(A-I) = -2\neq 0,
\]
hence $\det(I-A)\neq 0$ and all fixed points are nondegenerate. Thus $f_A$ is Lefschetz.

\medskip
\noindent\textbf{Fixed points and their number.}
A point $[x]\in T^2$ is fixed iff $[Ax]=[x]$, i.e.
\[
(A-I)x\in \mathbb{Z}^2.
\]
For linear torus maps with $\det(A-I)\neq 0$ one has
\[
\#\mathrm{Fix}(f_A)=|\det(A-I)|.
\]
Hence $\#\mathrm{Fix}(f_A)=2$.

Moreover, writing $x=(x,y)\in\mathbb{R}^2$ we have
\[
(A-I)(x,y)=(2x+2y,\; x).
\]
Thus $(A-I)(x,y)\in\mathbb{Z}^2$ implies $x\in\mathbb{Z}$, so $x\equiv 0\ (\mathrm{mod}\ 1)$, and then
$2y\in\mathbb{Z}$, so $y\equiv 0$ or $1/2\ (\mathrm{mod}\ 1)$.
Therefore the fixed points are
\[
[(0,0)],\qquad \left[(0,\tfrac12)\right].
\]
\end{problem}

\begin{problem}[CP-I-0008]
\label{prob:cp-i-0008}
\textbf{Nonnegative Matrices: Explicit Perron--Frobenius Examples}

\hspace*{1.5em}Problem 11 only asks for existence of \emph{some} real eigenvalue $\lambda\ge 0$.

\noindent\textbf{Example 1.}
Let
\[
A=\begin{pmatrix}2&1\\1&2\end{pmatrix}\ge 0.
\]
Then
\[
\det(A-\lambda I)=\det\begin{pmatrix}2-\lambda&1\\1&2-\lambda\end{pmatrix}
=(2-\lambda)^2-1=\lambda^2-4\lambda+3=(\lambda-3)(\lambda-1).
\]
Hence the eigenvalues are $\lambda=3$ and $\lambda=1$, both nonnegative.
For $\lambda=3$ one may take $v=(1,1)$, since
\[
Av=(3,3)=3(1,1).
\]
Moreover, with $x=(\tfrac12,\tfrac12)\in\Delta$ one has $Ax=(\tfrac32,\tfrac32)$ and
\[
F(x)=\frac{Ax}{\sum_i (Ax)_i}=\frac{(\tfrac32,\tfrac32)}{3}=\Bigl(\tfrac12,\tfrac12\Bigr)=x,
\]
so $x$ is a Brouwer fixed point and produces the eigenvalue $3$.

\medskip
\noindent\textbf{Example 2.}
Let
\[
A=\begin{pmatrix}0&1\\1&0\end{pmatrix}\ge 0.
\]
Then
\[
\det(A-\lambda I)=\det\begin{pmatrix}-\lambda&1\\1&-\lambda\end{pmatrix}=\lambda^2-1,
\]
so $\lambda=1,-1$. The nonnegative eigenvalue is $\lambda=1$ with eigenvector $v=(1,1)$, since
\[
Av=(1,1)=1\cdot(1,1).
\]

\medskip
\noindent\textbf{Example 3 (zero column).}
Let
\[
A=\begin{pmatrix}1&0\\2&0\end{pmatrix}\ge 0.
\]
The second column is zero, hence $Ae_2=0$, so $\lambda=0$ is an eigenvalue with eigenvector $e_2=(0,1)\ge 0$.
(Indeed,
\[
\det(A-\lambda I)=\det\begin{pmatrix}1-\lambda&0\\2&-\lambda\end{pmatrix}=-\lambda(1-\lambda),
\]
so the eigenvalues are $0$ and $1$.)

\noindent\textbf{Example.}
Let
\[
A=\begin{pmatrix}
	1&1&0\\
	0&1&1\\
	1&0&1
\end{pmatrix}\ge 0.
\]
Take $v=(1,1,1)^T$. Then
\[
Av=
\begin{pmatrix}
	1+1+0\\
	0+1+1\\
	1+0+1
\end{pmatrix}
=
\begin{pmatrix}
	2\\2\\2
\end{pmatrix}
=2v.
\]
Hence $A$ has the nonnegative real eigenvalue $\lambda=2$ with a strictly positive eigenvector $v>0$.

\medskip
\noindent\textbf{Brouwer normalization viewpoint.}
Let $\Delta=\{x\in\mathbb{R}^3:\ x_i\ge 0,\ \sum_i x_i=1\}$ and define
\[
F(x)=\frac{Ax}{\sum_{i=1}^3 (Ax)_i}.
\]
For $x=(\tfrac13,\tfrac13,\tfrac13)\in\Delta$ we have
\[
Ax=\tfrac13 Av=\Bigl(\tfrac23,\tfrac23,\tfrac23\Bigr),\qquad \sum_i (Ax)_i = 2,
\]
so
\[
F(x)=\frac{(\tfrac23,\tfrac23,\tfrac23)}{2}=\Bigl(\tfrac13,\tfrac13,\tfrac13\Bigr)=x.
\]
Thus $x$ is a fixed point of $F$, yielding the eigenpair $Ax=\lambda x$ with $\lambda=2\ge 0$.
\end{problem}

\begin{problem}[CP-I-0009]
\label{prob:cp-i-0009}
\textbf{Dual bases under basis changes in dimension 4}

\textbf{Exact statement.}
\hspace*{1.5em}Let \(V\) be 4-dimensional over \(\mathbb R\), and let \(\{\xi_1,\xi_2,\xi_3,\xi_4\}\) be the dual basis to \(\{x_1,x_2,x_3,x_4\}\).
Find the dual bases to:

	\par\noindent (a)\quad \(\{x_2,x_1,x_4,x_3\}\),
	\par\noindent (b)\quad \(\{x_1,2x_2,\tfrac12 x_3,x_4\}\),
	\par\noindent (c)\quad \(\{x_1+x_2,\;x_2+x_3,\;x_3+x_4,\;x_4\}\),
	\par\noindent (d)\quad \(\{x_1,\;x_2-x_1,\;x_3-x_2+x_1,\;x_4-x_3+x_2-x_1\}\).

\end{problem}

\begin{solution}
If \(y_j=\sum_i p_{ij}x_i\) and \(P=(p_{ij})\), then the dual basis \(\eta\) satisfies
\[
\eta_i=\sum_{k=1}^4 (P^{-1})_{ik}\,\xi_k.
\]
(a) \((\eta_1,\eta_2,\eta_3,\eta_4)=(\xi_2,\xi_1,\xi_4,\xi_3).\)

(b) \((\eta_1,\eta_2,\eta_3,\eta_4)=\big(\xi_1,\tfrac12\xi_2,2\xi_3,\xi_4\big).\)

(c)
\[
\eta_1=\xi_1,\quad
\eta_2=\xi_2-\xi_1,\quad
\eta_3=\xi_1-\xi_2+\xi_3,\quad
\eta_4=-\xi_1+\xi_2-\xi_3+\xi_4.
\]

(d)
\[
\eta_1=\xi_1+\xi_2,\quad
\eta_2=\xi_2+\xi_3,\quad
\eta_3=\xi_3+\xi_4,\quad
\eta_4=\xi_4.
\]
\end{solution}

\begin{problem}[CP-I-0010]
\label{prob:cp-i-0010}
\textbf{A linear system depending on a parameter \(t\)}

\textbf{Exact statement.}
\hspace*{1.5em}For each real value of \(t\), determine whether or not there exist solutions to the simultaneous equations
\[
\begin{cases}
	x+y+z=t,\\
	tx+2z=3,\\
	3x+ty+5z=7,
\end{cases}
\]
exhibiting the most general form of such solutions when they exist.

\end{problem}

\begin{solution}
Let
\[
A(t)=\begin{pmatrix}
	1&1&1\\
	t&0&2\\
	3&t&5
\end{pmatrix},\qquad
b(t)=\begin{pmatrix}t\\3\\7\end{pmatrix}.
\]
Then
\[
\det A(t)=
\begin{vmatrix}
	1&1&1\\
	t&0&2\\
	3&t&5
\end{vmatrix}
=t^2-7t+6=(t-1)(t-6).
\]

\textbf{Case 1: \(t\neq 1,6\).}
Then \(\det A(t)\neq 0\), so there is a unique solution. Solving (e.g.\ by elimination) gives
\[
x=\frac{1-2t}{t-6},\qquad
y=\frac{8-5t}{t-6},\qquad
z=\frac{t^2+t-9}{t-6}.
\]

\textbf{Case 2: \(t=1\).}
The system becomes
\[
\begin{cases}
	x+y+z=1,\\
	x+2z=3,\\
	3x+y+5z=7.
\end{cases}
\]
From \(x+2z=3\) we get \(x=3-2z\). Substituting into \(3x+y+5z=7\) yields
\[
3(3-2z)+y+5z=7 \ \Rightarrow\ y-z=-2 \ \Rightarrow\ y=z-2.
\]
Thus, writing \(z=s\),
\[
(x,y,z)=(3-2s,\ s-2,\ s),\qquad s\in\mathbb R,
\]
equivalently \((x,y,z)=(3,-2,0)+s(-2,1,1)\).

\textbf{Case 3: \(t=6\).}
The system becomes
\[
\begin{cases}
	x+y+z=6,\\
	6x+2z=3,\\
	3x+6y+5z=7.
\end{cases}
\]
From \(6x+2z=3\) we get \(3x+z=\tfrac{3}{2}\).
Also,
\[
(3x+6y+5z)-3(x+y+z)=7-18 \ \Rightarrow\ 3y+2z=-11.
\]
Using \(y=6-x-z\) gives
\[
3(6-x-z)+2z=-11 \ \Rightarrow\ 18-3x-z=-11 \ \Rightarrow\ 3x+z=29,
\]
contradicting \(3x+z=\tfrac{3}{2}\). Hence there are no solutions when \(t=6\).
\end{solution}

\begin{problem}[CP-I-0011]
\label{prob:cp-i-0011}
\textbf{Completing an orthogonal matrix}

\textbf{Exact statement.}
\hspace*{1.5em}Find $a,b,c$ such that
\[
Q=\begin{pmatrix}
	\frac13&0&a\\[2pt]
	\frac23&\frac1{\sqrt2}&b\\[2pt]
	\frac23&-\frac1{\sqrt2}&c
\end{pmatrix}
\]
is an orthogonal matrix. Does this condition determine $a,b,c$ uniquely?

\end{problem}

\begin{solution}
Let the columns be
\[
v_1=\begin{pmatrix}\frac13\\\frac23\\\frac23\end{pmatrix},\qquad
v_2=\begin{pmatrix}0\\\frac1{\sqrt2}\\-\frac1{\sqrt2}\end{pmatrix},\qquad
v_3=\begin{pmatrix}a\\b\\c\end{pmatrix}.
\]
One checks
\[
\|v_1\|^2=\frac19+\frac49+\frac49=1,\qquad
\|v_2\|^2=\frac12+\frac12=1,\qquad
v_1\cdot v_2=0,
\]
so $v_1,v_2$ are already orthonormal.
Thus $v_3$ must be a unit vector orthogonal to both $v_1$ and $v_2$, hence
\[
v_3=\pm(v_1\times v_2).
\]
Compute
\[
v_1\times v_2=\frac{\sqrt2}{6}\begin{pmatrix}-4\\1\\1\end{pmatrix}.
\]
Therefore
\[
(a,b,c)=\left(-\frac{2\sqrt2}{3},\,\frac{\sqrt2}{6},\,\frac{\sqrt2}{6}\right)
\quad\text{or}\quad
(a,b,c)=\left(\frac{2\sqrt2}{3},\,-\frac{\sqrt2}{6},\,-\frac{\sqrt2}{6}\right).
\]
So the condition does not determine $a,b,c$ uniquely: there are exactly two choices (sign $\pm$).
\end{solution}

\begin{problem}[CP-I-0012]
\label{prob:cp-i-0012}
\textbf{Diagonalise a quadric; minimum distance to the origin}

\textbf{Exact statement.}
\hspace*{1.5em}Let $\Sigma\subset\mathbb R^3$ be given by
\[
2x^2+2xy+4yz+z^2=1.
\]
By considering a suitable real symmetric matrix, show that there is a new orthonormal basis with associated coordinates $u,v,w$ such that
\[
\Sigma:\ \lambda u^2+\mu v^2+\nu w^2=1,
\]
for constants $\lambda,\mu,\nu$ to be determined. Find the minimum distance from a point on $\Sigma$ to the origin.

\end{problem}

\begin{solution}
Write the quadratic form as $(x,y,z)M(x,y,z)^T$ with
\[
M=\begin{pmatrix}
	2&1&0\\
	1&0&2\\
	0&2&1
\end{pmatrix},
\]
which is real symmetric. By the spectral theorem there is an orthogonal matrix $P$ such that
\[
P^T M P=\operatorname{diag}(\lambda,\mu,\nu).
\]
The eigenvalues are
\[
\lambda=3,\qquad \mu=\sqrt3,\qquad \nu=-\sqrt3,
\]
so in orthonormal coordinates $(u,v,w)$ one has
\[
3u^2+\sqrt3\,v^2-\sqrt3\,w^2=1.
\]

To minimize the distance to the origin, minimize $u^2+v^2+w^2$ subject to
$3u^2+\sqrt3 v^2-\sqrt3 w^2=1$. Using Lagrange multipliers,
the minimum occurs for $v=w=0$ and $3u^2=1$, hence $u^2=1/3$ and
\[
d_{\min}=\sqrt{u^2}=\frac1{\sqrt3}.
\]
(Indeed, the closest points lie on the eigendirection of eigenvalue $3$.)
\end{solution}

\begin{problem}[CP-I-0013]
\label{prob:cp-i-0013}
\textbf{Rodrigues' rotation formula and extracting $\theta,\mathbf n$ from $R$}

\textbf{Exact statement.}
\hspace*{1.5em}The linear map $\mathbb R^3\to\mathbb R^3$ defined by
\[
\mathbf x\mapsto \mathbf x'=\cos\theta\,\mathbf x+(\mathbf x\cdot\mathbf n)(1-\cos\theta)\,\mathbf n-\sin\theta\,(\mathbf x\times\mathbf n)
\]
is a rotation by angle $\theta$ in a positive sense about the unit vector $\mathbf n$.
Check this in the case $\mathbf n=(0,0,1)$.

Show that the expression above can be written as $\mathbf x'=R\mathbf x$, where $R$ is a matrix with entries $R_{ij}$ found explicitly (in terms of $\theta,n_i,\delta_{ij},\varepsilon_{ijk}$). Hence show that
\[
R_{ii}=2\cos\theta+1,\qquad \varepsilon_{ijk}R_{jk}=-2n_i\sin\theta.
\]
Determine $\theta$ and $\mathbf n$ for
\[
R=\frac13\begin{pmatrix}
	2&-1&2\\
	2&2&-1\\
	-1&2&2
\end{pmatrix}.
\]

\end{problem}

\begin{solution}
For $\mathbf n=(0,0,1)$ and $\mathbf x=(x,y,z)$, we have $\mathbf x\cdot\mathbf n=z$ and
\[
\mathbf x\times\mathbf n=(y,-x,0),
\]
so
\[
\mathbf x'=(x\cos\theta-y\sin\theta,\ x\sin\theta+y\cos\theta,\ z),
\]
the standard rotation about the $z$-axis.

In components,
\[
x'_i=\cos\theta\,x_i+(1-\cos\theta)(n_jx_j)n_i-\sin\theta\,\varepsilon_{ijk}x_jn_k,
\]
so $x'_i=R_{ij}x_j$ with
\[
R_{ij}=\cos\theta\,\delta_{ij}+(1-\cos\theta)n_in_j-\sin\theta\,\varepsilon_{ijk}n_k.
\]
Taking the trace gives
\[
R_{ii}=\cos\theta\,\delta_{ii}+(1-\cos\theta)n_in_i-\sin\theta\,\varepsilon_{iik}n_k
=3\cos\theta+(1-\cos\theta)\cdot 1
=1+2\cos\theta.
\]
Also,
\[
\varepsilon_{ijk}R_{jk}
=-\sin\theta\,\varepsilon_{ijk}\varepsilon_{jkm}n_m
=-\sin\theta\,(2\delta_{im})n_m
=-2n_i\sin\theta.
\]

For the given matrix, $\mathrm{tr}(R)=\frac13(2+2+2)=2$, hence
\[
1+2\cos\theta=2\quad\Rightarrow\quad \cos\theta=\frac12\quad\Rightarrow\quad \theta=\frac{\pi}{3}.
\]
Next,
\[
(\varepsilon_{ijk}R_{jk})_{i}=(R_{32}-R_{23},\,R_{13}-R_{31},\,R_{21}-R_{12})=(1,1,1).
\]
Since $\varepsilon_{ijk}R_{jk}=-2\mathbf n\sin\theta$ and $\sin(\pi/3)=\sqrt3/2$, we get
\[
(1,1,1)=-(\sqrt3)\,\mathbf n
\quad\Rightarrow\quad
\mathbf n=-\frac{1}{\sqrt3}(1,1,1).
\]
\end{solution}

\begin{problem}[CP-I-0015]
\label{prob:cp-i-0015}
\textbf{Triangular Matrices and Their Spectrum}

\textbf{Canonical authored completion of the retained triangular-matrix fragment.}

\hspace*{1.5em}Let \(A=(a_{ij})\in M_n(\mathbb F)\) be upper triangular.
Prove that
\[
\chi_A(t)=\det(A-tI)=\prod_{i=1}^{n}(a_{ii}-t).
\]
Deduce that, over a field in which \(\chi_A\) splits, the eigenvalues of \(A\), counted with
algebraic multiplicity, are precisely its diagonal entries. Show that the same conclusion holds
for lower triangular matrices.
\end{problem}

\begin{solution}
The matrix \(A-tI\) is upper triangular, with diagonal entries
\[
a_{11}-t,\dots,a_{nn}-t.
\]
The determinant of a triangular matrix is the product of its diagonal entries, hence
\[
\chi_A(t)=\det(A-tI)=\prod_{i=1}^{n}(a_{ii}-t).
\]
Therefore the roots of the characteristic polynomial are exactly
\[
t=a_{11},\dots,a_{nn},
\]
with repetitions matching their occurrences on the diagonal. These are precisely the eigenvalues
when the characteristic polynomial splits. The proof for lower triangular matrices is identical,
because \(A-tI\) is again triangular with the same diagonal entries.
\end{solution}

\begin{problem}[CP-I-0017]
\label{prob:cp-i-0017}
\textbf{Jordan Blocks and Algebraic Multiplicity}

\textbf{Canonical authored completion of the retained Jordan-multiplicity fragment.}

\hspace*{1.5em}Suppose \(A\in M_n(\mathbb C)\) is in Jordan normal form.
For an eigenvalue \(\lambda\), prove that its algebraic multiplicity equals the sum of the sizes
of all Jordan blocks belonging to \(\lambda\). Prove also that
\[
\dim E_\lambda(A)
\]
equals the number of Jordan blocks belonging to \(\lambda\).
\end{problem}

\begin{solution}
Write the Jordan form as a direct sum of blocks. A block \(J_m(\lambda)\) contributes
\[
(t-\lambda)^m
\]
to the characteristic polynomial, because
\[
\det(tI-J_m(\lambda))=(t-\lambda)^m.
\]
Hence, if the blocks for \(\lambda\) have sizes \(m_1,\dots,m_r\), their total contribution is
\[
(t-\lambda)^{m_1+\cdots+m_r}.
\]
Thus the algebraic multiplicity of \(\lambda\) is
\[
m_1+\cdots+m_r.
\]

For one Jordan block \(J_m(\lambda)\),
\[
\ker(J_m(\lambda)-\lambda I)
\]
is one-dimensional. Since kernels add across direct sums, \(r\) Jordan blocks for \(\lambda\)
contribute \(r\) independent eigenvectors. Therefore
\[
\dim E_\lambda(A)=r,
\]
the number of Jordan blocks for \(\lambda\).
\end{solution}

\begin{problem}[CP-I-0020]
\label{prob:cp-i-0020}
\textbf{Determinant factorization by row operations}

\textbf{Exact statement.}
\hspace*{1.5em}For the real matrix
\[
M=\begin{pmatrix}
	a & a^2 & bc\\
	b & b^2 & ca\\
	c & c^2 & ab
\end{pmatrix},
\]
show (using row operations) that
\[
\det M=(a-b)(b-c)(c-a)(ab+bc+ca).
\]

\end{problem}

\begin{solution}
Let the rows be \(R_1,R_2,R_3\). Perform
\[
R_2\leftarrow R_2-R_1,\qquad R_3\leftarrow R_3-R_1.
\]
Then
\[
R_2-R_1=(b-a,\ b^2-a^2,\ ca-bc)=(b-a)(1,\ a+b,\ -c),
\]
\[
R_3-R_1=(c-a,\ c^2-a^2,\ ab-bc)=(c-a)(1,\ a+c,\ -b).
\]
Hence
\[
\det M=(b-a)(c-a)\det\begin{pmatrix}
	a & a^2 & bc\\
	1 & a+b & -c\\
	1 & a+c & -b
\end{pmatrix}.
\]
Now do \(R_3\leftarrow R_3-R_2\):
\[
(1,a+c,-b)-(1,a+b,-c)=(0,c-b,c-b)=(c-b)(0,1,1),
\]
so
\[
\det M=(b-a)(c-a)(c-b)\det\begin{pmatrix}
	a & a^2 & bc\\
	1 & a+b & -c\\
	0 & 1 & 1
\end{pmatrix}.
\]
Compute the remaining determinant by expansion along the first row:
\[
\det\begin{pmatrix}
	a & a^2 & bc\\
	1 & a+b & -c\\
	0 & 1 & 1
\end{pmatrix}
=
a\bigl((a+b)\cdot 1-(-c)\cdot 1\bigr)-a^2\cdot 1+bc\cdot 1
\]
\[
=a(a+b+c)-a^2+bc=a(b+c)+bc=ab+ac+bc.
\]
Therefore
\[
\det M=(b-a)(c-a)(c-b)(ab+bc+ca)=(a-b)(b-c)(c-a)(ab+bc+ca).
\]
\end{solution}

\begin{problem}[CP-I-0022]
\label{prob:cp-i-0022}
\textbf{Eigenvalues of an upper triangular matrix}

\textbf{Exact statement.}
\hspace*{1.5em}A square matrix $A$ is \emph{upper triangular} if $A_{ij}=0$ for $i>j$.
Show that the eigenvalues of such a matrix are its diagonal entries:
\[
\lambda_i=A_{ii}\qquad(\text{no sum over }i).
\]

\end{problem}

\begin{solution}
The characteristic polynomial is
\[
p_A(t)=\det(A-tI).
\]
If $A$ is upper triangular, then $A-tI$ is also upper triangular with diagonal entries
$A_{11}-t,\dots,A_{nn}-t$. The determinant of a triangular matrix equals the product
of its diagonal entries, hence
\[
p_A(t)=\prod_{i=1}^n (A_{ii}-t).
\]
Therefore the roots are $t=A_{ii}$, so the eigenvalues are exactly the diagonal entries.
\end{solution}

\begin{problem}[CP-I-0023]
\label{prob:cp-i-0023}
\textbf{Repeated Eigenvalues and the Jordan Obstruction}

\textbf{Canonical authored completion of the retained repeated-eigenvalue fragment.}

\hspace*{1.5em}Let \(A\in M_n(\mathbb C)\) have characteristic polynomial
\[
\chi_A(t)=(t-\lambda)^n.
\]
Prove that \(A\) is diagonalizable if and only if
\[
\dim\ker(A-\lambda I)=n.
\]
Show equivalently that \(A\) is diagonalizable if and only if \(A=\lambda I\). Explain the
Jordan-form obstruction when the eigenspace has dimension \(<n\).
\end{problem}

\begin{solution}
Since \(\lambda\) is the only eigenvalue, a diagonalizable \(A\) would have a basis consisting
entirely of \(\lambda\)-eigenvectors. Thus
\[
A=\lambda I,
\]
and consequently
\[
\ker(A-\lambda I)=\mathbb C^n.
\]
Hence diagonalizability implies
\[
\dim\ker(A-\lambda I)=n.
\]

Conversely, if the kernel has dimension \(n\), then \(A-\lambda I=0\), so \(A=\lambda I\),
which is diagonal.

In Jordan form, every block belongs to \(\lambda\). The matrix is diagonalizable exactly when
all Jordan blocks have size \(1\). If
\[
\dim\ker(A-\lambda I)<n,
\]
at least one block has size \(>1\), and this nontrivial Jordan block prevents diagonalizability.
\end{solution}

\begin{problem}[CP-I-0024]
\label{prob:cp-i-0024}
\textbf{Minimum Distance from a Positive-Definite Quadric to the Origin}

\textbf{Canonical authored completion of the retained minimum-distance fragment.}

\hspace*{1.5em}Let \(A=A^T\in M_3(\mathbb R)\) be positive definite, with eigenvalues
\[
0<\lambda_1\le \lambda_2\le \lambda_3.
\]
For the ellipsoid
\[
Q=\{x\in\mathbb R^3:x^TAx=1\},
\]
prove that the minimum distance from \(Q\) to the origin is
\[
r_{\min}=\frac1{\sqrt{\lambda_3}}
=\frac1{\sqrt{\lambda_{\max}}}.
\]
Identify the points where the minimum is attained.
\end{problem}

\begin{solution}
Choose an orthonormal eigenbasis and write \(x=Pu\). Then
\[
x^TAx=\lambda_1u_1^2+\lambda_2u_2^2+\lambda_3u_3^2=1,
\]
while orthogonality of \(P\) gives
\[
\|x\|^2=\|u\|^2=u_1^2+u_2^2+u_3^2.
\]

Since \(\lambda_i\le\lambda_3\),
\[
1=\sum_{i=1}^3\lambda_i u_i^2
\le
\lambda_3\sum_{i=1}^3u_i^2
=
\lambda_3\|x\|^2.
\]
Hence
\[
\|x\|^2\ge\frac1{\lambda_3},
\qquad
\|x\|\ge\frac1{\sqrt{\lambda_3}}.
\]

Equality holds exactly when \(u_i=0\) for every eigendirection whose eigenvalue is strictly less
than \(\lambda_3\). In particular, if \(\lambda_3\) is simple, the closest points are
\[
x=\pm\frac1{\sqrt{\lambda_3}}\,v_3,
\]
where \(v_3\) is a unit eigenvector for \(\lambda_3\).
\end{solution}

\begin{problem}[CP-I-0025]
\label{prob:cp-i-0025}
\textbf{A circulant determinant and its factorization}

\textbf{Exact statement.}
\hspace*{1.5em}Show by direct evaluation that
\[
\Delta(x,y,z)\equiv
\begin{vmatrix}
	x&y&z\\
	z&x&y\\
	y&z&x
\end{vmatrix}
= x^3+y^3+z^3-3xyz.
\]
Now show that \(x+y+z\), \(x+\omega y+\omega^2 z\), \(x+\omega^2 y+\omega z\)
are factors of \(\Delta\), where \(\omega^3=1\), \(\omega\neq 1\).

\end{problem}

\begin{solution}
\textbf{Direct evaluation.} Expanding along the first row,
\[
\Delta=
x\begin{vmatrix}x&y\\ z&x\end{vmatrix}
- y\begin{vmatrix}z&y\\ y&x\end{vmatrix}
+ z\begin{vmatrix}z&x\\ y&z\end{vmatrix}
=
x(x^2-yz)-y(zx-y^2)+z(z^2-xy),
\]
so
\[
\Delta=x^3+y^3+z^3-3xyz.
\]

\textbf{Factors via eigenvectors.}
Let
\[
C=\begin{pmatrix}
	x&y&z\\
	z&x&y\\
	y&z&x
\end{pmatrix}.
\]
Using \(\omega^3=1\) and \(1+\omega+\omega^2=0\), check:
\[
C\begin{pmatrix}1\\1\\1\end{pmatrix}=(x+y+z)\begin{pmatrix}1\\1\\1\end{pmatrix},\quad
C\begin{pmatrix}1\\\omega\\\omega^2\end{pmatrix}=(x+\omega y+\omega^2 z)\begin{pmatrix}1\\\omega\\\omega^2\end{pmatrix},
\]
\[
C\begin{pmatrix}1\\\omega^2\\\omega\end{pmatrix}=(x+\omega^2 y+\omega z)\begin{pmatrix}1\\\omega^2\\\omega\end{pmatrix}.
\]
Hence \(\det(C)=0\) whenever any of the three linear factors vanishes, so each factor divides \(\Delta\).
Since \(\Delta\) is homogeneous of degree \(3\), we must have
\[
\Delta=c\,(x+y+z)(x+\omega y+\omega^2 z)(x+\omega^2 y+\omega z)
\]
for a constant \(c\). Setting \(y=z=0\) gives \(\Delta(x,0,0)=x^3\) and the product equals \(x^3\), so \(c=1\).
\end{solution}

\begin{problem}[CP-I-0027]
\label{prob:cp-i-0027}
\textbf{Determinant of an odd-dimensional antisymmetric matrix}

\textbf{Exact statement.}
\hspace*{1.5em}If \(A\) is a \((2n+1)\times(2n+1)\) antisymmetric matrix, find \(\det A\).

\end{problem}

\begin{solution}
Antisymmetric means \(A^T=-A\). Then
\[
\det(A)=\det(A^T)=\det(-A)=(-1)^{2n+1}\det(A)=-\det(A),
\]
so \(\det(A)=0\).
\end{solution}

\begin{problem}[CP-I-0028]
\label{prob:cp-i-0028}
\textbf{Hadamard's inequality (entrywise bound version)}

\textbf{Exact statement.}
\hspace*{1.5em}Prove Hadamard's Inequality: if $A$ is a real $n\times n$ matrix and $k>0$ satisfies $|a_{ij}|\le k$ for all $1\le i,j\le n$, then
\[
|\det(A)|\le k^n\,n^{n/2}.
\]
\end{problem}

\begin{solution}
Let $r_1,\dots,r_n\in\mathbb R^n$ be the row vectors of $A$.
Hadamard's inequality states that
\[
|\det(A)|\le \prod_{i=1}^n \|r_i\|_2,
\]
where $\|\cdot\|_2$ is the Euclidean norm.

Now $|a_{ij}|\le k$ implies
\[
\|r_i\|_2^2=\sum_{j=1}^n a_{ij}^2 \le \sum_{j=1}^n k^2 = nk^2,
\]
so $\|r_i\|_2\le \sqrt n\,k$ for every $i$.
Therefore
\[
|\det(A)|
\le \prod_{i=1}^n (\sqrt n\,k)
=(\sqrt n\,k)^n
=k^n n^{n/2}.
\]
\end{solution}

\begin{problem}[CP-I-0029]
\label{prob:cp-i-0029}
\textbf{a non-degenerate critical point (a saddle)}

\hspace*{1.5em}Consider
	\[
	g(u,v)=uv.
	\]
	\noindent\textbf{Step 1: critical point.}
	The gradient is
	\[
	\nabla g(u,v)=(g_u(u,v),g_v(u,v))=(v,u).
	\]
	Thus $\nabla g(0,0)=(0,0)$, so $(0,0)$ is a critical point.

	\smallskip
	\noindent\textbf{Step 2: non-degeneracy.}
	The second derivatives are
	\[
	g_{uu}=0,\qquad g_{uv}=1,\qquad g_{vv}=0,
	\]
	so the Hessian matrix is constant:
	\[
	H_g=
	\begin{pmatrix}
		0 & 1\\
		1 & 0
	\end{pmatrix},
	\qquad
	\det(H_g)=-1\neq 0.
	\]
	Therefore $(0,0)$ is a \emph{non-degenerate} critical point (in fact, a saddle point). This is the local model
	that appears in the case $n=2$ for the determinant function on $M(2)$ (restricted to diagonal matrices).

	\noindent
	Let $f:\mathbb{R}^k\to\mathbb{R}$ be smooth and let $x_0\in\mathbb{R}^k$.
	The condition that $x_0$ is a critical point is
	\[
	\nabla f(x_0)=0
	\qquad\Longleftrightarrow\qquad
	df_{x_0}=0.
	\]
	Geometrically, this means that the \emph{first-order} (linear) change of $f$ at $x_0$ vanishes in every direction.

	\medskip
	\noindent\textbf{(1) The tangent approximation.}
	For a small vector $h\in\mathbb{R}^k$, Taylor's formula gives
	\[
	f(x_0+h)=f(x_0)+\nabla f(x_0)\cdot h+\frac12\,h^T H_f(x_0)\,h+\text{(higher order terms)}.
	\]
	The term $\nabla f(x_0)\cdot h$ is the \emph{first-order slope}.
	Thus if $\nabla f(x_0)=0$, then
	\[
	f(x_0+h)=f(x_0)+\frac12\,h^T H_f(x_0)\,h+\cdots,
	\]
	so $f$ does not change to first order when we move away from $x_0$; the first nontrivial change begins at second order
	(or at even higher order if the critical point is degenerate).

	\medskip
	\noindent\textbf{(2) The case $k=2$: a surface in $\mathbb{R}^3$.}
	If $f(x,y)$ is a function of two variables, its graph is the surface
	\[
	S=\{(x,y,z)\in\mathbb{R}^3:\ z=f(x,y)\}.
	\]
	At a point $(x_0,y_0)$ the tangent plane to $S$ is
	\[
	z=f(x_0,y_0)+f_x(x_0,y_0)(x-x_0)+f_y(x_0,y_0)(y-y_0).
	\]
	Hence $(x_0,y_0)$ is critical iff $f_x(x_0,y_0)=f_y(x_0,y_0)=0$, and then the tangent plane simplifies to
	\[
	z=f(x_0,y_0),
	\]
	which is a \emph{horizontal plane}. This is the precise meaning of the phrase
	``the tangent plane is horizontal at a critical point.''

	\medskip
	\noindent\textbf{(3) What you see on a contour plot.}
	A contour plot draws level sets $f(x,y)=c$. If $\nabla f(x_0,y_0)\neq 0$, then nearby level curves are smooth and
	their normal direction is given by the gradient. If $\nabla f(x_0,y_0)=0$, the point can become a special point of
	the level sets: for a non-degenerate minimum/maximum one typically sees small closed curves around $(x_0,y_0)$, for a
	saddle one sees level curves crossing, and for degenerate critical points one can see multi-armed or ``pinched''
	patterns (e.g.\ the monkey saddle).

	\medskip
	\noindent\textbf{Figures.}

	\noindent\textbf{1D: critical points and the second derivative test.}
	Let $f:\mathbb{R}\to\mathbb{R}$ be $C^2$ and let $x_0\in\mathbb{R}$.
	A \emph{critical point} is a point where
	\[
	f'(x_0)=0.
	\]
	Taylor's theorem (with remainder) gives, for $x$ near $x_0$,
	\[
	f(x)=f(x_0)+f'(x_0)(x-x_0)+\frac12 f''(x_0)(x-x_0)^2 + o\!\left((x-x_0)^2\right).
	\]
	At a critical point $f'(x_0)=0$, so
	\[
	f(x)=f(x_0)+\frac12 f''(x_0)(x-x_0)^2 + o\!\left((x-x_0)^2\right).
	\]
	Hence:

		\par\noindent\textbullet\quad If $f''(x_0)>0$, then $x_0$ is a \emph{strict local minimum}.
		\par\noindent\textbullet\quad If $f''(x_0)<0$, then $x_0$ is a \emph{strict local maximum}.
		\par\noindent\textbullet\quad If $f''(x_0)=0$, the second derivative test is \emph{inconclusive}; one must use higher-order terms or other methods.

	\medskip
	\noindent\textbf{1D: why $x^2$ and $x^4$ differ near $0$.}
	Consider $f(x)=x^2$ and $g(x)=x^4$ at $0$:
	\[
	f(0)=0,\quad f'(0)=0,\quad f''(0)=2\neq 0,
	\]
	so the \emph{first nonzero Taylor term} is quadratic:
	\[
	f(x)=x^2.
	\]
	For $g(x)=x^4$,
	\[
	g(0)=0,\quad g'(0)=0,\quad g''(0)=0,\quad g^{(3)}(0)=0,\quad g^{(4)}(0)=24\neq 0,
	\]
	so the first nonzero Taylor term appears at order $4$:
	\[
	g(x)=\frac{1}{4!}g^{(4)}(0)x^4=x^4.
	\]
	This explains the geometric behavior: $x^4$ is \emph{much flatter near $0$} (since for small $|x|$ we have $|x|^4\ll |x|^2$), even though it grows faster for large $|x|$.

	\bigskip
	\noindent\textbf{2D: extended Taylor expansion and the Hessian.}
	Let $f:\mathbb{R}^2\to\mathbb{R}$ be $C^2$ and let $p=(a,b)$.
	Write $h=(h_1,h_2)=(x-a,y-b)$. Then Taylor's theorem in $\mathbb{R}^2$ gives
	\[
	f(p+h)
	=
	f(p)
	+
	\nabla f(p)\cdot h
	+
	\frac12\, h^\top H_f(p)\, h
	+
	o\!\left(\|h\|^2\right),
	\]
	where
	\[
	\nabla f(p)=
	\begin{pmatrix}
		f_x(p)\\[2pt]
		f_y(p)
	\end{pmatrix},
	\qquad
	H_f(p)=
	\begin{pmatrix}
		f_{xx}(p) & f_{xy}(p)\\[2pt]
		f_{yx}(p) & f_{yy}(p)
	\end{pmatrix},
	\]
	and
	\[
	h^\top H_f(p)\,h
	=
	f_{xx}(p)\,h_1^2 + 2f_{xy}(p)\,h_1h_2 + f_{yy}(p)\,h_2^2.
	\]

	\medskip
	\noindent\textbf{Critical points and leading-order shape.}
	A point $p$ is a \emph{critical point} iff
	\[
	\nabla f(p)=0.
	\]
	At a critical point the linear term vanishes, and the first term controlling the local geometry is the quadratic form
	\[
	Q(h)=\frac12\, h^\top H_f(p)\,h.
	\]
	This is the multivariable analogue of using $f''(x_0)$ in 1D.

	\medskip
	\noindent\textbf{Nondegenerate vs degenerate critical points.}

		\par\noindent\textbullet\quad $p$ is \emph{nondegenerate} if the Hessian is invertible:
		\[
		\det H_f(p)\neq 0.
		\]
		\par\noindent\textbullet\quad $p$ is \emph{degenerate} if the Hessian is singular:
		\[
		\det H_f(p)=0.
		\]
		In this case the quadratic approximation does \emph{not} fully determine the local shape; higher-order terms are needed.

	\medskip
	\noindent\textbf{2D classification via the determinant test (nondegenerate case).}
	At a critical point $p$, set
	\[
	D=\det H_f(p)=f_{xx}(p)f_{yy}(p)-f_{xy}(p)^2.
	\]
	Then:

		\par\noindent\textbullet\quad If $D>0$ and $f_{xx}(p)>0$, then $p$ is a \emph{strict local minimum}.
		\par\noindent\textbullet\quad If $D>0$ and $f_{xx}(p)<0$, then $p$ is a \emph{strict local maximum}.
		\par\noindent\textbullet\quad If $D<0$, then $H_f(p)$ is \emph{indefinite} and $p$ is a \emph{saddle point}.
		\par\noindent\textbullet\quad If $D=0$, the test is \emph{inconclusive} (the point is degenerate).

	\bigskip
	\noindent\textbf{Saddles vs degeneracy: key examples.}

	\medskip
	\noindent\emph{(1) Nondegenerate saddle.}
	Let
	\[
	f(x,y)=x^2-y^2.
	\]
	Then $(0,0)$ is critical and
	\[
	H_f(0,0)=
	\begin{pmatrix}
		2 & 0\\
		0 & -2
	\end{pmatrix},
	\qquad
	\det H_f(0,0)=-4<0.
	\]
	Hence $(0,0)$ is a \emph{nondegenerate saddle}. Geometrically, the surface bends upward in the $x$-direction and downward in the $y$-direction.

	\medskip
	\noindent\emph{(2) Degenerate saddle.}
	Let
	\[
	g(x,y)=x^4-y^4.
	\]
	Then $(0,0)$ is critical, and all second partial derivatives vanish at $(0,0)$, so
	\[
	H_g(0,0)=0,
	\qquad
	\det H_g(0,0)=0,
	\]
	i.e.\ $(0,0)$ is \emph{degenerate}. Nevertheless it is still a \emph{saddle} because
	\[
	g(x,0)=x^4\ge 0,
	\qquad
	g(0,y)=-y^4\le 0,
	\]
	so arbitrarily close to $(0,0)$ one finds points where $g>0$ and points where $g<0$.

	\medskip
	\noindent\emph{(3) Degenerate minimum (flat bottom).}
	Let
	\[
	h(x,y)=x^4+y^4.
	\]
	Then $(0,0)$ is critical and degenerate ($H_h(0,0)=0$), but $h(x,y)\ge 0$ with equality only at $(0,0)$, so $(0,0)$ is a \emph{(degenerate) strict local minimum}.

	\medskip
	\noindent\textbf{Moral.}

		\par\noindent\textbullet\quad \emph{Saddle} means: the function takes values both above and below $f(p)$ arbitrarily close to $p$.
		\par\noindent\textbullet\quad \emph{Degenerate} means: the quadratic (Hessian) term is insufficient ($\det H_f(p)=0$).
		\par\noindent\textbullet\quad Saddles can be nondegenerate (e.g.\ $x^2-y^2$) or degenerate (e.g.\ $x^4-y^4$).

	\subsection*{Why $d(\det)_A(H)=\mathrm{tr}(\mathrm{adj}(A)^T H)$}

	\noindent
	Let $\det:M(n)\to\mathbb{R}$ be the determinant function. Fix $A\in M(n)$ and a direction $H\in M(n)$.
	Define a one-variable function
	\[
	g(t)=\det(A+tH).
	\]
	By definition of the directional derivative,
	\[
	d(\det)_A(H)=g'(0).
	\]

	\medskip
	\noindent\textbf{Step 1: Multilinearity in columns.}
	Write the columns of $A$ as $a_1,\dots,a_n$ and the columns of $H$ as $h_1,\dots,h_n$.
	Then
	\[
	A+tH=[\,a_1+th_1,\ a_2+th_2,\ \dots,\ a_n+th_n\,].
	\]
	Since $\det$ is multilinear in the columns,
	\[
	\det(A+tH)=\det(a_1+th_1,\dots,a_n+th_n)
	\]
	\[
	=\det(a_1,\dots,a_n)
	+t\sum_{j=1}^n \det(a_1,\dots,h_j,\dots,a_n)
	+\text{(terms of order }t^2\text{ and higher)}.
	\]
	Therefore the coefficient of $t$ equals $g'(0)$, and hence
	\[
	d(\det)_A(H)=g'(0)=\sum_{j=1}^n \det(a_1,\dots,h_j,\dots,a_n).
	\tag{$\ast$}
	\]

	\medskip
	\noindent\textbf{Step 2: Laplace expansion along the replaced column.}
	Fix an index $j$. Let $B=(a_1,\dots,h_j,\dots,a_n)$ be the matrix obtained by replacing the $j$-th column of $A$ by $h_j$.
	Expand $\det(B)$ along the $j$-th column:
	\[
	\det(B)=\sum_{i=1}^n B_{ij}\,C_{ij}(B),
	\]
	where $C_{ij}(B)=(-1)^{i+j}\det(B_{\widehat{i},\widehat{j}})$ is the $(i,j)$-cofactor.
	Deleting row $i$ and column $j$ removes the replaced column, hence
	\[
	B_{\widehat{i},\widehat{j}}=A_{\widehat{i},\widehat{j}}
	\qquad\Longrightarrow\qquad
	C_{ij}(B)=(-1)^{i+j}\det(A_{\widehat{i},\widehat{j}})=C_{ij}(A).
	\]
	Moreover $B_{ij}=(h_j)_i=H_{ij}$. Therefore
	\[
	\det(a_1,\dots,h_j,\dots,a_n)=\det(B)=\sum_{i=1}^n H_{ij}\,C_{ij}(A).
	\]

	\medskip
	\noindent\textbf{Step 3: Sum over all $j$ and identify the adjugate.}
	Substituting the previous formula into $(\ast)$ gives
	\[
	d(\det)_A(H)=\sum_{j=1}^n\sum_{i=1}^n H_{ij}\,C_{ij}(A)
	=\sum_{i,j} C_{ij}(A)\,H_{ij}.
	\]
	Recall that the adjugate matrix satisfies
	\[
	\mathrm{adj}(A)_{ji}=C_{ij}(A)
	\quad\text{(adjugate is the transpose of the cofactor matrix).}
	\]
	Hence
	\[
	\sum_{i,j} C_{ij}(A)\,H_{ij}
	=\sum_{i,j} \mathrm{adj}(A)_{ji}\,H_{ij}.
	\]
	Finally, using the identity
	\[
	\mathrm{tr}(X^T Y)=\sum_{i,j} X_{ij}Y_{ij},
	\]
	we obtain
	\[
	\mathrm{tr}\big(\mathrm{adj}(A)^T H\big)
	=\sum_{i,j} \big(\mathrm{adj}(A)^T\big)_{ij}H_{ij}
	=\sum_{i,j} \mathrm{adj}(A)_{ji}H_{ij}.
	\]
	Therefore,
	\[
	d(\det)_A(H)=\mathrm{tr}\big(\mathrm{adj}(A)^T H\big).
	\]

	\noindent\textbf{Multilinearity (in columns).}
	Write an $n\times n$ matrix by its columns
	\[
	A=[c_1,\dots,c_n],\qquad c_j\in\mathbb{R}^n.
	\]
	A map $F(c_1,\dots,c_n)$ is called \emph{multilinear} if it is linear in each argument separately, i.e.\ for every
	fixed $j$ and all vectors $u,v\in\mathbb{R}^n$ and scalars $\lambda\in\mathbb{R}$,
	\begin{center}
	\resizebox{\linewidth}{!}{$\displaystyle
	F(c_1,\dots,c_{j-1},u+v,c_{j+1},\dots,c_n)
	=
	F(c_1,\dots,c_{j-1},u,c_{j+1},\dots,c_n)
	+
	F(c_1,\dots,c_{j-1},v,c_{j+1},\dots,c_n),
	$}
	\end{center}
	and
	\[
	F(c_1,\dots,c_{j-1},\lambda u,c_{j+1},\dots,c_n)
	=
	\lambda\,F(c_1,\dots,c_{j-1},u,c_{j+1},\dots,c_n).
	\]
	For the determinant, this property follows from the Leibniz formula:
	\[
	\det(A)=\sum_{\sigma\in S_n}\mathrm{sgn}(\sigma)\prod_{i=1}^n a_{i,\sigma(i)}.
	\]
	Fix a column index $j$. In each product $\prod_{i=1}^n a_{i,\sigma(i)}$, the permutation $\sigma$ hits each column
	index exactly once, so \emph{exactly one factor} in the product comes from the $j$-th column. Therefore each summand
	depends linearly on the entries of the $j$-th column, and so does the sum. Since $j$ was arbitrary, $\det$ is linear
	in each column separately, i.e.\ $\det$ is multilinear.

	\medskip
	\noindent\textbf{Where the powers $t^2,t^3,\dots$ come from.}
	Let $A,H\in M(n)$ and write their columns as
	\[
	A=[a_1,\dots,a_n],\qquad H=[h_1,\dots,h_n].
	\]
	Then
	\[
	A+tH=[a_1+th_1,\dots,a_n+th_n].
	\]
	Using multilinearity, we can expand $\det(A+tH)$ by distributing over each column (analogous to expanding a product).
	From each column $a_j+th_j$ we choose either $a_j$ or $th_j$. This yields a sum of $2^n$ determinants:
	\[
	\det(A+tH)
	=
	\sum_{S\subset\{1,\dots,n\}}
	\det\big(b_1(S),\dots,b_n(S)\big),
	\]
	where
	\[
	b_j(S)=
	\begin{cases}
		th_j, & j\in S,\\
		a_j, & j\notin S.
	\end{cases}
	\]
	If $|S|=m$, then exactly $m$ columns contribute a factor $t$, and by linearity we can factor these out:
	\[
	\det\big(b_1(S),\dots,b_n(S)\big)
	=
	t^m\,\det\big(c_1(S),\dots,c_n(S)\big),
	\]
	where $c_j(S)=h_j$ if $j\in S$ and $c_j(S)=a_j$ otherwise.
	Hence the expansion groups naturally by the size $m$ of $S$:
	\begin{center}
	\resizebox{\linewidth}{!}{$\displaystyle
	\det(A+tH)
	=
	\det(A)
	+
	t\sum_{|S|=1}\det(\text{replace one column of }A\text{ by the corresponding column of }H)
	+
	t^2(\cdots)
	+\cdots+t^n\det(H).
	$}
	\end{center}
	In particular, the terms with $t^2,t^3,\dots$ come from choosing $th_j$ in two, three, $\dots$ columns respectively.

	\medskip
	\noindent\textbf{Why only the linear term matters for the derivative at $t=0$.}
	If we write
	\[
	\det(A+tH)=\det(A)+tL+t^2Q(t),
	\]
	then
	\[
	\frac{d}{dt}\Big|_{t=0}\det(A+tH)=L,
	\]
	because $\frac{d}{dt}\big(t^2Q(t)\big)\big|_{t=0}=0$.
	Therefore $d(\det)_A(H)$ is exactly the coefficient of $t$ in the multilinear expansion above.

	\noindent\textbf{$3\times 3$ case (fully expanded).}
	Let
	\[
	A=\begin{pmatrix}
		a&b&c\\
		d&e&f\\
		g&h&i
	\end{pmatrix},\qquad
	H=\begin{pmatrix}
		p&q&r\\
		s&t&u\\
		v&w&x
	\end{pmatrix}.
	\]
	Writing $A=[a_1,a_2,a_3]$ and $H=[h_1,h_2,h_3]$ by columns, multilinearity gives
	\[
	d(\det)_A(H)=\frac{d}{dt}\Big|_{0}\det(A+tH)
	=\det(h_1,a_2,a_3)+\det(a_1,h_2,a_3)+\det(a_1,a_2,h_3).
	\]
	Expanding each term along its replaced column yields
	\[
	d(\det)_A(H)=\sum_{i=1}^3 C_{i1}(A)H_{i1}+\sum_{i=1}^3 C_{i2}(A)H_{i2}+\sum_{i=1}^3 C_{i3}(A)H_{i3}
	=\sum_{i,j=1}^3 C_{ij}(A)H_{ij}.
	\]
	The cofactors of $A$ are
	\[
	\begin{aligned}
		C_{11}&=ei-fh, & C_{12}&=fg-di, & C_{13}&=dh-eg,\\
		C_{21}&=ch-bi, & C_{22}&=ai-cg, & C_{23}&=bg-ah,\\
		C_{31}&=bf-ce, & C_{32}&=cd-af, & C_{33}&=ae-bd.
	\end{aligned}
	\]
	Hence
	\begin{center}
	\resizebox{\linewidth}{!}{$\displaystyle
	\boxed{
		d(\det)_A(H)=
		(ei-fh)p+(fg-di)q+(dh-eg)r
		+(ch-bi)s+(ai-cg)t+(bg-ah)u
		+(bf-ce)v+(cd-af)w+(ae-bd)x.}
	$}
	\end{center}

	\medskip
	\noindent\textbf{$4\times 4$ case (column-replacement formula).}
	Write $A=[a_1,a_2,a_3,a_4]$ and $H=[h_1,h_2,h_3,h_4]$. Then
	\begin{center}
	\resizebox{\linewidth}{!}{$\displaystyle
	\boxed{
		d(\det)_A(H)=
		\det(h_1,a_2,a_3,a_4)+
		\det(a_1,h_2,a_3,a_4)+
		\det(a_1,a_2,h_3,a_4)+
		\det(a_1,a_2,a_3,h_4).}
	$}
	\end{center}
	Equivalently,
	\[
	d(\det)_A(H)=\sum_{i,j=1}^4 C_{ij}(A)H_{ij},
	\qquad
	C_{ij}(A)=(-1)^{i+j}\det(A_{\widehat{i},\widehat{j}}).
	\]

	\medskip
	\noindent\textbf{Example (diagonal $A$).}
	If $A=\mathrm{diag}(d_1,d_2,d_3,d_4)$ and $H=(h_{ij})$, then
	\[
	\boxed{
		d(\det)_A(H)=d_2d_3d_4\,h_{11}+d_1d_3d_4\,h_{22}+d_1d_2d_4\,h_{33}+d_1d_2d_3\,h_{44}.}
	\]

	\subsection*{Visualizing $\det:M(2)\cong\mathbb{R}^4\to\mathbb{R}$ and its derivative via slices}

	\noindent
	A $2\times 2$ matrix can be parametrized by $(a,b,c,d)\in\mathbb{R}^4$:
	\[
	A=\begin{pmatrix}a&b\\ c&d\end{pmatrix},
	\qquad
	\det(A)=ad-bc.
	\]
	Since the domain is $4$-dimensional, we cannot plot $\det$ directly. Instead, we fix two coordinates and visualize
	$2$-dimensional slices. For example, if $b,c$ are fixed then $\det$ becomes a function of $(a,d)$:
	\[
	\det(A)=ad-(bc),
	\]
	so its level sets in the $(a,d)$-plane are hyperbolas. Similarly, if $a,d$ are fixed then
	\[
	\det(A)=(ad)-bc,
	\]
	and level sets in the $(b,c)$-plane are again hyperbolas. To visualize the \emph{derivative} on a slice, we plot the
	gradient vector field of the sliced function (arrows give the direction of steepest increase \emph{within the slice}).

	\medskip
	\noindent\textbf{Figures (determinant slices).}

	\medskip
	\noindent\textbf{Figures (derivative on slices).}
	\noindent
	On the $(a,d)$-slice (with fixed $b,c$), the sliced function is $F(a,d)=ad-(bc)$, hence
	\[
	\nabla_{(a,d)}F(a,d)=(\partial_a F,\partial_d F)=(d,a).
	\]
	On the $(b,c)$-slice (with fixed $a,d$), the sliced function is $G(b,c)=(ad)-bc$, hence
	\[
	\nabla_{(b,c)}G(b,c)=(\partial_b G,\partial_c G)=(-c,-b).
	\]

	\noindent
	Identify $M(2)\cong\mathbb{R}^4$ via coordinates $(a,b,c,d)$:
	\[
	F(a,b,c,d)=\det\!\begin{pmatrix}a&b\\ c&d\end{pmatrix}=ad-bc.
	\]
	Since $F$ has a $4$-dimensional domain, we visualize it by fixing two variables and taking a $2$-dimensional slice.
	Fix $b,c\in\mathbb{R}$. Then the restricted (slice) function is
	\[
	f_{b,c}(a,d):=F(a,b,c,d)=ad-bc,
	\qquad (a,d)\in\mathbb{R}^2.
	\]

	\medskip
	\noindent\textbf{Derivative (differential) on the slice.}
	The derivative of $f_{b,c}$ at a point $(a,d)$ is the linear map
	\[
	df_{b,c}\big|_{(a,d)}:\mathbb{R}^2\to\mathbb{R},
	\]
	defined by
	\[
	df_{b,c}\big|_{(a,d)}(\Delta a,\Delta d)
	=\frac{\partial f_{b,c}}{\partial a}(a,d)\,\Delta a
	+\frac{\partial f_{b,c}}{\partial d}(a,d)\,\Delta d.
	\]
	For $f_{b,c}(a,d)=ad-bc$ we have
	\[
	\frac{\partial f_{b,c}}{\partial a}(a,d)=d,
	\qquad
	\frac{\partial f_{b,c}}{\partial d}(a,d)=a,
	\]
	hence
	\[
	\boxed{
		df_{b,c}\big|_{(a,d)}(\Delta a,\Delta d)=d\,\Delta a+a\,\Delta d.}
	\]

	\medskip
	\noindent\textbf{Gradient as a representation of the derivative.}
	In $\mathbb{R}^2$ with the standard dot product, the differential can be written as
	\[
	df_{b,c}\big|_{(a,d)}(\Delta a,\Delta d)=\nabla f_{b,c}(a,d)\cdot(\Delta a,\Delta d),
	\]
	so the gradient vector is
	\[
	\boxed{\nabla f_{b,c}(a,d)=\big(d,a\big).}
	\]
	Therefore, the arrow field in the figure is a visualization of the derivative:
	each arrow at $(a,d)$ points in the direction of fastest increase of $f_{b,c}$ within the $(a,d)$-slice.

	\medskip
	\noindent\textbf{Full $4$-dimensional derivative (for completeness).}
	For $F(a,b,c,d)=ad-bc$,
	\[
	\nabla F(a,b,c,d)=\big(d,-c,-b,a\big),
	\]
	and the differential at $(a,b,c,d)$ is the linear map
	\[
	\boxed{
		dF\big|_{(a,b,c,d)}(\Delta a,\Delta b,\Delta c,\Delta d)
		=d\,\Delta a-c\,\Delta b-b\,\Delta c+a\,\Delta d.}
	\]
	The $(a,d)$-slice figure corresponds to restricting this differential to variations with $\Delta b=\Delta c=0$.

	\subsection*{Why $df_{(a,d)}(\Delta a,\Delta d)=\nabla f(a,d)\cdot(\Delta a,\Delta d)$?}

	Let $f:\mathbb{R}^2\to\mathbb{R}$ be differentiable and let $p=(a,d)$.
	The \emph{differential} of $f$ at $p$ is the linear map (a linear functional)
	\[
	df_p:\mathbb{R}^2\to\mathbb{R}
	\]
	characterized by the first--order Taylor approximation
	\[
	f(p+h)=f(p)+df_p(h)+o(\|h\|)\qquad (h\to 0).
	\]
	Writing $h=(\Delta a,\Delta d)$ and using partial derivatives, one obtains the explicit formula
	\[
	df_{(a,d)}(\Delta a,\Delta d)=
	\frac{\partial f}{\partial x}(a,d)\,\Delta a
	+\frac{\partial f}{\partial y}(a,d)\,\Delta d.
	\]

	Now introduce the \emph{gradient} (with respect to the standard dot product on $\mathbb{R}^2$)
	\[
	\nabla f(a,d)=\left(\frac{\partial f}{\partial x}(a,d),\ \frac{\partial f}{\partial y}(a,d)\right).
	\]
	Then the same expression can be rewritten as a dot product:
	\[
	\frac{\partial f}{\partial x}(a,d)\,\Delta a
	+\frac{\partial f}{\partial y}(a,d)\,\Delta d
	=
	\left(\frac{\partial f}{\partial x}(a,d),\frac{\partial f}{\partial y}(a,d)\right)\cdot(\Delta a,\Delta d)
	=
	\nabla f(a,d)\cdot(\Delta a,\Delta d).
	\]
	Hence
	\[
	df_{(a,d)}(\Delta a,\Delta d)=\nabla f(a,d)\cdot(\Delta a,\Delta d).
	\]

	\medskip
	\noindent\textbf{Conceptual reason.}
	The differential $df_p$ is a covector (linear functional). In Euclidean space with the standard inner product
	$\langle\cdot,\cdot\rangle$, every covector $L$ can be represented uniquely as
	\[
	L(h)=\langle v,h\rangle
	\quad\text{for a unique vector }v.
	\]
	The gradient is precisely this representing vector: it is defined by
	\[
	df_p(h)=\langle \nabla f(p),h\rangle\qquad\text{for all }h\in\mathbb{R}^2.
	\]
	For the standard inner product, $\langle u,v\rangle=u\cdot v$, giving the stated formula.
\end{problem}

\begin{problem}[CP-I-0030]
\label{prob:cp-i-0030}
\textbf{Cauchy--Schwarz as a positivity principle}

\hspace*{1.5em}Cauchy--Schwarz:
\[
|\langle u,v\rangle|\le \|u\|\,\|v\|,
\]
with equality iff $u$ and $v$ are linearly dependent.
A standard proof uses
\[
\|u-tv\|^2\ge 0\ \ \forall t\in\mathbb R
\]
and discriminant $\le 0$.
Many inequalities reduce to ``sum of squares'', e.g.
\[
x^2+y^2+z^2-(xy+yz+zx)=\tfrac12\big((x-y)^2+(y-z)^2+(z-x)^2\big)\ge 0.
\]

For unit $n$,
\[
T(x)=x-(x\cdot n)n
\]
is orthogonal projection onto $n^\perp$:
\[
\ker T=\mathrm{span}\{n\},\qquad \mathrm{Im}\,T=n^\perp,\qquad T^2=T.
\]
The cross-product map
\[
Q(x)=n\times x
\]
is represented by the skew-symmetric matrix
\[
[n]_\times=\begin{pmatrix}
	0&-n_3&n_2\\
	n_3&0&-n_1\\
	-n_2&n_1&0
\end{pmatrix},
\qquad Q(x)=[n]_\times x.
\]
Using $n\times(n\times x)=n(n\cdot x)-x$, one gets
\[
Q^2=-T,\qquad Q^3=-Q,\qquad Q^4=T.
\]
For any $A:\mathbb R^3\to\mathbb R^3$, $\mathrm{Im}\,A$ is the column space and $\ker A$ solves $Ax=0$; the rank determines whether the image is a line (rank $1$), plane (rank $2$), or all $\mathbb R^3$ (rank $3$).
\end{problem}

\begin{problem}[CP-I-0032]
\label{prob:cp-i-0032}
\textbf{Cofactors, inverse, and solving a linear system}

\textbf{Exact statement.}
\hspace*{1.5em}Calculate the cofactors \(\Delta_{ij}\) for
\[
A=\begin{pmatrix}
	1&1&1\\
	1&2&3\\
	3&-2&2
\end{pmatrix}.
\]
Convert \(A_{ij}\Delta_{ik}=\delta_{jk}\det A\) into matrix notation and check it.
Find \(A^{-1}\).
Use it to solve
\[
\begin{cases}
	x+y+z=1\\
	x+2y+3z=-5\\
	3x-2y+2z=4.
\end{cases}
\]

\end{problem}

\begin{solution}
The cofactor matrix \(C=(\Delta_{ij})\) is
\[
C=\begin{pmatrix}
	10&7&-8\\
	-4&-1&5\\
	1&-2&1
\end{pmatrix},
\qquad \det(A)=9.
\]
The suffix identity \(A_{ij}\Delta_{ik}=\delta_{jk}\det(A)\) (sum over \(i\)) is exactly the matrix identity
\[
A^T\,C=\det(A)\,I.
\]
Equivalently, with \(\operatorname{adj}(A)=C^T\),
\[
A\,\operatorname{adj}(A)=\det(A)\,I.
\]
Thus
\[
A^{-1}=\frac{1}{\det(A)}\operatorname{adj}(A)=\frac{1}{9}C^T
=\frac{1}{9}\begin{pmatrix}
	10&-4&1\\
	7&-1&-2\\
	-8&5&1
\end{pmatrix}.
\]
To solve \(A(x,y,z)^T=(1,-5,4)^T\),
\[
\begin{pmatrix}x\\y\\z\end{pmatrix}
=A^{-1}\begin{pmatrix}1\\-5\\4\end{pmatrix}
=\begin{pmatrix}\frac{34}{9}\\[2pt]\frac{4}{9}\\[2pt]-\frac{29}{9}\end{pmatrix}.
\]
A direct substitution verifies the system.
\end{solution}

\begin{problem}[CP-I-0034]
\label{prob:cp-i-0034}
\textbf{Geometric Multiplicity and Jordan-Block Count}

\textbf{Source-recovered canonical problem from the Jordan-form subsection.}

\hspace*{1.5em}Let \(A\in M_n(\mathbb C)\), and fix an eigenvalue \(\lambda\).
Suppose the Jordan blocks belonging to \(\lambda\) have sizes
\[
m_1,\dots,m_r.
\]
Prove that the algebraic multiplicity of \(\lambda\) is
\[
m_1+\cdots+m_r,
\]
while its geometric multiplicity is
\[
\dim\ker(A-\lambda I)=r.
\]
Deduce that \(A\) is diagonalizable on the generalized \(\lambda\)-eigenspace if and only if
\(r=m_1+\cdots+m_r\).
\end{problem}

\begin{solution}
A Jordan block \(J_{m_j}(\lambda)\) contributes the factor
\[
(t-\lambda)^{m_j}
\]
to the characteristic polynomial. Hence the total power of \(t-\lambda\) is
\[
m_1+\cdots+m_r,
\]
which is the algebraic multiplicity of \(\lambda\).

For a single block,
\[
J_{m_j}(\lambda)-\lambda I
\]
is the nilpotent Jordan block with one-dimensional kernel. Therefore each Jordan block contributes
exactly one independent eigenvector. Since the Jordan form is a direct sum, the eigenspace
dimension is
\[
\dim\ker(A-\lambda I)=r.
\]

On the generalized \(\lambda\)-eigenspace the operator is diagonalizable exactly when every
Jordan block has size \(1\). This is equivalent to
\[
m_1=\cdots=m_r=1,
\]
and hence to
\[
r=m_1+\cdots+m_r.
\]
\end{solution}

\begin{problem}[CP-I-0037]
\label{prob:cp-i-0037}
\textbf{Matrices of a linear map in nonstandard bases (and change of basis)}

\textbf{Exact statement.}
\hspace*{1.5em}Define the $m\times n$ matrix $A$ that represents a linear map $T:\mathbb R^n\to\mathbb R^m$ with respect to general bases
$\{e_1,\dots,e_n\}$ and $\{f_1,\dots,f_m\}$.

\textbf{(a)} Taking $n=2,m=3$, let $T$ be defined by
\[
T\binom10=\begin{pmatrix}2\\1\\5\end{pmatrix},\qquad
T\binom01=\begin{pmatrix}7\\0\\3\end{pmatrix}.
\]
Find the matrix $A$ with respect to the bases
\[
e_1=\binom{1}{-1},\quad e_2=\binom{-3}{2};\qquad
f_1=\begin{pmatrix}1\\0\\0\end{pmatrix},\
f_2=\begin{pmatrix}0\\1\\0\end{pmatrix},\
f_3=\begin{pmatrix}0\\0\\1\end{pmatrix}.
\]

\textbf{(b)} Taking $n=m=3$, let $T$ be reflection in the plane $x_1\sin\theta=x_2\cos\theta$.
Find the matrix $A$ with respect to a convenient choice of bases (to be specified) such that $e_i=f_i$ $(i=1,2,3)$.

\textbf{(c)} Taking $n=m=2$, let $T$ be the shear with parameter $\lambda$ defined by
\[
T\binom10=\binom10,\qquad T\binom01=\binom{\lambda}{1}.
\]
Find the matrix $A$ when $e_1=f_1$ and $e_2=f_2$ are the standard basis vectors for $\mathbb R^2$; find also the matrix $A'$
with respect to the new basis $e_1'=f_1'=-e_2$ and $e_2'=f_2'=e_1$.
Show that $A'=R^{-1}AR$ for a rotation matrix $R$ and comment on this result.

\end{problem}

\begin{solution}
\textbf{General definition.} The representing matrix $A=(a_{ij})$ is defined by
\[
T(e_j)=\sum_{i=1}^m a_{ij} f_i\qquad (j=1,\dots,n).
\]
Equivalently, for all $v\in\mathbb R^n$,
\[
[T(v)]_F = A\,[v]_E.
\]

\textbf{(a)} In the standard basis of $\mathbb R^2$, the matrix of $T$ is
\[
M=\begin{pmatrix}2&7\\1&0\\5&3\end{pmatrix}.
\]
Since $(f_1,f_2,f_3)$ is the standard basis of $\mathbb R^3$, coordinates in $F$ are just components. Compute
\[
T(e_1)=T\binom{1}{-1}
=\begin{pmatrix}2\\1\\5\end{pmatrix}-\begin{pmatrix}7\\0\\3\end{pmatrix}
=\begin{pmatrix}-5\\1\\2\end{pmatrix},
\]
\[
T(e_2)=T\binom{-3}{2}
=-3\begin{pmatrix}2\\1\\5\end{pmatrix}+2\begin{pmatrix}7\\0\\3\end{pmatrix}
=\begin{pmatrix}8\\-3\\-9\end{pmatrix}.
\]
Hence the representing matrix (columns are $[T(e_1)]_F$ and $[T(e_2)]_F$) is
\[
\boxed{
	A=\begin{pmatrix}
		-5&8\\
		1&-3\\
		2&-9
	\end{pmatrix}.}
\]

\textbf{(b)} The reflecting plane is $x_1\sin\theta-x_2\cos\theta=0$, so a unit normal is
\[
e_1=n=(\sin\theta,-\cos\theta,0).
\]
Choose a basis adapted to the plane:
\[
e_2=(\cos\theta,\sin\theta,0)\in \text{plane},\qquad e_3=(0,0,1)\in \text{plane}.
\]
Reflection fixes vectors in the plane and flips the normal, hence
\[
T(e_1)=-e_1,\qquad T(e_2)=e_2,\qquad T(e_3)=e_3.
\]
With respect to the basis $(e_1,e_2,e_3)$ (and taking $f_i=e_i$), the matrix is
\[
\boxed{A=\operatorname{diag}(-1,1,1).}
\]
(Optionally, in the standard basis one may write $T(x)=x-2(x\cdot n)n$ to obtain
$\begin{psmallmatrix}
	\cos 2\theta & \sin 2\theta & 0\\
	\sin 2\theta & -\cos 2\theta & 0\\
	0&0&1
\end{psmallmatrix}$.)

\textbf{(c)} In the standard basis, the columns are $T(e_1)$ and $T(e_2)$, so
\[
\boxed{
	A=\begin{pmatrix}
		1&\lambda\\
		0&1
	\end{pmatrix}.}
\]
For the new basis $e_1'=-e_2$ and $e_2'=e_1$, the change-of-basis matrix is
\[
R=[e_1'\ e_2']=\begin{pmatrix}0&1\\-1&0\end{pmatrix},
\qquad R^{-1}=R^T.
\]
Therefore the representing matrix in the primed basis is
\[
A'=R^{-1}AR,
\]
and a direct multiplication gives
\[
\boxed{
	A'=\begin{pmatrix}
		1&0\\
		-\lambda&1
	\end{pmatrix}.}
\]
This illustrates the general rule: matrices of the same linear map in different bases are similar, so they have the same invariants
(eigenvalues, trace, determinant).
\end{solution}

\begin{problem}[CP-I-0038]
\label{prob:cp-i-0038}
\textbf{Determinant of a “\(p\) on diagonal, \(1\) elsewhere” matrix}

\textbf{Exact statement.}
\hspace*{1.5em}Let \(D\) be the \(n\times n\) matrix (\(n>1\)) which has entries \(p\) on the main diagonal and entries \(1\) elsewhere.
Show that
\[
\det D=(p+n-1)(p-1)^{n-1}.
\]

\end{problem}

\begin{solution}
Write
\[
D=(p-1)I+J,
\]
where \(J\) is the all-ones matrix.
The eigenvalues of \(J\) are \(n\) (with eigenvector \(\mathbf{1}=(1,\dots,1)^T\)) and \(0\) with multiplicity \(n-1\).
Hence the eigenvalues of \(D\) are \(p+n-1\) (once) and \(p-1\) (multiplicity \(n-1\)).
Therefore,
\[
\det D=(p+n-1)(p-1)^{n-1}.
\]
\end{solution}

\begin{problem}[CP-I-0039]
\label{prob:cp-i-0039}
\textbf{Cayley--Hamilton and Recurrences for Matrix Powers}

\textbf{Canonical authored completion of the retained Cayley--Hamilton fragment.}

\hspace*{1.5em}Let \(A\in M_2(\mathbb F)\). Use the Cayley--Hamilton theorem to prove that every
power \(A^n\) with \(n\ge 2\) can be written in the form
\[
A^n=\alpha_n A+\beta_n I.
\]
Derive a recurrence for \(\alpha_n,\beta_n\) in terms of \(\operatorname{tr}A\) and \(\det A\).
Specialize to the repeated-root case
\[
\chi_A(t)=(t-\lambda)^2.
\]
\end{problem}

\begin{solution}
For a \(2\times2\) matrix,
\[
\chi_A(t)=t^2-(\operatorname{tr}A)t+\det A.
\]
Cayley--Hamilton gives
\[
A^2=(\operatorname{tr}A)A-(\det A)I.
\]
Suppose
\[
A^n=\alpha_nA+\beta_nI.
\]
Multiplying by \(A\) and reducing \(A^2\),
\[
A^{n+1}
=\alpha_nA^2+\beta_nA
=\bigl((\operatorname{tr}A)\alpha_n+\beta_n\bigr)A
-(\det A)\alpha_n I.
\]
Hence
\[
\alpha_{n+1}=(\operatorname{tr}A)\alpha_n+\beta_n,
\qquad
\beta_{n+1}=-(\det A)\alpha_n,
\]
with
\[
(\alpha_0,\beta_0)=(0,1),\qquad
(\alpha_1,\beta_1)=(1,0).
\]

If \(\chi_A(t)=(t-\lambda)^2\), write \(A=\lambda I+N\). Cayley--Hamilton implies
\[
N^2=0.
\]
Therefore
\[
A^n=(\lambda I+N)^n
=\lambda^nI+n\lambda^{n-1}N
=n\lambda^{n-1}A+(1-n)\lambda^nI
\]
for \(\lambda\neq0\); the nilpotent case \(\lambda=0\) follows directly from \(A^2=0\).
\end{solution}

\begin{problem}[CP-I-0040]
\label{prob:cp-i-0040}
\textbf{Eigenvalues of $\alpha^2$ and comparison of eigenspaces}

\textbf{Exact statement.}
\hspace*{1.5em}Let $\alpha\in\mathrm{End}(V)$ for a complex vector space $V$.
Show: if $\lambda$ is an eigenvalue of $\alpha$, then $\lambda^2$ is an eigenvalue of $\alpha^2$.
Show every eigenvalue of $\alpha^2$ arises this way.
Are $\ker(\alpha-\lambda I)$ and $\ker(\alpha^2-\lambda^2 I)$ necessarily the same?

\end{problem}

\begin{solution}
If $\alpha v=\lambda v$ with $v\neq 0$, then $\alpha^2v=\lambda^2v$, so $\lambda^2$ is an eigenvalue of $\alpha^2$.
Over $\mathbb C$, $\alpha$ is triangularisable, so in some basis $\alpha$ is upper triangular with diagonal entries the eigenvalues $\lambda_i$.
Then $\alpha^2$ is upper triangular with diagonal entries $\lambda_i^2$, hence the eigenvalues of $\alpha^2$ are exactly $\{\lambda_i^2\}$.

Moreover,
\[
\alpha^2-\lambda^2 I=(\alpha-\lambda I)(\alpha+\lambda I),
\]
so $\ker(\alpha-\lambda I)\subseteq \ker(\alpha^2-\lambda^2 I)$, but equality need not hold.
$\alpha=\mathrm{diag}(1,-1)$ gives $\alpha^2=I$ and
\[
\ker(\alpha-I)=\mathrm{span}\{(1,0)\}\neq \mathbb C^2=\ker(\alpha^2-I).
\]
\end{solution}

\begin{problem}[CP-I-0042]
\label{prob:cp-i-0042}
\textbf{Maximize a cyclic quadratic form under $\sum a_i=0$ and $\sum a_i^2=1$}

\textbf{Exact statement.}
\hspace*{1.5em}Let $a_1,a_2,\dots,a_n$ be real numbers such that
\[
a_1+\cdots+a_n=0,\qquad a_1^2+\cdots+a_n^2=1.
\]
What is the maximum value of
\[
a_1a_2+a_2a_3+\cdots+a_{n-1}a_n+a_na_1\ ?
\]
\end{problem}

\begin{solution}
Let $a=(a_1,\dots,a_n)^T$ and define the symmetric matrix
\[
Q=\frac12\begin{pmatrix}
	0&1&0&\cdots&0&1\\
	1&0&1&\cdots&0&0\\
	0&1&0&\ddots&0&0\\
	\vdots&\vdots&\ddots&\ddots&1&\vdots\\
	0&0&0&1&0&1\\
	1&0&0&\cdots&1&0
\end{pmatrix}.
\]
Then
\[
a_1a_2+\cdots+a_na_1=a^TQa.
\]
The constraints are $\|a\|^2=1$ and $a\perp \mathbf 1$ where $\mathbf 1=(1,\dots,1)$.

Let $\omega=e^{2\pi i/n}$ and $v_m=(1,\omega^m,\omega^{2m},\dots,\omega^{(n-1)m})^T$.
A direct computation using $(Qx)_j=\tfrac12(x_{j-1}+x_{j+1})$ (indices mod $n$) gives
\[
Qv_m=\frac12(\omega^m+\omega^{-m})v_m=\cos\!\left(\frac{2\pi m}{n}\right)v_m.
\]
Hence the eigenvalues of $Q$ are
\[
\lambda_m=\cos\!\left(\frac{2\pi m}{n}\right),\qquad m=0,1,\dots,n-1.
\]
The top eigenvalue is $\lambda_0=1$ with eigenvector $\mathbf 1$, but the constraint $\sum a_i=0$ excludes this direction.
Therefore the maximum of the Rayleigh quotient $a^TQa$ on $\mathbf 1^\perp$ is the next largest eigenvalue,
\[
\max_{\|a\|=1,\ a\perp \mathbf 1} a^TQa=\cos\!\left(\frac{2\pi}{n}\right).
\]
It is achieved by any real eigenvector in the $m=1$ eigenspace, e.g.
\[
a_j=\sqrt{\frac{2}{n}}\cos\!\left(\frac{2\pi j}{n}\right)\qquad (j=0,1,\dots,n-1),
\]
(up to phase shift), which satisfies $\sum a_j=0$ and $\sum a_j^2=1$.

\medskip
\noindent
\textbf{Answer.} $\boxed{\cos(2\pi/n)}$.
\end{solution}

\begin{problem}[CP-I-0043]
\label{prob:cp-i-0043}
\textbf{Jordan form of $J_n(\lambda)^k$}

\textbf{Exact statement.}
\hspace*{1.5em}Let $J_n(\lambda)$ be the $n\times n$ Jordan block with diagonal entry $\lambda\in\mathbb C$.
Find the Jordan normal form of $J_n(\lambda)^k$ for each $k\ge 1$.

\end{problem}

\begin{solution}
Write $J_n(\lambda)=\lambda I+N$ where $N=J_n(0)$ is nilpotent.

\emph{Case $\lambda\neq 0$.}
\[
J_n(\lambda)^k=(\lambda I+N)^k=\lambda^k\Bigl(I+\frac{1}{\lambda}N\Bigr)^k
=\lambda^k\Bigl(I+\frac{k}{\lambda}N+\cdots\Bigr).
\]
The nilpotent part is a polynomial in $N$ with nonzero linear coefficient, which preserves the single-block Jordan structure.
Hence
\[
J_n(\lambda)^k \sim J_n(\lambda^k)\qquad(\lambda\neq 0).
\]

\emph{Case $\lambda=0$.}
Then $J_n(0)=N$ and $N^k$ shifts by $k$ steps.
Write $n=qk+r$ with $0\le r<k$. Then $N^k$ decomposes into $k$ chains of lengths $q+1$ (for $r$ chains) and $q$ (for $k-r$ chains),
so
\[
J_n(0)^k \sim \underbrace{J_{q+1}(0)\oplus\cdots\oplus J_{q+1}(0)}_{r\text{ times}}
\;\oplus\;
\underbrace{J_q(0)\oplus\cdots\oplus J_q(0)}_{(k-r)\text{ times}},
\]
ignoring any $J_0(0)$ summands if $q=0$.
\end{solution}

\begin{problem}[CP-I-0044]
\label{prob:cp-i-0044}
\textbf{Brouwer Normalization for a Nonnegative Matrix}

\textbf{Source-recovered canonical problem from the Brouwer--Perron discussion.}

\hspace*{1.5em}Let \(A\ge0\) be an \(n\times n\) real matrix. Let
\[
\Delta=\left\{x\in\mathbb R^n:x_i\ge0,\ \sum_{i=1}^n x_i=1\right\}.
\]
Suppose
\[
s(x):=\sum_{i=1}^n(Ax)_i>0
\qquad (x\in\Delta),
\]
and define
\[
F(x)=\frac{Ax}{s(x)}.
\]
Show that \(F:\Delta\to\Delta\) is continuous. Use Brouwer's fixed-point theorem to prove that
\(A\) has a real eigenvalue \(\lambda\ge0\) with a nonnegative eigenvector.
\end{problem}

\begin{solution}
Since \(A\ge0\) and \(x\in\Delta\) has nonnegative coordinates, \(Ax\ge0\). By assumption,
\(s(x)>0\), so \(F\) is well defined and has nonnegative coordinates. Moreover,
\[
\sum_{i=1}^n F_i(x)
=
\frac{\sum_i(Ax)_i}{s(x)}
=
1.
\]
Hence \(F(x)\in\Delta\).

The map \(x\mapsto Ax\) is linear and continuous, and \(s(x)\) is continuous and strictly
positive on \(\Delta\). Therefore \(F\) is continuous.

The simplex \(\Delta\) is compact and convex. Brouwer's fixed-point theorem gives
\(x\in\Delta\) such that
\[
F(x)=x.
\]
Thus
\[
Ax=s(x)x.
\]
Therefore \(x\neq0\) is an eigenvector of \(A\) with eigenvalue
\[
\lambda=s(x)=\sum_i(Ax)_i\ge0.
\]
Since \(x\in\Delta\), the eigenvector is nonnegative.
\end{solution}

\begin{problem}[CP-I-0045]
\label{prob:cp-i-0045}
\textbf{Three Equivalent Tests for Diagonalizability}

\textbf{Canonical authored completion of the retained diagonalizability fragment.}

\hspace*{1.5em}Let \(A\in M_n(\mathbb C)\). Prove that the following are equivalent:
\[
\text{(i) }A\text{ is diagonalizable};
\]
\[
\text{(ii) }\sum_{\lambda}\dim\ker(A-\lambda I)=n;
\]
\[
\text{(iii) the minimal polynomial of }A\text{ is a product of distinct linear factors}.
\]
\end{problem}

\begin{solution}
If \(A\) is diagonalizable, it has a basis of eigenvectors. Grouping this basis by eigenvalue gives
\[
\mathbb C^n=\bigoplus_\lambda E_\lambda(A),
\]
so
\[
\sum_\lambda\dim E_\lambda(A)=n.
\]
Thus (i) implies (ii). Conversely, eigenspaces for distinct eigenvalues are linearly independent,
so if their dimensions sum to \(n\), their union contains a basis of eigenvectors. Hence (ii)
implies (i).

For (i) and (iii), if \(A=PDP^{-1}\) with \(D\) diagonal, then the minimal polynomial is the
product
\[
\prod_{\lambda\in\sigma(A)}(t-\lambda),
\]
with no repeated factor. Conversely, if the minimal polynomial splits with no repeated factor,
the primary decomposition has no nontrivial nilpotent part; equivalently every Jordan block has
size \(1\). Hence \(A\) is diagonalizable.
\end{solution}

\begin{problem}[CP-I-0046]
\label{prob:cp-i-0046}
\textbf{Odd-dimensional orthogonal matrices with \(\det=1\) have eigenvalue \(1\)}

\textbf{Exact statement.}
\hspace*{1.5em}Let \(Q\) be a \((2n+1)\times(2n+1)\) orthogonal matrix with \(\det Q=1\).
Show that \(1\) is an eigenvalue of \(Q\) and give a geometric interpretation of this result when \(Q\) is a \(3\times3\) matrix.
What can be said if \(\det Q=-1\)?

\end{problem}

\begin{solution}
If \(Q\) is real orthogonal then \(Q^TQ=I\), hence every eigenvalue \(\lambda\) satisfies \(|\lambda|=1\).
Non-real eigenvalues occur in complex conjugate pairs \(\lambda,\overline\lambda\), and each pair contributes product
\(\lambda\overline\lambda=|\lambda|^2=1\) to \(\det Q\).
Since the dimension is odd, at least one eigenvalue is real; real eigenvalues of an orthogonal matrix are \(\pm 1\).
Thus the sign of \(\det Q\) is determined by the number of eigenvalues \(-1\).
If \(\det Q=1\), the number of \(-1\) eigenvalues is even, so among the (odd number of) real eigenvalues there must be at least one \(+1\).
Hence \(1\) is an eigenvalue.

For \(3\times3\), \(\det Q=1\) means \(Q\) is a rotation: the eigenvector with eigenvalue \(1\) is the rotation axis (a fixed direction).
If \(\det Q=-1\), then \(-1\) is an eigenvalue; geometrically this corresponds to an improper orthogonal map (reflection or roto-reflection).
\end{solution}

\begin{problem}[CP-I-0049]
\label{prob:cp-i-0049}
\textbf{Invertibility and inverse via elementary row operations}

\textbf{Exact statement.}
\hspace*{1.5em}Use elementary row operations to prove that the real matrix
\[
A=\begin{pmatrix}
	1&-1&0\\
	0&0&1\\
	0&3&-1
\end{pmatrix}
\]
is invertible, and find \(A^{-1}\).

\end{problem}

\begin{solution}
Row-reduce the augmented matrix \([A\mid I_3]\):
\[
\left[\begin{array}{ccc|ccc}
	1&-1&0&1&0&0\\
	0&0&1&0&1&0\\
	0&3&-1&0&0&1
\end{array}\right]
\overset{R_2\leftrightarrow R_3}{\longrightarrow}
\left[\begin{array}{ccc|ccc}
	1&-1&0&1&0&0\\
	0&3&-1&0&0&1\\
	0&0&1&0&1&0
\end{array}\right]
\]
\[
\overset{R_2\leftarrow \frac13 R_2}{\longrightarrow}
\left[\begin{array}{ccc|ccc}
	1&-1&0&1&0&0\\
	0&1&-\frac13&0&0&\frac13\\
	0&0&1&0&1&0
\end{array}\right]
\overset{R_2\leftarrow R_2+\frac13 R_3}{\longrightarrow}
\left[\begin{array}{ccc|ccc}
	1&-1&0&1&0&0\\
	0&1&0&0&\frac13&\frac13\\
	0&0&1&0&1&0
\end{array}\right]
\]
\[
\overset{R_1\leftarrow R_1+R_2}{\longrightarrow}
\left[\begin{array}{ccc|ccc}
	1&0&0&1&\frac13&\frac13\\
	0&1&0&0&\frac13&\frac13\\
	0&0&1&0&1&0
\end{array}\right].
\]
Thus \(A\) is invertible and
\[
A^{-1}=\begin{pmatrix}
	1&\frac13&\frac13\\
	0&\frac13&\frac13\\
	0&1&0
\end{pmatrix}.
\]
\end{solution}

\begin{problem}[CP-I-0050]
\label{prob:cp-i-0050}
\textbf{Cayley transform of an anti-Hermitian matrix}

\textbf{Exact statement.}
\hspace*{1.5em}If $S$ is a real symmetric matrix and $A$ is a real antisymmetric matrix, show that $A+iS$ is anti-hermitian and deduce that
\[
\det(A+iS-I)\neq 0.
\]
Show that
\[
U=(I+A+iS)(I-A-iS)^{-1}
\]
is unitary. Find $U$ when
\[
S=\begin{pmatrix}1&1\\1&1\end{pmatrix},\qquad
A=\begin{pmatrix}0&1\\-1&0\end{pmatrix},
\]
and show that it has eigenvalues $\pm(1-i)/\sqrt2$.
\end{problem}

\begin{solution}
Let $K=A+iS$. Since $A^T=-A$ and $S^T=S$,
\[
K^\dagger=(A+iS)^\dagger=A^T-iS^T=-A-iS=-K,
\]
so $K$ is anti-hermitian. Hence all eigenvalues of $K$ are purely imaginary. Therefore eigenvalues of $K-I$ are of the form
$\mu-1$ with $\mu\in i\mathbb R$, which can never be $0$; thus $K-I$ is invertible and $\det(K-I)\neq 0$.

Define
\[
U=(I+K)(I-K)^{-1}.
\]
Using $K^\dagger=-K$,
\[
U^\dagger=\bigl((I-K)^{-1}\bigr)^\dagger (I+K)^\dagger
=(I-K^\dagger)^{-1}(I+K^\dagger)
=(I+K)^{-1}(I-K),
\]
and hence
\[
U^\dagger U=(I+K)^{-1}(I-K)(I+K)(I-K)^{-1}=I,
\]
so $U$ is unitary.

For
\[
S=\begin{pmatrix}1&1\\1&1\end{pmatrix},\qquad
A=\begin{pmatrix}0&1\\-1&0\end{pmatrix},
\]
a direct computation gives
\[
U=\begin{pmatrix}0&i\\-1&0\end{pmatrix}.
\]
Its characteristic polynomial is
\[
\det(U-\lambda I)=\det\begin{pmatrix}-\lambda&i\\-1&-\lambda\end{pmatrix}
=\lambda^2+i,
\]
so $\lambda^2=-i$, and therefore
\[
\lambda=\pm \sqrt{-i}=\pm e^{-i\pi/4}=\pm\frac{1-i}{\sqrt2}.
\]
\end{solution}

\begin{problem}[CP-I-0053]
\label{prob:cp-i-0053}
\textbf{Inverses and transposes}

\hspace*{1.5em}\textbf{(a) Exact statement.}
Let $A$ be an invertible square matrix. Describe the eigenvalues and the characteristic and minimal polynomials of $A^{-1}$ in terms of those of $A$.

\textbf{Solution to (a).}
If $Av=\lambda v$ with $v\neq 0$, then $A^{-1}v=\lambda^{-1}v$. Hence the eigenvalues of $A^{-1}$ are $\lambda^{-1}$, where $\lambda$ runs over eigenvalues of $A$ (with the same algebraic multiplicities).
For $n=\dim A$,
\[
\chi_{A^{-1}}(t)=\det(tI-A^{-1})
=\det(A)^{-1}\det(tA-I)
=\det(A)^{-1}(-1)^n t^n\,\chi_A(1/t).
\]
If
\[
m_A(t)=\prod_{\lambda\neq 0}(t-\lambda)^{s_\lambda},
\]
then
\[
m_{A^{-1}}(t)=\prod_{\lambda\neq 0}(t-\lambda^{-1})^{s_\lambda}.
\]

\medskip
\textbf{(b) Exact statement.}
Prove that the inverse of a Jordan block $J_m(\lambda)$ with $\lambda\neq 0$ has Jordan normal form $J_m(\lambda^{-1})$.
Use this to find the Jordan normal form of $A^{-1}$ for an invertible $A$.

\textbf{Solution to (b).}
Write $J_m(\lambda)=\lambda I+N$ with $N=J_m(0)$ nilpotent ($N^m=0$). Then
\[
J_m(\lambda)^{-1}
=(\lambda I+N)^{-1}
=\lambda^{-1}(I+\lambda^{-1}N)^{-1}
=\lambda^{-1}\sum_{k=0}^{m-1}(-\lambda^{-1}N)^k.
\]
Thus $J_m(\lambda)^{-1}=\lambda^{-1}I+q(N)$ where $q(N)$ is nilpotent and its linear term is nonzero, so $q(N)$ has the same nilpotent Jordan type as $N$.
Hence $J_m(\lambda)^{-1}$ is similar to $\lambda^{-1}I+N=J_m(\lambda^{-1})$.
If
\[
A\sim \bigoplus_i J_{m_i}(\lambda_i)\qquad(\lambda_i\neq 0),
\]
then
\[
A^{-1}\sim \bigoplus_i J_{m_i}(\lambda_i^{-1}).
\]

\medskip
\textbf{(c) Exact statement.}
Prove that any square complex matrix is similar to its transpose.

\end{problem}

\begin{solution}
Let $A=PJP^{-1}$ be a Jordan form over $\mathbb C$. Then
\[
A^T=(P^{-1})^T J^T P^T,
\]
so $A^T$ is similar to $J^T$. For each Jordan block $J_m(\lambda)$, let $R_m$ be the reversal permutation matrix (ones on the anti-diagonal). Then
\[
R_m J_m(\lambda) R_m^{-1}=J_m(\lambda)^T,
\]
so each block is similar to its transpose, hence $J\sim J^T$ and therefore $A\sim A^T$.
\end{solution}

\begin{problem}[CP-I-0054]
\label{prob:cp-i-0054}
\textbf{Minimal polynomial over $\mathbb C$ has real coefficients}

\textbf{Exact statement.}
\hspace*{1.5em}If $A$ is a real $n\times n$ matrix, show that its minimal polynomial over $\mathbb C$ has real coefficients.

\end{problem}

\begin{solution}
Let $m(t)\in\mathbb C[t]$ be the monic minimal polynomial of $A$ over $\mathbb C$, so $m(A)=0$.
Since $A$ has real entries, taking complex conjugates gives $\overline{m}(A)=0$.
Both $m$ and $\overline{m}$ are monic and annihilate $A$. By minimality, $m\mid \overline{m}$.
Conjugating again gives $\overline{m}\mid m$. Hence $m=\overline{m}$, so $m$ has real coefficients.
\end{solution}

\begin{problem}[CP-I-0055]
\label{prob:cp-i-0055}
\textbf{A Hermitian form on $\mathbb C^3$ and a roots-of-unity polarization identity}

\textbf{Exact statement.}
\hspace*{1.5em}Let $\psi:V\times V\to\mathbb C$ be a Hermitian form on a complex vector space $V$.

\textbf{(i)} Find the rank and signature of $\psi$ in the case $V=\mathbb C^3$ and
\[
\psi(x,x)=|x_1+ix_2|^2+|x_2+ix_3|^2+|x_3+ix_1|^2-|x_1+x_2+x_3|^2.
\]

\textbf{(ii)} Show in general that if $n>2$ then
\[
\psi(u,v)=\frac1n\sum_{k=1}^n \zeta^k\,\psi(u+\zeta^k v,\ u+\zeta^k v)
\quad\text{where}\quad \zeta=e^{2\pi i/n}.
\]



\textbf{Quadratic form from a symmetric matrix.}
Over a field $F$ with $\mathrm{char}(F)\neq 2$, a symmetric matrix $A\in M_n(F)$ defines a quadratic form
\[
q(x)=x^T A x,\qquad x\in F^n,
\]
and the associated symmetric bilinear form
\[
b(x,y)=\frac12\bigl(q(x+y)-q(x)-q(y)\bigr)=x^T A y.
\]

\textbf{Change of variables = congruence.}
If we change coordinates $x=P\,x'$ with $P\in GL_n(F)$, then
\[
q(x)=x^T A x = (x')^T (P^T A P)\,x'.
\]
Thus \emph{classification of quadratic forms up to change of basis} is exactly classification of symmetric matrices up to \emph{congruence} $A\sim B \iff B=P^TAP$.

\textbf{Rank.}
$\mathrm{rank}(q):=\mathrm{rank}(A)$ (independent of basis). The \emph{radical} is
\[
\mathrm{rad}(b)=\{x:\ b(x,y)=0\ \forall y\}=\ker(A),
\]
so $\dim \mathrm{rad}(b)=n-\mathrm{rank}(A)$.

\textbf{Diagonalization over $\mathbb R$.}
Every real symmetric matrix is congruent to a diagonal matrix with entries in $\{1,-1,0\}$:
\[
A \sim_{\mathbb R} \mathrm{diag}(\underbrace{1,\dots,1}_{p},
\underbrace{-1,\dots,-1}_{q},
\underbrace{0,\dots,0}_{r}),
\quad p+q+r=n.
\]

\textbf{Law of inertia (invariants).}
The triple $(p,q,r)$ is \emph{invariant} under congruence and uniquely determined by $A$.
In particular:
\[
\mathrm{rank}(A)=p+q,\qquad \text{signature}(A)=p-q,
\]
and $A$ is positive definite iff $q=r=0$ (equivalently $p=n$).

\textbf{Quick definiteness tests (minimum computation).}
For $2\times 2$ symmetric $A=\begin{pmatrix}a&b\\b&c\end{pmatrix}$:
\[
A>0\ \Longleftrightarrow\ a>0\ \text{and}\ \det(A)>0.
\]
More generally, \emph{Sylvester's criterion} says:
\[
A>0\ \Longleftrightarrow\ \Delta_k>0\ \text{for all leading principal minors}\ \Delta_k.
\]

\textbf{Why this solves Problem 1 over $\mathbb R$.}
$I_2$ has inertia $(2,0,0)$, so $A\sim_{\mathbb R}I_2$ iff $A$ has $(2,0,0)$, i.e.\ is positive definite.

\textbf{Key fact (classification over $\mathbb C$).}
Over an algebraically closed field of characteristic $\neq 2$ (e.g.\ $\mathbb C$),
every symmetric matrix is congruent to a diagonal matrix with entries in $\{1,0\}$:
\[
A \sim_{\mathbb C} \mathrm{diag}(\underbrace{1,\dots,1}_{r},\underbrace{0,\dots,0}_{n-r}),
\qquad r=\mathrm{rank}(A).
\]
In particular, any \emph{nonsingular} symmetric matrix ($r=n$) is congruent to $I_n$.

\textbf{Reason (one-line).}
Diagonalize to $\mathrm{diag}(\lambda_1,\dots,\lambda_n)$ and rescale each coordinate by $\sqrt{\lambda_i}$,
possible because $\mathbb C$ contains square roots of every nonzero number.

Over $\mathbb Q$, congruence classification is \emph{arithmetical} (not just rank/signature).
Two indispensable invariants:

\textbf{1) Determinant mod squares (discriminant).}
If $B=P^TAP$, then
\[
\det(B)=\det(P)^2\det(A),
\]
so $\det(A)$ is preserved in $\mathbb Q^*/(\mathbb Q^*)^2$.
Thus if $A\sim_{\mathbb Q}I_n$ then $\det(A)$ must be a square in $\mathbb Q^*$.

\textbf{2) Hasse invariants / Hilbert symbols (deeper).}
For a diagonal form $q=\langle a_1,\dots,a_n\rangle$ over $\mathbb Q$,
the \emph{Hasse invariant} is
\[
\epsilon(q)=\prod_{i<j}(a_i,a_j)_{\mathbb Q},
\]
a product of Hilbert symbols. Together with rank and discriminant, it completes the classification over $\mathbb Q$
(Hasse--Minkowski principle: a rational quadratic form is determined by its local data over $\mathbb R$ and $\mathbb Q_p$).

\textbf{Practical takeaway for small $2\times 2$.}
Often it suffices to use:
\[
\det \text{ mod squares} \quad + \quad \text{whether a represented value is possible (e.g.\ }a^2+b^2=-1\text{)}.
\]

\textbf{Core technique.}
Write the quadratic form as a sum/difference of squares using linear substitutions:
\[
q(x)=\sum \pm (\ell_i(x))^2.
\]
Each $+(\ell_i)^2$ contributes one positive direction; each $-(\ell_i)^2$ contributes one negative direction.
Any variable that disappears contributes to the nullity (rank drops).

\textbf{Matrix language (LDL$^T$).}
Over $\mathbb R$, symmetric Gaussian elimination gives
\[
A = L D L^T,
\]
where $L$ is unit lower triangular and $D$ is block diagonal with $1\times 1$ and $2\times 2$ blocks.
The signs of $D$ encode inertia. Completing squares is the hand version of this.

\textbf{Frobenius product.}
On $\mathrm{Mat}_n(\mathbb R)$ define
\[
\langle A,B\rangle_F := \mathrm{tr}(AB^T).
\]
Then
\[
\langle A,A\rangle_F = \mathrm{tr}(AA^T)=\sum_{i,j}a_{ij}^2=\|A\|_F^2>0\ (A\neq 0),
\]
so it is an inner product.

\textbf{Cauchy--Schwarz.}
\[
|\mathrm{tr}(AB^T)| = |\langle A,B\rangle_F|
\le \|A\|_F\,\|B\|_F
= \mathrm{tr}(AA^T)^{1/2}\mathrm{tr}(BB^T)^{1/2}.
\]

\textbf{Geometric meaning.}
$\mathrm{Mat}_n(\mathbb R)\cong \mathbb R^{n^2}$ with the Euclidean dot product on entries:
$\mathrm{tr}(AB^T)=\sum_{i,j}a_{ij}b_{ij}$.

\subsubsection*{G. The quadratic form $q(A)=\mathrm{tr}(A^2)$ and the symmetric/skew split}

\textbf{Why it is quadratic.}
$q(tA)=t^2q(A)$ and its polarization is
\[
b(A,B)=\frac12(q(A+B)-q(A)-q(B))=\mathrm{tr}(AB),
\]
a symmetric bilinear form on $\mathrm{Mat}_n(\mathbb R)$.

\textbf{Decomposition.}
Every real matrix decomposes uniquely as
\[
A=S+K,\qquad S^T=S,\quad K^T=-K.
\]
Moreover the subspaces are orthogonal for $\mathrm{tr}(AB)$:
\[
\mathrm{tr}(SK)=0.
\]

\textbf{Sign behavior.}
Eigenvalues of $S$ are real, so $\mathrm{tr}(S^2)=\sum \lambda_i(S)^2>0$ if $S\neq 0$.
Eigenvalues of $K$ are purely imaginary in conjugate pairs $\pm i\mu$, giving $\mathrm{tr}(K^2)=-2\sum \mu^2<0$ if $K\neq 0$.
Hence
\[
\mathrm{tr}(A^2)=\mathrm{tr}(S^2)+\mathrm{tr}(K^2)
\]
has inertia
\[
\Bigl(\dim \mathrm{Sym}_n,\ \dim \mathrm{Skew}_n\Bigr)
=
\left(\frac{n(n+1)}2,\ \frac{n(n-1)}2\right),
\]
and rank $n^2$.

\textbf{Skew-symmetric form.}
$\varphi(u,v)=-\varphi(v,u)$ corresponds to a matrix $A$ with $A^T=-A$.

\textbf{Even dimension obstruction.}
If $A^T=-A$ then
\[
\det(A)=\det(A^T)=\det(-A)=(-1)^n\det(A),
\]
so if $n$ is odd then $\det(A)=0$.
Thus a \emph{nondegenerate} skew form forces $n=2m$.

\textbf{Darboux (canonical) form.}
There exists a basis $(e_1,f_1,\dots,e_m,f_m)$ such that
\[
[\varphi]=\mathrm{diag}(J,\dots,J),\qquad J=\begin{pmatrix}0&1\\-1&0\end{pmatrix}.
\]
This is the linear-algebra version of Darboux's theorem in symplectic geometry.

\textbf{Isotropic and Lagrangian subspaces.}
A subspace $W$ is \emph{isotropic} if $\varphi|_W\equiv 0$, equivalently $W\subseteq W^\perp$.
Nondegeneracy gives $\dim W+\dim W^\perp=2m$, hence $\dim W\le m$.
If $\dim W=m$, then $W$ is \emph{maximal isotropic} (a \emph{Lagrangian} subspace).
So the maximum dimension on which $\varphi$ vanishes is $m=n/2$.

\textbf{Hermitian form.}
A Hermitian form $\psi:V\times V\to\mathbb C$ satisfies
\[
\psi(u,v)=\overline{\psi(v,u)},
\]
linear in the first variable and conjugate-linear in the second (convention used here).

\textbf{Hermitian matrices.}
On $V=\mathbb C^n$, $\psi(x,y)=x^*Hy$ with $H^*=H$.
Then $\psi(x,x)=x^*Hx\in \mathbb R$.

\textbf{Unitary diagonalization and inertia.}
There exists a unitary $U$ such that
\[
U^*HU=\mathrm{diag}(\lambda_1,\dots,\lambda_n),\qquad \lambda_i\in\mathbb R.
\]
Thus Hermitian congruence classification parallels the real case:
one defines inertia $(p,q,r)$ by the signs of eigenvalues ($+$, $-$, $0$),
and $p,q,r$ are invariants (signature is $p-q$).

\textbf{Quick strategy for structured $3\times 3$ examples.}
If $H$ is circulant, the Fourier basis diagonalizes it immediately.
If $H$ has a ``sum'' direction $(1,1,1)$, test it first; then restrict to its orthogonal complement.

\textbf{Real/complex standard polarization.}
For a Hermitian form,
\[
\psi(u,v)=\frac14\Bigl(\psi(u+v,u+v)-\psi(u-v,u-v)\Bigr)
+\frac{i}{4}\Bigl(\psi(u+iv,u+iv)-\psi(u-iv,u-iv)\Bigr).
\]
This shows that knowing $x\mapsto \psi(x,x)$ determines $\psi$.

\textbf{Roots-of-unity averaging (Problem 5(ii)).}
Let $\zeta=e^{2\pi i/n}$ with $n>2$.
Expanding $\psi(u+\zeta^k v,u+\zeta^k v)$ and summing with weights $\zeta^k$ uses:
\[
\sum_{k=1}^n \zeta^k=0,\qquad \sum_{k=1}^n \zeta^{2k}=0\ (n>2),\qquad \sum_{k=1}^n 1=n,
\]
so all terms cancel except the $\psi(u,v)$-term, yielding
\[
\psi(u,v)=\frac1n\sum_{k=1}^n \zeta^k\,\psi(u+\zeta^k v,\ u+\zeta^k v).
\]
Conceptually: you project out the ``mixed'' coefficient via Fourier orthogonality.\end{problem}

\begin{solution}
\textbf{(i)} Expanding $|\,\cdot\,|^2$ and collecting coefficients of $x_j\overline{x_k}$ gives $\psi(x,x)=x^*Hx$ with Hermitian matrix
\[
H=\begin{pmatrix}
	1&-1-i&-1+i\\
	-1+i&1&-1-i\\
	-1-i&-1+i&1
\end{pmatrix}.
\]
This is circulant, hence diagonalized by the Fourier vectors
\[
(1,1,1),\quad (1,\omega,\omega^2),\quad (1,\omega^2,\omega),\qquad \omega=e^{2\pi i/3}.
\]
The corresponding eigenvalues are
\[
-1,\qquad 2+\sqrt3,\qquad 2-\sqrt3,
\]
(in some order for the last two). Hence $\psi$ has rank $3$ and inertia/signature $(2,1)$ (two positive, one negative).

\textbf{(ii)} For any $\alpha\in\mathbb C$, sesquilinearity (Hermitian: linear in the first slot, conjugate-linear in the second) gives
\[
\psi(u+\alpha v,u+\alpha v)=\psi(u,u)+\alpha\,\psi(v,u)+\overline{\alpha}\,\psi(u,v)+|\alpha|^2\psi(v,v).
\]
Take $\alpha=\zeta^k$ (so $|\zeta^k|=1$), multiply by $\zeta^k$, and sum over $k=1,\dots,n$:
\[
\sum_{k=1}^n \zeta^k\psi(u+\zeta^k v,u+\zeta^k v)
=
\psi(u,u)\sum_{k=1}^n \zeta^k
+\psi(v,u)\sum_{k=1}^n \zeta^{2k}
+\psi(u,v)\sum_{k=1}^n 1
+\psi(v,v)\sum_{k=1}^n \zeta^k.
\]
For $n>2$ we have $\sum_{k=1}^n\zeta^k=0$ and $\sum_{k=1}^n\zeta^{2k}=0$, while $\sum_{k=1}^n 1=n$.
Thus the right-hand side equals $n\,\psi(u,v)$, proving
\[
\psi(u,v)=\frac1n\sum_{k=1}^n \zeta^k\,\psi(u+\zeta^k v,\ u+\zeta^k v).
\]
\end{solution}

\begin{problem}[CP-I-0058]
\label{prob:cp-i-0058}
\textbf{A symmetric matrix of finite order is an involution}

\textbf{Exact statement.}
\hspace*{1.5em}Let $S$ be an $n\times n$ real symmetric matrix with $S^k=I_n$ for some $k\ge 1$.
Show that $S^2=I_n$.

\end{problem}

\begin{solution}
Since $S$ is real symmetric, it is orthogonally diagonalizable:
\[
S=U\Lambda U^T,\qquad \Lambda=\mathrm{diag}(\lambda_1,\dots,\lambda_n),\ \lambda_i\in\mathbb R.
\]
Then
\[
I=S^k=U\Lambda^kU^T \quad\Rightarrow\quad \Lambda^k=I \quad\Rightarrow\quad \lambda_i^k=1\ \ \forall i.
\]
A real number $\lambda$ with $\lambda^k=1$ must satisfy $\lambda=1$ (always) or $\lambda=-1$ (only if $k$ is even).
Thus each $\lambda_i\in\{\pm 1\}$, so $\Lambda^2=I$ and hence
\[
S^2=U\Lambda^2U^T=I.
\]
\end{solution}

\begin{problem}[CP-I-0059]
\label{prob:cp-i-0059}
\textbf{Eigenvalues, eigenspace dimensions, diagonalizability}

\textbf{Exact statement.}
\hspace*{1.5em}For each of the matrices
\[
A=\begin{pmatrix}
	5&-3&2\\
	6&-4&4\\
	4&-4&5
\end{pmatrix},\qquad
B=\begin{pmatrix}
	1&-3&4\\
	4&-7&8\\
	6&-7&7
\end{pmatrix},\qquad
C=\begin{pmatrix}
	7&-12&6\\
	10&-19&10\\
	12&-24&13
\end{pmatrix},
\]
(i) compute their eigenvalues;
(ii) for each eigenvalue \(\lambda\) compute \(\dim\{x\in\mathbb R^3: Mx=\lambda x\}\);
(iii) hence determine whether or not the matrix is diagonalizable.

\end{problem}

\begin{solution}
\textbf{Matrix \(A\).}
\[
\chi_A(\lambda)=\det(\lambda I-A)=\lambda^3-6\lambda^2+11\lambda-6=(\lambda-1)(\lambda-2)(\lambda-3).
\]
Hence the eigenvalues are \(1,2,3\), each with algebraic multiplicity \(1\), so each eigenspace has dimension \(1\).
For example,
\[
\lambda=1:\ (1,2,1),\qquad
\lambda=2:\ (1,1,0),\qquad
\lambda=3:\ (1,2,2)
\]
are eigenvectors. Therefore \(A\) is diagonalizable.

\textbf{Matrix \(B\).}
\[
\chi_B(\lambda)=\lambda^3-\lambda^2-5\lambda-3=(\lambda-3)(\lambda+1)^2.
\]
So the eigenvalues are \(3\) and \(-1\) (with algebraic multiplicity \(2\)).
One finds
\[
E_3=\ker(B-3I)=\operatorname{span}\{(1,2,2)\},\qquad \dim E_3=1,
\]
\[
E_{-1}=\ker(B+I)=\operatorname{span}\{(1,2,1)\},\qquad \dim E_{-1}=1.
\]
Since \(\dim E_{-1}=1<2\), the matrix \(B\) is not diagonalizable.

\textbf{Matrix \(C\).}
\[
\chi_C(\lambda)=\lambda^3-\lambda^2-\lambda+1=(\lambda-1)^2(\lambda+1).
\]
So the eigenvalues are \(1\) (algebraic multiplicity \(2\)) and \(-1\).
One finds
\[
E_1=\ker(C-I)=\operatorname{span}\{(2,1,0),\,(-1,0,1)\},\qquad \dim E_1=2,
\]
\[
E_{-1}=\ker(C+I)=\operatorname{span}\{(3,5,6)\},\qquad \dim E_{-1}=1.
\]
Thus \(\dim E_1+\dim E_{-1}=2+1=3\), so \(C\) is diagonalizable.
\end{solution}

\begin{problem}[CP-I-0064]
\label{prob:cp-i-0064}
\textbf{Modulus–argument forms on the unit circle; a Cayley transform locus}

\textbf{Exact statement.}
\hspace*{1.5em}Given \(|z|=1\) and \(\arg z=\theta\), find both algebraically and geometrically the modulus-argument forms of
(i) \(1+z\), (ii) \(1-\bar z\).

Show that the locus of \(w\) as \(z\) varies with \(|z|=1\), where \(w\) is given by
\[
w^2=\frac{1-z}{1+z},
\]
is a pair of straight lines.

\end{problem}

\begin{solution}
Since \(|z|=1\) and \(\arg z=\theta\), write \(z=e^{i\theta}\) and \(\bar z=e^{-i\theta}\).

\medskip
\noindent\textbf{(i) Modulus–argument of \(1+z\).}
\[
1+z=1+e^{i\theta}=e^{i\theta/2}\bigl(e^{-i\theta/2}+e^{i\theta/2}\bigr)
=2\cos\!\left(\frac\theta2\right)e^{i\theta/2}.
\]
So \(|1+z|=2\bigl|\cos(\theta/2)\bigr|\), and an argument is \(\theta/2\) (adjust by \(\pi\) if \(\cos(\theta/2)<0\)).

Geometrically: \(1+z\) is the diagonal of the parallelogram with sides \(1\) and \(z\); its direction bisects the angle between them, hence angle \(\theta/2\).

\medskip
\noindent\textbf{(ii) Modulus–argument of \(1-\bar z\).}
\[
1-\bar z = 1-e^{-i\theta}
= e^{-i\theta/2}\bigl(e^{i\theta/2}-e^{-i\theta/2}\bigr)
=2i\sin\!\left(\frac\theta2\right)e^{-i\theta/2}.
\]
Equivalently,
\[
1-\bar z = 2\sin\!\left(\frac\theta2\right)\,e^{\,i(\pi/2-\theta/2)}.
\]
So \(|1-\bar z|=2\bigl|\sin(\theta/2)\bigr|\), and an argument is \(\pi/2-\theta/2\) (adjust by \(\pi\) if \(\sin(\theta/2)<0\)).

Geometrically: \(|1-\bar z|=|1-e^{-i\theta}|\) is the chord length between \(1\) and \(e^{-i\theta}\) on the unit circle.

\medskip
\noindent\textbf{Locus of \(w\).}
For \(z=e^{i\theta}\) (excluding \(z=-1\)),
\[
\frac{1-z}{1+z}
=\frac{1-e^{i\theta}}{1+e^{i\theta}}
=\frac{e^{-i\theta/2}-e^{i\theta/2}}{e^{-i\theta/2}+e^{i\theta/2}}
=\frac{-2i\sin(\theta/2)}{2\cos(\theta/2)}
=-i\tan\!\left(\frac\theta2\right).
\]
Thus \(w^2\in i\mathbb{R}\) (purely imaginary). Write \(w=u+iv\). Then
\[
w^2=(u^2-v^2)+2uvi
\]
is purely imaginary iff \(u^2-v^2=0\), i.e. \(v=\pm u\). Therefore \(w\) lies on the two straight lines
\[
v=u \quad\text{or}\quad v=-u.
\]
\end{solution}

\begin{problem}[CP-I-0065]
\label{prob:cp-i-0065}
\textbf{Non-similarity of unipotent matrices; identify a given matrix}

\textbf{Exact statement.}
\hspace*{1.5em}Without appealing directly to the uniqueness of Jordan normal form show that none of the following real matrices are similar:
\[
U_1=\begin{pmatrix}1&1&0\\0&1&1\\0&0&1\end{pmatrix},\quad
U_2=\begin{pmatrix}1&1&0\\0&1&0\\0&0&1\end{pmatrix},\quad
U_3=\begin{pmatrix}1&0&0\\0&1&0\\0&0&1\end{pmatrix}.
\]
Is the matrix
\[
A=\begin{pmatrix}-2&-2&-1\\ 3&3&1\\ 3&2&2\end{pmatrix}
\]
similar to any of them? If so, which? Find a basis such that it is in Jordan normal form.

\end{problem}

\begin{solution}
If $B=P^{-1}AP$, then for every $k\ge 1$,
\[
(B-I)^k=P^{-1}(A-I)^kP,
\]
so $\dim\ker(A-I)^k$ (and $\operatorname{rank}(A-I)^k$) are similarity invariants.

Let $N_i=U_i-I$.
For $U_3=I$ we have $N_3=0$, hence $\dim\ker(N_3)=3$.
For $U_2$ one has $N_2\neq 0$ but $N_2^2=0$; moreover $\operatorname{rank}(N_2)=1$, hence $\dim\ker(N_2)=2$.
For $U_1$ one has $N_1\neq 0$, $N_1^2\neq 0$, but $N_1^3=0$; in particular $\dim\ker(N_1)=1$ and $\dim\ker(N_1^2)=2$.
Since these kernel dimensions differ, the matrices $U_1,U_2,U_3$ are pairwise non-similar.

Now consider $A$. Its characteristic polynomial is
\[
\chi_A(t)=(t-1)^3,
\]
so the only eigenvalue is $1$. Let $N=A-I$. One checks that $N\neq 0$ and $N^2=0$, and $\operatorname{rank}(N)=1$,
so $\dim\ker(N)=2$. This matches the invariants of $U_2$, hence $A$ is similar to $U_2$.

To produce a Jordan basis, take
\[
w=\begin{pmatrix}1\\0\\0\end{pmatrix},\qquad
u=(A-I)w=\begin{pmatrix}-3\\3\\3\end{pmatrix}.
\]
Then $(A-I)u=0$ since $(A-I)^2=0$, so $u$ is an eigenvector.
Choose an independent eigenvector $z\in\ker(A-I)$, e.g.
\[
z=\begin{pmatrix}-2\\3\\0\end{pmatrix}\in\ker(A-I).
\]
In the basis $(u,w,z)$ we have $Au=u$, $Aw=w+u$, $Az=z$, hence
\[
[A]_{(u,w,z)}=\begin{pmatrix}1&1&0\\0&1&0\\0&0&1\end{pmatrix}.
\]
\end{solution}

\begin{problem}[CP-I-0067]
\label{prob:cp-i-0067}
\textbf{Eigenvalues and determinant of the ``$\lambda$ on diagonal, $1$ elsewhere'' matrix}

\textbf{Exact statement.}
\hspace*{1.5em}Find the eigenvalues of the $n\times n$ real matrix $A$ with diagonal entries $\lambda$ and all other entries $1$.
Hence compute $\det(A)$.

\end{problem}

\begin{solution}
Let $J$ be the all-ones matrix. Then $A=(\lambda-1)I+J$.
We have $J\mathbf 1=n\mathbf 1$ for $\mathbf 1=(1,\dots,1)^T$, and $Jx=0$ for all $x$ with $\sum_i x_i=0$.
Hence the eigenvalues of $A$ are
\[
\lambda+n-1 \quad (\text{mult. }1),\qquad \lambda-1 \quad (\text{mult. }n-1).
\]
Thus
\[
\det(A)=(\lambda+n-1)(\lambda-1)^{n-1}.
\]
\end{solution}

\begin{problem}[CP-I-0070]
\label{prob:cp-i-0070}
\textbf{$2\times2$ matrices for plane transformations and complex multiplication}

\textbf{Exact statement.}

	\par\hspace*{1.5em}\textbullet\quad Give examples of $2\times 2$ real matrices representing the following types of transformations in $\mathbb R^2$:
	(i) reflection; (ii) dilation (or scaling); (iii) shear; and (iv) rotation.
	Which of these types of transformation are always represented by a $2\times 2$ matrix with determinant $+1$?
	For which types (i)--(iv) do transformations $A$ and $B$ of the same type obey $AB=BA$ in general?
	\par\noindent\textbullet\quad A linear map $\mathbb R^2\to\mathbb R^2$ with $\mathbf x\mapsto \mathbf x'=M\mathbf x$ is defined by
	$z' = cz$ where $z=x_1+i x_2$, $z'=x_1'+i x_2'$ and $c=a+ib$ is a fixed complex number.
	Find the real $2\times 2$ matrix $M$ in terms of $a$ and $b$.
	Which types of transformations (i)--(iv) can be obtained for particular choices of $c=a+ib$?

\end{problem}

\begin{solution}
\par\noindent\textbullet\quad Examples:
	\[
	\text{reflection across $x$-axis: }\;
	\begin{pmatrix}1&0\\0&-1\end{pmatrix}\quad(\det=-1),
	\]
	\[
	\text{dilation by }s:\;
	\begin{pmatrix}s&0\\0&s\end{pmatrix}\quad(\det=s^2),
	\]
	\[
	\text{shear (horizontal) by }k:\;
	\begin{pmatrix}1&k\\0&1\end{pmatrix}\quad(\det=1),
	\]
	\[
	\text{rotation by }\theta:\;
	\begin{pmatrix}\cos\theta&-\sin\theta\\ \sin\theta&\cos\theta\end{pmatrix}\quad(\det=1).
	\]
	Thus rotations and (standard) shears always have determinant $+1$; reflections always have determinant $-1$; dilations have determinant $s^2$ (not always $1$).

	Commutation: dilations commute (scalar matrices), and rotations in $\mathbb R^2$ commute:
	$Rot_\alpha Rot_\beta=Rot_{\alpha+\beta}=Rot_\beta Rot_\alpha$.
	Reflections do not commute in general, and shears do not commute in general (only shears along the same direction commute).
	\par\noindent\textbullet\quad Compute
	\[
	z'=(a+ib)(x_1+i x_2)=(ax_1-bx_2)+i(bx_1+ax_2).
	\]
	Hence
	\[
	\begin{pmatrix}x_1'\\x_2'\end{pmatrix}
	=
	\begin{pmatrix}a&-b\\ b&a\end{pmatrix}
	\begin{pmatrix}x_1\\x_2\end{pmatrix},
	\qquad
	M=\begin{pmatrix}a&-b\\ b&a\end{pmatrix},
	\qquad
	\det M=a^2+b^2.
	\]
	Writing $c=re^{i\varphi}$ gives
	$M=r\begin{pmatrix}\cos\varphi&-\sin\varphi\\ \sin\varphi&\cos\varphi\end{pmatrix}$,
	i.e.\ a rotation by $\varphi$ and a uniform scaling by $r$.
	Thus one obtains dilations (take $b=0$, $a>0$) and rotations (take $a^2+b^2=1$), but not reflections (which require $\det=-1$) nor shears.
\end{solution}

\begin{problem}[CP-I-0071]
\label{prob:cp-i-0071}
\textbf{Eigenvalues/eigenspaces; simultaneous diagonalisation}

\textbf{Exact statement.}
\hspace*{1.5em}Find the eigenvalues and give bases for the eigenspaces of
\[
A_1=\begin{pmatrix}1&1&0\\0&3&-2\\0&1&0\end{pmatrix},\quad
A_2=\begin{pmatrix}1&1&-1\\0&3&-2\\0&1&0\end{pmatrix},\quad
A_3=\begin{pmatrix}1&1&-1\\-1&3&-1\\-1&1&1\end{pmatrix}.
\]
The second and third matrices commute; find a basis with respect to which they are both diagonal.

\end{problem}

\begin{solution}
For $A_1$,
\[
\chi_{A_1}(t)=\det(A_1-tI)=(t-2)(t-1)^2.
\]
Hence $\lambda=1,2$. Solving $(A_1-I)x=0$ gives
\[
E_1(A_1)=\ker(A_1-I)=\mathrm{span}\{(1,0,0)^T\},
\]
and solving $(A_1-2I)x=0$ gives
\[
E_2(A_1)=\ker(A_1-2I)=\mathrm{span}\{(2,2,1)^T\}.
\]
(So $A_1$ is not diagonalizable since $\dim E_1=1<2$.)

For $A_2$,
\[
\chi_{A_2}(t)=(t-2)(t-1)^2.
\]
Solving yields
\begin{center}
\resizebox{\linewidth}{!}{$\displaystyle
E_1(A_2)=\ker(A_2-I)=\mathrm{span}\{(1,0,0)^T,\,(0,1,1)^T\},
\quad
E_2(A_2)=\ker(A_2-2I)=\mathrm{span}\{(1,2,1)^T\}.
$}
\end{center}
Thus $A_2$ is diagonalizable.

For $A_3$,
\[
\chi_{A_3}(t)=(t-1)(t-2)^2,
\]
and
\begin{center}
\resizebox{\linewidth}{!}{$\displaystyle
E_1(A_3)=\ker(A_3-I)=\mathrm{span}\{(1,1,1)^T\},
\quad
E_2(A_3)=\ker(A_3-2I)=\mathrm{span}\{(1,1,0)^T,\,(-1,0,1)^T\}.
$}
\end{center}
Thus $A_3$ is diagonalizable.

\medskip
\textbf{Simultaneous diagonalisation of $A_2$ and $A_3$.}
Since $A_2A_3=A_3A_2$, the operator $A_3$ preserves each eigenspace of $A_2$.
A common eigenbasis is
\[
u=(1,1,1)^T,\quad w=(0,1,1)^T,\quad v=(1,2,1)^T,
\]
with
\[
A_2u=u,\ A_3u=u;\qquad A_2w=w,\ A_3w=2w;\qquad A_2v=2v,\ A_3v=2v.
\]
Therefore, in the basis $(u,w,v)$,
\[
[A_2]=\mathrm{diag}(1,1,2),\qquad [A_3]=\mathrm{diag}(1,2,2).
\]
\end{solution}

\begin{problem}[CP-I-0075]
\label{prob:cp-i-0075}
\textbf{Diagonalizability and the Direct Sum of Eigenspaces}

\textbf{Canonical authored completion of the retained diagonalization fragment.}

\hspace*{1.5em}Let \(A\in\operatorname{End}(V)\) on a finite-dimensional vector space over a field
for which the characteristic polynomial of \(A\) splits.
Prove that \(A\) is diagonalizable if and only if
\[
V=\bigoplus_{\lambda}E_\lambda(A).
\]
Equivalently, show that
\[
A\text{ is diagonalizable}
\quad\Longleftrightarrow\quad
\sum_{\lambda}\dim E_\lambda(A)=\dim V.
\]
\end{problem}

\begin{solution}
If \(A\) is diagonalizable, \(V\) has a basis consisting of eigenvectors. Grouping the basis vectors
according to their eigenvalues gives
\[
V=\sum_\lambda E_\lambda(A).
\]
Eigenspaces corresponding to distinct eigenvalues are linearly independent, so the sum is direct:
\[
V=\bigoplus_\lambda E_\lambda(A).
\]

Conversely, if
\[
V=\bigoplus_\lambda E_\lambda(A),
\]
choose a basis of each eigenspace and concatenate these bases. Their union is a basis of \(V\)
consisting entirely of eigenvectors of \(A\). Therefore the matrix of \(A\) in this basis is diagonal,
so \(A\) is diagonalizable.

Taking dimensions of the direct sum gives the equivalent criterion
\[
\sum_\lambda \dim E_\lambda(A)=\dim V.
\]
\end{solution}

\begin{problem}[CP-I-0077]
\label{prob:cp-i-0077}
\textbf{One eigenvalue, not diagonalizable}

\textbf{Exact statement.}
\hspace*{1.5em}Show that
\[
M=\begin{pmatrix}
	0&1&0\\
	-4&4&0\\
	-2&1&2
\end{pmatrix}
\]
has characteristic equation $(t-2)^3=0$. Explain, as simply as possible, why $M$ is not diagonalizable.

\end{problem}

\begin{solution}
Compute
\[
tI-M=\begin{pmatrix}
	t&-1&0\\
	4&t-4&0\\
	2&-1&t-2
\end{pmatrix}.
\]
Expanding along the third column gives
\[
\det(tI-M)=(t-2)\det\begin{pmatrix}t&-1\\4&t-4\end{pmatrix}
=(t-2)\bigl(t(t-4)+4\bigr)
=(t-2)(t^2-4t+4)=(t-2)^3.
\]
Thus $2$ is the only eigenvalue, with algebraic multiplicity $3$.

To test diagonalizability, compute the eigenspace:
\[
M-2I=\begin{pmatrix}
	-2&1&0\\
	-4&2&0\\
	-2&1&0
\end{pmatrix}.
\]
All rows are multiples of $(-2,1,0)$, so $\operatorname{rank}(M-2I)=1$, hence
\[
\dim\ker(M-2I)=3-1=2.
\]
The geometric multiplicity is $2<3$, so $M$ does not have $3$ linearly independent eigenvectors
and is not diagonalizable.
\end{solution}

\begin{problem}[CP-I-0078]
\label{prob:cp-i-0078}
\textbf{Anti-hermitian and unitary matrices: eigenvalues and orthogonality}

\textbf{Exact statement.}
\hspace*{1.5em}(i) A matrix $A$ is anti-hermitian, $A^\dagger=-A$; show that the eigenvalues of $A$ are pure-imaginary.
(ii) A matrix $U$ is unitary, $U^\dagger U=I$; show that the eigenvalues of $U$ have unit modulus.
(iii) In each of the cases (i) and (ii), show that eigenvectors with distinct eigenvalues are orthogonal.

\end{problem}

\begin{solution}
\textbf{(i)} Let $Av=\lambda v$ with $v\neq 0$. Then
\[
v^\dagger Av=\lambda\, v^\dagger v.
\]
Taking complex conjugates,
\[
(v^\dagger Av)^* = v^\dagger A^\dagger v = v^\dagger(-A)v = -\,v^\dagger Av.
\]
Thus $z:=v^\dagger Av$ satisfies $z^*=-z$, hence $z$ is purely imaginary. Since $v^\dagger v>0$ is real,
\[
\lambda=\frac{v^\dagger Av}{v^\dagger v}
\]
is purely imaginary.

\textbf{(ii)} If $Uv=\lambda v$, then
\[
v^\dagger v = v^\dagger U^\dagger U v = (Uv)^\dagger(Uv)
=(\lambda v)^\dagger(\lambda v)=|\lambda|^2\, v^\dagger v.
\]
Since $v\neq 0$, we get $|\lambda|=1$.

\textbf{(iii)} For the anti-hermitian case, let $Av=\lambda v$ and $Aw=\mu w$. Then
\[
v^\dagger Aw=\mu\, v^\dagger w,
\]
but also
\[
v^\dagger Aw=(A^\dagger v)^\dagger w = (-Av)^\dagger w = -(\lambda v)^\dagger w
= -\overline{\lambda}\, v^\dagger w.
\]
Hence $(\mu+\overline{\lambda})v^\dagger w=0$. By (i), $\lambda=i\alpha$ and $\mu=i\beta$ with
$\alpha,\beta\in\mathbb R$, so $\mu+\overline{\lambda}=i(\beta-\alpha)$. If $\lambda\neq \mu$ then
$\beta\neq \alpha$, hence $v^\dagger w=0$.

For the unitary case, if $Uv=\lambda v$ then applying $U^\dagger$ gives
$v=\lambda U^\dagger v$, hence $U^\dagger v=\lambda^{-1}v=\overline{\lambda}\,v$ (since $|\lambda|=1$).
Now with $Uw=\mu w$,
\[
v^\dagger Uw=\mu\, v^\dagger w,
\qquad
v^\dagger Uw=(U^\dagger v)^\dagger w=(\overline{\lambda}v)^\dagger w=\lambda\, v^\dagger w.
\]
Thus $(\mu-\lambda)v^\dagger w=0$, and if $\mu\neq\lambda$ then $v^\dagger w=0$.
\end{solution}

\begin{problem}[CP-I-0079]
\label{prob:cp-i-0079}
\textbf{Constructing a Jordan Chain}

\textbf{Canonical authored completion of the retained Jordan-chain fragment.}

\hspace*{1.5em}Let \(A\in M_n(\mathbb C)\), let \(\lambda\) be an eigenvalue, and let
\(v\neq0\) satisfy
\[
(A-\lambda I)v=0.
\]
Suppose there exists \(w\) such that
\[
(A-\lambda I)w=v.
\]
Show that \(v,w\) are linearly independent and that, in the ordered basis \((v,w)\) of their
span, the restriction of \(A\) is represented by the Jordan block
\[
\begin{pmatrix}
\lambda&1\\
0&\lambda
\end{pmatrix}.
\]
Explain how this is the first step in building a Jordan basis.
\end{problem}

\begin{solution}
If \(w=cv\) for some scalar \(c\), then
\[
(A-\lambda I)w
=
c(A-\lambda I)v
=
0,
\]
contradicting
\[
(A-\lambda I)w=v\neq0.
\]
Hence \(v,w\) are linearly independent.

The defining equations give
\[
Av=\lambda v,
\qquad
Aw=\lambda w+v.
\]
Therefore, relative to the ordered basis \((v,w)\), the coordinate columns of \(Av\) and \(Aw\)
are
\[
\binom{\lambda}{0},
\qquad
\binom{1}{\lambda},
\]
so
\[
[A]_{(v,w)}
=
\begin{pmatrix}
\lambda&1\\
0&\lambda
\end{pmatrix}.
\]

More generally, one constructs longer chains
\[
(A-\lambda I)v_1=0,\qquad
(A-\lambda I)v_{j+1}=v_j,
\]
and the chain vectors produce a Jordan block. Combining chains for all eigenvalues yields a
Jordan basis.
\end{solution}

\begin{problem}[CP-I-0080]
\label{prob:cp-i-0080}
\textbf{Column operations on a block matrix and \(\det(AB)=\det(A)\det(B)\)}

\textbf{Exact statement.}
\hspace*{1.5em}Let \(A,B\) be \(n\times n\) matrices over a field \(\mathbb F\).
Show that the \(2n\times 2n\) block matrix
\[
C=\begin{pmatrix}
	I & -B\\
	A & 0
\end{pmatrix}
\]
can be transformed into
\[
D=\begin{pmatrix}
	I & 0\\
	A & AB
\end{pmatrix}
\]
by elementary column operations (specify them).
By considering \(\det(C)\) and \(\det(D)\), obtain another proof that \(\det(AB)=\det(A)\det(B)\).



Row operations on a matrix \(A\) are equivalent to left-multiplication \(EA\) by an elementary matrix \(E\).
Column operations are equivalent to right-multiplication \(AE\).

For square matrices, determinant effects are:
\[
\begin{array}{c|c}
	\text{operation} & \text{effect on }\det \\ \hline
	\text{swap two rows/columns} & \det \mapsto -\det\\
	\text{scale a row/column by }\lambda\neq0 & \det \mapsto \lambda\,\det\\
	\text{add }\lambda\text{ times one row/column to another} & \det\ \text{unchanged}
\end{array}
\]

If \(A\in M_n(\mathbb F)\), row-reducing
\[
[A\mid I] \xrightarrow{\text{row ops}} [EA\mid E]
\]
for some invertible \(E\).
If the left block becomes \(I\), then \(EA=I\), hence \(E=A^{-1}\), and therefore
\[
[A\mid I] \xrightarrow{\text{row ops}} [I\mid A^{-1}].
\]

Let \(V\) be \(n\)-dimensional over \(\mathbb F\). For a basis \(B\), write \([v]_B\in\mathbb F^n\) for coordinates.
For bases \(B,C\), define \([id_V]^C_B\in GL_n(\mathbb F)\) by
\[
[v]_C = [id_V]^C_B\,[v]_B \quad \text{for all }v\in V.
\]
The \(j\)-th column of \([id_V]^C_B\) is \([c_j]_B\).
Composition holds:
\[
[id_V]^D_C\,[id_V]^C_B = [id_V]^D_B.
\]

If \(A\in M_{m\times n}(\mathbb F)\) has column rank \(r\), then \(A\) can be written
\[
A=BC,\qquad B\in M_{m\times r}(\mathbb F),\quad C\in M_{r\times n}(\mathbb F),
\]
by choosing \(B\) as \(r\) independent columns of \(A\) and letting \(C\) record coefficients expressing each column of \(A\)
in terms of those columns.
If \(A=BC\) with \(B\in M_{m\times k}\), then \(\operatorname{colrank}(A)\le k\), so the least such \(k\) equals \(r\).
Transposing gives
\[
A^T=(BC)^T=C^T B^T,
\]
so \(\operatorname{colrank}(A^T)\le r\). Applying the same argument to \(A^T\) and transposing back yields
\(\operatorname{colrank}(A^T)=r\).

If \(X,Y\in M_n(\mathbb F)\), then
\[
\det\begin{pmatrix}X&0\\ *&Y\end{pmatrix}=\det(X)\det(Y),
\qquad
\det\begin{pmatrix}X&*\\ 0&Y\end{pmatrix}=\det(X)\det(Y).
\]
Elementary column operations of the form \(\mathrm{Col}_j\leftarrow \mathrm{Col}_j+\lambda\,\mathrm{Col}_i\) do not change determinant.
This is why the column operations in Problem 4 preserve \(\det\).\end{problem}

\begin{solution}
Write \(C\) in block-columns:
the first \(n\) columns form \(\binom{I}{A}\), the last \(n\) columns form \(\binom{-B}{0}\).
For each \(j=1,\dots,n\), perform the elementary column operation
\[
\mathrm{Col}_{n+j}\ \leftarrow\ \mathrm{Col}_{n+j}+\sum_{k=1}^n b_{kj}\,\mathrm{Col}_k.
\]
Equivalently, in block form: \((\text{right block})\leftarrow(\text{right block})+(\text{left block})\cdot B\).
Then the top-right block becomes \(-B+IB=0\) and the bottom-right block becomes \(0+AB=AB\), giving \(D\).
These operations preserve determinant, so \(\det(C)=\det(D)\).

Since \(D\) is block lower-triangular,
\[
\det(D)=\det(I)\det(AB)=\det(AB).
\]

Now swap the two block-columns of size \(n\) in \(C\), obtaining
\[
C'=\begin{pmatrix}-B&I\\ 0&A\end{pmatrix}.
\]
This swap uses \(n^2\) transpositions, so \(\det(C')=(-1)^{n^2}\det(C)=(-1)^n\det(C)\).
But \(C'\) is block upper-triangular, hence
\[
\det(C')=\det(-B)\det(A)=(-1)^n\det(B)\det(A).
\]
Therefore \(\det(C)=\det(A)\det(B)\).
Combining \(\det(C)=\det(D)\) yields \(\det(AB)=\det(A)\det(B)\).
\end{solution}

\begin{problem}[CP-I-0081]
\label{prob:cp-i-0081}
\textbf{plane transformations and complex multiplication}

\hspace*{1.5em}Orthogonal maps satisfy $A^TA=I$:
\[
\det A=+1 \Rightarrow \text{rotation},\qquad \det A=-1 \Rightarrow \text{reflection (or reflection+rotation)}.
\]
Uniform scaling is $sI$; shears are typically unipotent (determinant $1$) but not angle-preserving.
Complex multiplication $z\mapsto cz$ with $c=a+ib$ corresponds to
\[
\begin{pmatrix}a&-b\\ b&a\end{pmatrix}
=r\begin{pmatrix}\cos\varphi&-\sin\varphi\\ \sin\varphi&\cos\varphi\end{pmatrix},
\]
so it yields rotation + uniform scaling only (no shear, no reflection).
In $\mathbb R^2$ rotations commute; scalings commute with everything; reflections and general shears do not commute in general.
\end{problem}

\begin{problem}[CP-I-0082]
\label{prob:cp-i-0082}
\textbf{Reflections $H(\mathbf n)$ and rotations $R(\mathbf n,\theta)$ in $\mathbb R^3$}

\textbf{Exact statement.}
\hspace*{1.5em}Let $R(\mathbf n,\theta)$ be the matrix corresponding to a rotation with angle $\theta$ and axis $\mathbf n$, as given in $(*)$ of Question~10.
Let $H(\mathbf n)$ be the matrix corresponding to reflection in a plane through the origin with unit normal $\mathbf n$, as defined by
\[
\mathbf x\mapsto \mathbf x'=H(\mathbf n)\mathbf x=\mathbf x-2(\mathbf x\cdot\mathbf n)\mathbf n.
\]
In the following, $\mathbf i,\mathbf j,\mathbf k$ are the standard orthonormal basis vectors in $\mathbb R^3$.

	\par\noindent\textbullet\quad Find explicitly the matrices $R(\mathbf i,\pi/2)$ and $R(\mathbf j,\pi/2)$ and check that
	$R(\mathbf i,\pi/2)R(\mathbf j,\pi/2)\neq R(\mathbf j,\pi/2)R(\mathbf i,\pi/2)$.
	\par\noindent\textbullet\quad Show by both algebraic and geometrical means that the map $\mathbf x\mapsto \mathbf x'=-H(\mathbf n)\mathbf x$ is a rotation through angle $\pi$ about $\mathbf n$.
	\par\noindent\textbullet\quad Given that $\mathbf n_\pm=\cos(\tfrac{\theta}{2})\mathbf i \pm \sin(\tfrac{\theta}{2})\mathbf j$, prove that
	\[
	H(\mathbf i)H(\mathbf n_-)=H(\mathbf n_+)H(\mathbf i)=R(\mathbf k,\theta),
	\]
	and draw diagrams to explain the geometrical meaning of this result.
\end{problem}

\begin{solution}
\par\noindent\textbullet\quad The $90^\circ$ rotation about the $x$-axis is
	\[
	R(\mathbf i,\tfrac{\pi}{2})=
	\begin{pmatrix}
		1&0&0\\
		0&0&-1\\
		0&1&0
	\end{pmatrix},
	\]
	and the $90^\circ$ rotation about the $y$-axis is
	\[
	R(\mathbf j,\tfrac{\pi}{2})=
	\begin{pmatrix}
		0&0&1\\
		0&1&0\\
		-1&0&0
	\end{pmatrix}.
	\]
	Then
	\[
	R(\mathbf i,\tfrac{\pi}{2})R(\mathbf j,\tfrac{\pi}{2})
	=
	\begin{pmatrix}
		0&0&1\\
		1&0&0\\
		0&1&0
	\end{pmatrix}
	\neq
	\begin{pmatrix}
		0&1&0\\
		0&0&-1\\
		-1&0&0
	\end{pmatrix}
	=
	R(\mathbf j,\tfrac{\pi}{2})R(\mathbf i,\tfrac{\pi}{2}),
	\]
	so they do not commute.
	\par\noindent\textbullet\quad Decompose $\mathbf x=\mathbf x_\parallel+\mathbf x_\perp$ where
	$\mathbf x_\parallel=(\mathbf x\cdot\mathbf n)\mathbf n$ and $\mathbf x_\perp\cdot\mathbf n=0$.
	Then
	\[
	H(\mathbf n)\mathbf x=\mathbf x-2\mathbf x_\parallel=\mathbf x_\perp-\mathbf x_\parallel,
	\]
	so
	\[
	(-H(\mathbf n))\mathbf x=-\mathbf x_\perp+\mathbf x_\parallel.
	\]
	Thus $-H(\mathbf n)$ fixes the axis direction $\mathbf n$ and reverses every vector orthogonal to $\mathbf n$, i.e.\ it is a rotation by angle $\pi$ about axis $\mathbf n$.
	\par\noindent\textbullet\quad Let $\alpha=\theta/2$, so
	\[
	\mathbf n_-=(\cos\alpha,-\sin\alpha,0),\qquad \mathbf n_+=(\cos\alpha,\sin\alpha,0).
	\]
	Using $H(\mathbf u)=I-2\mathbf u\mathbf u^{\mathsf T}$,
	\[
	H(\mathbf i)=\mathrm{diag}(-1,1,1),
	\]
	and
	\[
	H(\mathbf n_-)=
	\begin{pmatrix}
		1-2\cos^2\alpha & 2\cos\alpha\sin\alpha & 0\\
		2\cos\alpha\sin\alpha & 1-2\sin^2\alpha & 0\\
		0&0&1
	\end{pmatrix}
	=
	\begin{pmatrix}
		-\cos(2\alpha) & \sin(2\alpha) & 0\\
		\sin(2\alpha) & \cos(2\alpha) & 0\\
		0&0&1
	\end{pmatrix}.
	\]
	Therefore
	\[
	H(\mathbf i)H(\mathbf n_-)=
	\begin{pmatrix}
		\cos(2\alpha) & -\sin(2\alpha) & 0\\
		\sin(2\alpha) & \cos(2\alpha) & 0\\
		0&0&1
	\end{pmatrix}
	=
	\begin{pmatrix}
		\cos\theta & -\sin\theta & 0\\
		\sin\theta & \cos\theta & 0\\
		0&0&1
	\end{pmatrix}
	=R(\mathbf k,\theta).
	\]
	Similarly one checks $H(\mathbf n_+)H(\mathbf i)=R(\mathbf k,\theta)$.
	Geometrically, $H(\mathbf i)$ and $H(\mathbf n_-)$ are reflections in two planes through the origin intersecting in the $\mathbf k$-axis; the angle between the planes is $\alpha=\theta/2$, hence their composition is a rotation about $\mathbf k$ by angle $2\alpha=\theta$.
\end{solution}

\begin{problem}[CP-I-0086]
\label{prob:cp-i-0086}
\textbf{Simultaneous Diagonalization of Commuting Diagonalizable Operators}

\textbf{Canonical authored completion of the retained invariant-eigenspace step.}

\hspace*{1.5em}Let \(A,B\in\operatorname{End}(V)\) over \(\mathbb C\). Assume that \(A\) and
\(B\) are diagonalizable and commute:
\[
AB=BA.
\]
Prove that \(A\) and \(B\) are simultaneously diagonalizable.
\end{problem}

\begin{solution}
Because \(A\) is diagonalizable,
\[
V=\bigoplus_\lambda E_\lambda(A).
\]
If \(x\in E_\lambda(A)\), then
\[
A(Bx)=B(Ax)=B(\lambda x)=\lambda Bx,
\]
so each \(E_\lambda(A)\) is \(B\)-invariant.

Since \(B\) is diagonalizable, its minimal polynomial is square-free. The minimal polynomial of
the restriction \(B|_{E_\lambda(A)}\) divides the minimal polynomial of \(B\), hence is also
square-free. Therefore each restriction \(B|_{E_\lambda(A)}\) is diagonalizable.

Choose a basis of \(B\)-eigenvectors in every \(E_\lambda(A)\). Their union is a basis of \(V\).
Every vector in this basis is simultaneously an \(A\)-eigenvector and a \(B\)-eigenvector, so
both matrices are diagonal in the same basis.
\end{solution}

\begin{problem}[CP-I-0088]
\label{prob:cp-i-0088}
\textbf{Congruence to the identity over $\mathbb R,\mathbb C,\mathbb Q$}

\textbf{Exact statement.}
\hspace*{1.5em}The square matrices $A$ and $B$ over the field $F$ are congruent if $B=P^{T}AP$ for some invertible matrix $P$ over $F$.
Which of the following symmetric matrices are congruent to the identity matrix over $\mathbb R$, and which over $\mathbb C$?
(Which, if any, over $\mathbb Q$?) Try to get away with the minimum calculation:
\[
A_1=\begin{pmatrix}2&0\\0&3\end{pmatrix},\quad
A_2=\begin{pmatrix}0&2\\2&0\end{pmatrix},\quad
A_3=\begin{pmatrix}-1&0\\0&-1\end{pmatrix},\quad
A_4=\begin{pmatrix}4&4\\4&5\end{pmatrix}.
\]

\end{problem}

\begin{solution}
\emph{Over $\mathbb R$.} By Sylvester's law of inertia, a real symmetric matrix is congruent to $I_2$ iff it is positive definite.
Thus:
\[
A_1 \text{ is positive definite} \Rightarrow A_1\sim_{\mathbb R} I_2.
\]
\[
\det(A_2)=-4<0 \Rightarrow A_2 \text{ is indefinite} \Rightarrow A_2\not\sim_{\mathbb R} I_2.
\]
\[
A_3=-I_2 \text{ is negative definite} \Rightarrow A_3\not\sim_{\mathbb R} I_2.
\]
\[
\text{For }A_4:\ \Delta_1=4>0,\ \det(A_4)=20-16=4>0 \Rightarrow A_4 \text{ positive definite} \Rightarrow A_4\sim_{\mathbb R} I_2.
\]

\emph{Over $\mathbb C$.} Every nonsingular symmetric matrix is congruent to $I_2$: diagonalize to $\mathrm{diag}(\lambda_1,\lambda_2)$ and rescale by square roots in $\mathbb C$.
All four have nonzero determinant, hence
\[
A_1\sim_{\mathbb C} I_2,\quad A_2\sim_{\mathbb C} I_2,\quad A_3\sim_{\mathbb C} I_2,\quad A_4\sim_{\mathbb C} I_2.
\]

\emph{Over $\mathbb Q$.} A necessary condition for $A\sim_{\mathbb Q} I_2$ is that $\det(A)$ be a square in $\mathbb Q^*$, since
\[
1=\det(I_2)=\det(P^TAP)=\det(P)^2\det(A).
\]
Here $\det(A_1)=6$ (not a square), $\det(A_2)=-4$ (not a square in $\mathbb Q^*$), $\det(A_3)=1$ (a square), $\det(A_4)=4$ (a square).
So only $A_3,A_4$ can possibly be congruent to $I_2$ over $\mathbb Q$.

For $A_3=-I_2$, $-I_2\sim_{\mathbb Q} I_2$ would imply $P^T(-I_2)P=I_2$, i.e.\ $P^TP=-I_2$.
Taking the first column $(a,b)^T$ gives $a^2+b^2=-1$, impossible in $\mathbb Q$.
Hence $A_3\not\sim_{\mathbb Q} I_2$.

For $A_4$, the quadratic form is $q(x,y)=4x^2+8xy+5y^2$ and
\[
q(x,y)=4(x+y)^2+y^2.
\]
With the rational change of variables $u=x+y$, $v=y$ we get $q=4u^2+v^2$, i.e.\ $A_4\sim_{\mathbb Q}\mathrm{diag}(4,1)$.
Then scaling $u\mapsto \tfrac12 u$ (rational) gives $\mathrm{diag}(1,1)=I_2$, so $A_4\sim_{\mathbb Q} I_2$.

\medskip
\noindent
\textbf{Conclusion.}
Over $\mathbb R$: $A_1,A_4$ only.
Over $\mathbb C$: all four.
Over $\mathbb Q$: $A_4$ only.
\end{solution}

\begin{problem}[CP-I-0089]
\label{prob:cp-i-0089}
\textbf{Similarity to a Unipotent Jordan Block and Fast Powers}

\textbf{Canonical authored completion of the retained similarity fragment.}

\hspace*{1.5em}Suppose
\[
A=PBP^{-1},\qquad B=I+N,\qquad N^2=0.
\]
Prove that for every integer \(m\ge0\),
\[
B^m=I+mN,
\qquad
A^m=I+m(A-I).
\]
Deduce that
\[
\det(A^m)=(\det A)^m.
\]
\end{problem}

\begin{solution}
Since \(N^2=0\), the binomial theorem truncates:
\[
B^m=(I+N)^m=I+mN.
\]
As \(B-I=N\),
\[
B^m-I=m(B-I).
\]
Similarity gives
\[
A^m=PB^mP^{-1}
=P(I+mN)P^{-1}
=I+mPNP^{-1}.
\]
But
\[
PNP^{-1}=P(B-I)P^{-1}=A-I,
\]
hence
\[
A^m=I+m(A-I).
\]

Finally, multiplicativity of determinant gives
\[
\det(A^m)=\underbrace{\det(A)\cdots\det(A)}_{m\text{ factors}}
=(\det A)^m.
\]
\end{solution}

\begin{problem}[CP-I-0091]
\label{prob:cp-i-0091}
\textbf{Change of basis to a Jordan block; compute $A^n$}

\textbf{Exact statement.}
\hspace*{1.5em}The linear map $S:\mathbb R^2\to\mathbb R^2$ is defined in terms of its matrix $A$ with respect to the standard basis by
\[
S\binom{x}{y}=A\binom{x}{y}=\binom{-5x+9y}{-4x+7y}.
\]
Find the matrix $B$ for $S$ with respect to the basis
\[
\left\{\binom{3}{2},\binom{1}{1}\right\}.
\]
Show that $B^n-I = n(B-I)$ for all positive integers $n$, and hence determine $A^n$.
Verify that $\det(A^n)=(\det A)^n$.

\end{problem}

\begin{solution}
From the definition,
\[
A=\begin{pmatrix}-5&9\\-4&7\end{pmatrix}.
\]
Let $v_1=\binom{3}{2}$ and $v_2=\binom{1}{1}$ and set
\[
P=[v_1\ v_2]=\begin{pmatrix}3&1\\2&1\end{pmatrix}.
\]
Compute
\[
Av_1=\binom{-15+18}{-12+14}=\binom{3}{2}=v_1,
\qquad
Av_2=\binom{-5+9}{-4+7}=\binom{4}{3}.
\]
Write $Av_2=a v_1+b v_2$:
\[
\binom{4}{3}=a\binom{3}{2}+b\binom{1}{1}
\ \Longrightarrow\
\begin{cases}
	3a+b=4\\
	2a+b=3
\end{cases}
\ \Longrightarrow\ a=1,\ b=1.
\]
Hence the columns of $B$ are $(1,0)^T$ and $(1,1)^T$, so
\[
\boxed{
	B=\begin{pmatrix}1&1\\0&1\end{pmatrix}.}
\]
Now $B=I+N$ with $N=\begin{psmallmatrix}0&1\\0&0\end{psmallmatrix}$ and $N^2=0$, so
\[
B^n=(I+N)^n=I+nN=\begin{pmatrix}1&n\\0&1\end{pmatrix}.
\]
Therefore
\[
\boxed{B^n-I=n(B-I).}
\]
Since $A=PBP^{-1}$, we have $A^n=PB^nP^{-1}$. Equivalently (by similarity),
\[
A^n=I+n(A-I)=nA+(1-n)I.
\]
Thus
\[
\boxed{
	A^n
	=
	n\begin{pmatrix}-5&9\\-4&7\end{pmatrix}+(1-n)\begin{pmatrix}1&0\\0&1\end{pmatrix}
	=
	\begin{pmatrix}
		1-6n & 9n\\
		-4n & 1+6n
	\end{pmatrix}.}
\]
Finally,
\[
\det A = (-5)(7)-9(-4)=1
\quad\Rightarrow\quad
\det(A^n)=1=(\det A)^n.
\]
\end{solution}

\begin{problem}[CP-I-0092]
\label{prob:cp-i-0092}
\textbf{Basic identities for Hermitian matrices and unitary similarity}

\textbf{Exact statement.}
\hspace*{1.5em}Let \(A\) and \(B\) be \(n\times n\) Hermitian matrices and \(U\) be an \(n\times n\) unitary matrix. Show that
(i) \(\operatorname{tr}A\) and \(\det A\) are real;
(ii) \(AB+BA\) is Hermitian;
(iii) \(i(AB-BA)\) is Hermitian and has zero trace;
(iv) \(\operatorname{tr}(AB)\) and \(\det(AB)\) are real;
(v) \(\operatorname{tr}(UAU^\ast)=\operatorname{tr}(A)\) and \(\det(UAU^\ast)=\det(A)\).

\end{problem}

\begin{solution}
Recall: Hermitian means \(A^\ast=A\), and unitary means \(U^\ast U=I\). Also
\(\overline{\operatorname{tr}(X)}=\operatorname{tr}(X^\ast)\),
\(\overline{\det(X)}=\det(X^\ast)\),
and \(\operatorname{tr}(XY)=\operatorname{tr}(YX)\).

	\par\noindent (i)\quad
	Since \(A^\ast=A\),
	\[
	\overline{\operatorname{tr}(A)}=\operatorname{tr}(A^\ast)=\operatorname{tr}(A),
	\]
	so \(\operatorname{tr}(A)\in\mathbb R\). Also
	\[
	\overline{\det(A)}=\det(A^\ast)=\det(A),
	\]
	so \(\det(A)\in\mathbb R\).

	\par\noindent (ii)\quad
	\[
	(AB+BA)^\ast=B^\ast A^\ast + A^\ast B^\ast=BA+AB=AB+BA.
	\]

	\par\noindent (iii)\quad
	\[
	(AB-BA)^\ast=B^\ast A^\ast-A^\ast B^\ast=BA-AB=-(AB-BA),
	\]
	so \(AB-BA\) is skew-Hermitian, hence
	\[
	(i(AB-BA))^\ast=-i(AB-BA)^\ast=i(AB-BA),
	\]
	so \(i(AB-BA)\) is Hermitian. Moreover,
	\[
	\operatorname{tr}(AB-BA)=\operatorname{tr}(AB)-\operatorname{tr}(BA)=0.
	\]

	\par\noindent (iv)\quad
	\[
	\overline{\operatorname{tr}(AB)}=\operatorname{tr}((AB)^\ast)=\operatorname{tr}(B^\ast A^\ast)=\operatorname{tr}(BA)=\operatorname{tr}(AB),
	\]
	so \(\operatorname{tr}(AB)\in\mathbb R\). Also \(\det(AB)=\det(A)\det(B)\), and by (i) both determinants are real, hence \(\det(AB)\in\mathbb R\).

	\par\noindent (v)\quad
	Using cyclicity,
	\[
	\operatorname{tr}(UAU^\ast)=\operatorname{tr}(AU^\ast U)=\operatorname{tr}(A).
	\]
	Also
	\[
	\det(UAU^\ast)=\det(U)\det(A)\det(U^\ast)=\det(U)\det(A)\overline{\det(U)}.
	\]
	For unitary \(U\), \(|\det(U)|=1\), so \(\det(U)\overline{\det(U)}=1\), hence \(\det(UAU^\ast)=\det(A)\).
\end{solution}

\begin{problem}[CP-I-0095]
\label{prob:cp-i-0095}
\textbf{Char.\ and minimal polynomials do not determine JNF in $4\times 4$}

\textbf{Exact statement.}
\hspace*{1.5em}Recall that the Jordan normal form of a $3\times 3$ complex matrix can be deduced from its characteristic and minimal polynomials.
Give an example to show that this is not so for $4\times 4$ complex matrices.

\end{problem}

\begin{solution}
In size $4$, the characteristic polynomial fixes the total algebraic multiplicity of each eigenvalue,
and the minimal polynomial fixes the size of the largest Jordan block for each eigenvalue, but this does not always determine the block partition.

Consider
\[
A=\mathrm{diag}(J_2(0),J_2(0))
=\begin{pmatrix}
	0&1&0&0\\
	0&0&0&0\\
	0&0&0&1\\
	0&0&0&0
\end{pmatrix},
\qquad
B=\mathrm{diag}(J_2(0),0,0)
=\begin{pmatrix}
	0&1&0&0\\
	0&0&0&0\\
	0&0&0&0\\
	0&0&0&0
\end{pmatrix}.
\]
Both satisfy $A^2=B^2=0$ but $A\neq 0$ and $B\neq 0$, so both have minimal polynomial $m(t)=t^2$.
Also both are nilpotent $4\times 4$, hence $\chi_A(t)=\chi_B(t)=t^4$.
However $\operatorname{rank}(A)=2$ while $\operatorname{rank}(B)=1$, and rank is a similarity invariant, so $A$ and $B$ are not similar.
Thus $(\chi,m)$ does not determine Jordan form in dimension $4$.
\end{solution}

\begin{problem}[CP-I-0096]
\label{prob:cp-i-0096}
\textbf{Half-Space Electrostatics by the Method of Images}

\textbf{Canonical authored reconstruction of the source's half-space electrostatics example.}

\hspace*{1.5em}Let
\[
H=\{x=(x_1,x_2,x_3)\in\mathbb R^3:x_3>0\},
\]
and let \(\xi=(\xi_1,\xi_2,\xi_3)\in H\). Reflect \(\xi\) across the grounded plane \(x_3=0\):
\[
\xi^\ast=(\xi_1,\xi_2,-\xi_3).
\]
Define
\[
G(x,\xi)
=
\frac{1}{4\pi|x-\xi|}
-
\frac{1}{4\pi|x-\xi^\ast|}.
\]
Show that \(G(\cdot,\xi)=0\) on \(x_3=0\), that
\[
-\Delta_x G(\cdot,\xi)=\delta_\xi
\]
in \(H\), and explain the interpretation in terms of a unit charge at \(\xi\) and an image charge
of opposite sign at \(\xi^\ast\).
\end{problem}

\begin{solution}
If \(x_3=0\), reflection symmetry gives
\[
|x-\xi|=|x-\xi^\ast|.
\]
Hence
\[
G(x,\xi)=0
\]
on the boundary plane.

In \(\mathbb R^3\),
\[
\Phi(x,\eta)=\frac{1}{4\pi|x-\eta|}
\]
is the fundamental solution of \(-\Delta\):
\[
-\Delta_x\Phi(x,\eta)=\delta_\eta.
\]
Therefore
\[
-\Delta_xG(\cdot,\xi)
=
\delta_\xi-\delta_{\xi^\ast}
\]
in \(\mathbb R^3\). Since \(\xi^\ast\notin H\), restriction to the upper half-space gives
\[
-\Delta_xG(\cdot,\xi)=\delta_\xi.
\]

Thus \(G\) is the Dirichlet Green function for the upper half-space. Physically, the grounded
plane is enforced by replacing the boundary-value problem with the real unit charge at \(\xi\)
and a fictitious image charge of magnitude \(-1\) at the reflected point \(\xi^\ast\).
\end{solution}

\begin{problem}[CP-I-0098]
\label{prob:cp-i-0098}
\textbf{Cayley transform}

\hspace*{1.5em}If $K^\dagger=-K$ then $I-K$ is invertible, and
\[
U=(I+K)(I-K)^{-1}
\]
is unitary. Eigenvalues map by $it\mapsto \frac{1+it}{1-it}$, sending the imaginary axis to the unit circle.

\renewcommand{\arraystretch}{1.25}
\noindent\resizebox{\linewidth}{!}{%
\begin{tabular}{|c|p{4.2cm}|p{6.6cm}|p{4.2cm}|}
	\hline
	\textbf{No.} & \textbf{Core idea / theorem} & \textbf{Key formula(s)} & \textbf{What to check fast} \\
	\hline
	1 &
	Triangular determinant \& spectrum &
	If $A$ triangular, $p_A(t)=\det(A-tI)=\prod_{i=1}^n (A_{ii}-t)$ &
	Eigenvalues are diagonal entries (with multiplicity). \\
	\hline
	2 &
	Diagonalizable $\Leftrightarrow$ enough eigenvectors (Jordan obstruction) &
	$m_{\mathrm{geo}}(\lambda)=\dim\ker(A-\lambda I)\le m_{\mathrm{alg}}(\lambda)$; diag iff equality for all $\lambda$ &
	If one eigenvalue $\lambda$ of mult.\ $n$, compute $\dim\ker(A-\lambda I)$; if $<n$, not diagonalizable. \\
	\hline
	3 &
	Orthogonal matrices and completing an ONB &
	$Q^TQ=I\iff$ columns orthonormal; in $\mathbb R^3$: $v_3=\pm(v_1\times v_2)$ &
	Two completions only ($\det=\pm1$). \\
	\hline
	4 &
	Spectral theorem (real symmetric) &
	$A=A^T\Rightarrow P^TAP=D$, $P$ orthogonal; $A^{-1}=PD^{-1}P^T$ if $\lambda_i\neq 0$ &
	Eigenvalues real; eigenvectors for distinct eigenvalues orthogonal. \\
	\hline
	5 &
	Quadratic forms: principal axes via diagonalization &
	$q(x)=x^TAx$ (take $A$ symmetric); $x=Py\Rightarrow q=y^T(P^TAP)y$ &
	Remove cross terms by orthogonal change; semi-axes $1/\sqrt{\lambda_i}$ (positive $\lambda_i$). \\
	\hline
	6 &
	Anti-Hermitian / unitary spectra and orthogonality &
	$A^\dagger=-A\Rightarrow \lambda\in i\mathbb R$;
	$U^\dagger U=I\Rightarrow |\lambda|=1$ &
	Distinct eigenvalues $\Rightarrow$ eigenvectors orthogonal (Hermitian inner product). \\
	\hline
	7 &
	Cayley--Hamilton (esp.\ $2\times2$) and power recurrences &
	$A^2-(\operatorname{tr} A)A+(\det A)I=0$;
	then $A^n=\alpha_n A+\beta_n I$ with recurrences &
	Use CH to reduce all powers to span$\{I,A\}$; repeated root $\Rightarrow$ Jordan/unipotent behavior. \\
	\hline
	8 &
	Matrix of a linear map in general bases; change-of-basis rule &
	$[T(v)]_F=A[v]_E$;
	under $E\to E'$, $F\to F'$:
	$A_{F',E'}=C_F^{-1}A_{F,E}C_E$ &
	Columns are coordinates of $T(e_j)$ in codomain basis. \\
	\hline
	9 &
	Similarity to a Jordan block; fast powers &
	If $A=PBP^{-1}$ then $A^n=PB^nP^{-1}$;
	if $B=I+N$, $N^2=0$, then $B^n=I+nN$ &
	If $(B-I)^2=0$, then $B^n-I=n(B-I)$; transfer by similarity. \\
	\hline
	10 &
	Quadrics from eigenvalues (classification) &
	For symmetric $A$: $x^TAx=1 \iff \sum \lambda_i u_i^2=1$ in ON eigenbasis &
	All $\lambda_i>0$: ellipsoid; one $\lambda_i=0$: cylinder; mixed signs: hyperboloid (signature test). \\
	\hline
	11 &
	Min distance to origin on a quadric &
	Diagonalize and minimize $u_1^2+u_2^2+u_3^2$ subject to $\sum \lambda_i u_i^2=1$ &
	If $A\succ0$: $r_{\min}=1/\sqrt{\lambda_{\max}}$ (achieved on eigenvector of $\lambda_{\max}$). \\
	\hline
	12 &
	Cayley transform: skew-Hermitian $\to$ unitary &
	If $K^\dagger=-K$, then $U=(I+K)(I-K)^{-1}$ is unitary; eigenvalue map:
	$it\mapsto\frac{1+it}{1-it}$ &
	$I-K$ invertible; unit circle image; invariants preserved by similarity/change of basis. \\
	\hline
\end{tabular}}

\clearpage

Let $A\in \mathbb R^{n\times n}$ (or $\mathbb C^{n\times n}$).

\textbf{Eigenvalue/eigenvector.}
A scalar $\lambda$ is an eigenvalue of $A$ if there exists $v\neq 0$ with
\[
Av=\lambda v.
\]
Such a vector $v$ is an eigenvector for $\lambda$.

\textbf{Characteristic polynomial.}
\[
p_A(t)=\det(tI-A).
\]
Then $\lambda$ is an eigenvalue iff $p_A(\lambda)=0$.

\textbf{Eigenspace and multiplicities.}
The eigenspace of $\lambda$ is
\[
E_\lambda=\ker(A-\lambda I).
\]
Algebraic multiplicity $m_{\rm alg}(\lambda)$ is the multiplicity of $\lambda$ as a root of $p_A(t)$.
Geometric multiplicity $m_{\rm geo}(\lambda)=\dim E_\lambda$.

\textbf{Diagonalizability.}
$A$ is diagonalizable iff it has $n$ linearly independent eigenvectors, equivalently
\[
\sum_{\lambda} \dim E_\lambda = n,
\]
equivalently $m_{\rm geo}(\lambda)=m_{\rm alg}(\lambda)$ for every eigenvalue $\lambda$.

\textbf{Standard workflow (practical).}

	\par\noindent\textbullet\quad Compute $p_A(t)=\det(tI-A)$ and factor it.
	\par\noindent\textbullet\quad For each eigenvalue $\lambda$, solve $(A-\lambda I)v=0$ to get eigenvectors (a basis of $E_\lambda$).
	\par\noindent\textbullet\quad Check diagonalizability by counting eigenvectors: total dimension $\sum_\lambda \dim E_\lambda$.
	\par\noindent\textbullet\quad If $A$ is symmetric (real): orthogonality is automatic for distinct eigenvalues; then build an orthonormal eigenbasis and use $P^TAP=D$.
	\par\noindent\textbullet\quad If needed: normalize eigenvectors, and form $P=[u_1\ \cdots\ u_n]$.

\textbf{Real symmetric matrices.}
If $A=A^T$ and $Av=\lambda v$, $Aw=\mu w$ with $\lambda\neq\mu$, then
\[
v^T w=0.
\]
\emph{Proof sketch:} $v^TAw=\mu v^Tw$, but also $v^TAw=(A^Tv)^Tw=(Av)^Tw=\lambda v^Tw$; hence $(\mu-\lambda)v^Tw=0$.

\textbf{Anti-Hermitian matrices.}
If $A^\dagger=-A$, then eigenvalues lie in $i\mathbb R$ and eigenvectors for distinct eigenvalues are orthogonal.

\textbf{Unitary matrices.}
If $U^\dagger U=I$, then $|\lambda|=1$ for eigenvalues and distinct-eigenvalue eigenvectors are orthogonal.

Let
\[
T=\begin{pmatrix}
	2&5&-1\\
	0&3&4\\
	0&0&-2
\end{pmatrix}.
\]

\textbf{Step 1: eigenvalues.}
$T-tI$ is triangular, so
\[
p_T(t)=\det(T-tI)=(2-t)(3-t)(-2-t).
\]
Hence eigenvalues are $\lambda=2,3,-2$.

\textbf{Step 2: eigenvectors.}
Solve $(T-\lambda I)v=0$.

\emph{For $\lambda=2$:}
\[
T-2I=\begin{pmatrix}0&5&-1\\0&1&4\\0&0&-4\end{pmatrix}.
\]
From bottom: $-4v_3=0\Rightarrow v_3=0$.
Then $v_2+4v_3=v_2=0$.
So $v_1$ free. Take $v=(1,0,0)^T$.

\emph{For $\lambda=3$:}
\[
T-3I=\begin{pmatrix}-1&5&-1\\0&0&4\\0&0&-5\end{pmatrix}.
\]
From $4v_3=0\Rightarrow v_3=0$.
Then $-v_1+5v_2=0\Rightarrow v_1=5v_2$.
Take $v=(5,1,0)^T$.

\emph{For $\lambda=-2$:}
\[
T+2I=\begin{pmatrix}4&5&-1\\0&5&4\\0&0&0\end{pmatrix}.
\]
Let $v_3=s$ free. From second row: $5v_2+4s=0\Rightarrow v_2=-\frac45 s$.
From first: $4v_1+5v_2-s=0\Rightarrow 4v_1+5(-\frac45 s)-s=0\Rightarrow 4v_1-5s=0\Rightarrow v_1=\frac54 s$.
Take $s=4$ to clear denominators: $v=(5,-16,4)^T$.

\textbf{Conclusion.} For distinct diagonal entries, triangular matrices are diagonalizable.

Let
\[
A=\begin{pmatrix}
	3&1&1\\
	1&2&0\\
	1&0&2
\end{pmatrix}.
\]

\textbf{Step 1: eigenvalues.}
Compute $p_A(t)=\det(tI-A)$ and factor (or use a CAS; here it factors as)
\[
p_A(t)=(t-1)(t-2)(t-4).
\]
So $\lambda=1,2,4$.

\textbf{Step 2: eigenvectors (solve each system).}

	\par\noindent\textbullet\quad For $\lambda=1$, a solution is $v_1=(-1,1,1)^T$.
	\par\noindent\textbullet\quad For $\lambda=2$, a solution is $v_2=(0,-1,1)^T$.
	\par\noindent\textbullet\quad For $\lambda=4$, a solution is $v_3=(2,1,1)^T$.

\textbf{Step 3: orthogonality and normalization.}
Because $A$ is symmetric, these eigenvectors are automatically orthogonal (check quickly):
\[
v_1\cdot v_2=0,\quad v_1\cdot v_3=0,\quad v_2\cdot v_3=0.
\]
Normalize:
\[
u_1=\frac{1}{\sqrt3}(-1,1,1),\quad
u_2=\frac{1}{\sqrt2}(0,-1,1),\quad
u_3=\frac{1}{\sqrt6}(2,1,1).
\]

\textbf{Step 4: diagonalization and inverse.}
Let $P=[u_1\ u_2\ u_3]$, $D=\mathrm{diag}(1,2,4)$. Then
\[
P^TAP=D,\qquad A^{-1}=P D^{-1} P^T.
\]

Let
\[
J=\begin{pmatrix}1&1\\0&1\end{pmatrix}.
\]

\textbf{Step 1: eigenvalues.}
\[
p_J(t)=\det(tI-J)=\det\begin{pmatrix}t-1&-1\\0&t-1\end{pmatrix}=(t-1)^2.
\]
Only eigenvalue is $\lambda=1$ with $m_{\rm alg}=2$.

\textbf{Step 2: eigenspace.}
\[
J-I=\begin{pmatrix}0&1\\0&0\end{pmatrix}
\Rightarrow (J-I)v=0 \Rightarrow v_2=0.
\]
So $E_1=\{(v_1,0)\}$ has dimension $1$.

\textbf{Conclusion.}
$m_{\rm geo}(1)=1<2=m_{\rm alg}(1)$, so $J$ is \textbf{not} diagonalizable.

\textbf{Bonus: fast powers.}
Write $J=I+N$ with $N=\begin{psmallmatrix}0&1\\0&0\end{psmallmatrix}$ and $N^2=0$. Then
\[
J^n=(I+N)^n=I+nN=\begin{pmatrix}1&n\\0&1\end{pmatrix}.
\]

\textbf{Anti-Hermitian example.}
Let
\[
K=\begin{pmatrix}0&i\\ i&0\end{pmatrix}.
\]
Then $K^\dagger=-K$, so eigenvalues are purely imaginary. Indeed,
\[
p_K(t)=\det(tI-K)=\det\begin{pmatrix}t&-i\\-i&t\end{pmatrix}=t^2+1,
\]
so $\lambda=\pm i$. Distinct eigenvalues $\Rightarrow$ eigenvectors orthogonal.

\textbf{Unitary example.}
Let
\[
U=\begin{pmatrix}0&i\\-1&0\end{pmatrix}.
\]
One checks $U^\dagger U=I$. Its eigenvalues satisfy $|\lambda|=1$; indeed,
\[
\det(U-\lambda I)=\lambda^2+i=0 \Rightarrow \lambda=\pm\frac{1-i}{\sqrt2},
\]
which lie on the unit circle.

Given a symmetric quadratic form $q(x)=x^TAx$, the surface $q(x)=1$ becomes
\[
\lambda_1 u_1^2+\lambda_2 u_2^2+\lambda_3 u_3^2=1
\]
in an orthonormal eigenbasis. The signs of $\lambda_i$ classify the quadric:

	\par\noindent\textbullet\quad all $\lambda_i>0$: ellipsoid,
	\par\noindent\textbullet\quad one $\lambda_i=0$ and others $>0$: cylinder,
	\par\noindent\textbullet\quad mixed signs: hyperboloid.

If $\lambda_i>0$, the semi-axis length along eigenvector $u_i$ is $1/\sqrt{\lambda_i}$.

\textbf{Distance minimization trick.}
For $x^TAx=1$ with $A\succ0$ (positive definite),
the nearest point to the origin lies on the eigenvector for $\lambda_{\max}$, and
\[
r_{\min}=\frac{1}{\sqrt{\lambda_{\max}}}.
\]
This comes from minimizing $\|x\|^2$ under the constraint $\sum \lambda_i u_i^2=1$.
\end{problem}

\begin{problem}[CP-I-0101]
\label{prob:cp-i-0101}
\textbf{Diagonalising a rotated quadratic form}

\textbf{Exact statement.}
\hspace*{1.5em}Diagonalise the quadratic form in $\mathbb R^2$ defined by
\[
F(x,y)=\bigl(a\cos^2\theta+b\sin^2\theta\bigr)x^2
+2(a-b)(\sin\theta\cos\theta)\,xy
+\bigl(a\sin^2\theta+b\cos^2\theta\bigr)y^2,
\]
i.e.\ find its eigenvalues and principal axes ($a,b,\theta$ are constants).

\end{problem}

\begin{solution}
Write $F(x,y)=\begin{pmatrix}x&y\end{pmatrix}Q\begin{pmatrix}x\\y\end{pmatrix}$ with
\[
Q=
\begin{pmatrix}
	a\cos^2\theta+b\sin^2\theta & (a-b)\sin\theta\cos\theta\\
	(a-b)\sin\theta\cos\theta & a\sin^2\theta+b\cos^2\theta
\end{pmatrix}.
\]
Let
\[
R=
\begin{pmatrix}
	\cos\theta&-\sin\theta\\
	\sin\theta&\cos\theta
\end{pmatrix}.
\]
Then one checks
\[
Q=R\begin{pmatrix}a&0\\0&b\end{pmatrix}R^T.
\]
Hence the eigenvalues are $a$ and $b$. The principal axes are the rotated directions
\[
(\cos\theta,\sin\theta)\quad\text{and}\quad(-\sin\theta,\cos\theta).
\]
Equivalently, in rotated coordinates
\[
\begin{pmatrix}u\\v\end{pmatrix}=R^T\begin{pmatrix}x\\y\end{pmatrix}
=
\begin{pmatrix}
	x\cos\theta+y\sin\theta\\
	-x\sin\theta+y\cos\theta
\end{pmatrix},
\]
the quadratic form becomes
\[
F(x,y)=a\,u^2+b\,v^2.
\]
\end{solution}

\begin{problem}[CP-I-0102]
\label{prob:cp-i-0102}
\textbf{Extracting Axis and Angle from a Rotation Matrix}

\textbf{Canonical authored completion of the retained axis--angle fragment.}

\hspace*{1.5em}Let \(R\in SO(3)\) be a nontrivial rotation with axis a unit vector \(n\) and
rotation angle \(\theta\). Starting from Rodrigues' formula
\[
R=\cos\theta\,I+(1-\cos\theta)nn^T-\sin\theta\,[n]_\times,
\]
prove that
\[
\operatorname{tr}(R)=1+2\cos\theta
\]
and
\[
R-R^T=-2\sin\theta\,[n]_\times.
\]
Deduce, when \(\sin\theta\neq0\), the extraction formulas
\[
\cos\theta=\frac{\operatorname{tr}(R)-1}{2},
\]
and
\[
(R_{32}-R_{23},\,R_{13}-R_{31},\,R_{21}-R_{12})
=
-2\sin\theta\,n.
\]
\end{problem}

\begin{solution}
Taking the trace of Rodrigues' formula,
\[
\operatorname{tr}(R)
=
3\cos\theta
+(1-\cos\theta)\operatorname{tr}(nn^T)
-\sin\theta\,\operatorname{tr}([n]_\times).
\]
Since \(n\) is a unit vector,
\[
\operatorname{tr}(nn^T)=n^Tn=1,
\]
and every skew-symmetric matrix has zero trace. Hence
\[
\operatorname{tr}(R)
=
3\cos\theta+1-\cos\theta
=
1+2\cos\theta.
\]

Now transpose Rodrigues' formula. Because \(nn^T\) is symmetric and
\[
[n]_\times^T=-[n]_\times,
\]
we get
\[
R^T
=
\cos\theta\,I
+(1-\cos\theta)nn^T
+\sin\theta\,[n]_\times.
\]
Subtracting,
\[
R-R^T=-2\sin\theta\,[n]_\times.
\]
Reading off the three independent skew entries gives
\[
(R_{32}-R_{23},R_{13}-R_{31},R_{21}-R_{12})
=
-2\sin\theta\,n.
\]
Thus the trace determines \(\cos\theta\), and when \(\sin\theta\neq0\) the antisymmetric part
determines the axis direction.
\end{solution}

\begin{problem}[CP-I-0103]
\label{prob:cp-i-0103}
\textbf{Vandermonde determinant; evaluation functionals; Lagrange dual basis}

\textbf{Exact statement.}

	\par\hspace*{1.5em} (a)\quad Let \(a_0,\dots,a_n\) be distinct real numbers and let
	\[
	A=\begin{pmatrix}
		1&1&\cdots&1\\
		a_0&a_1&\cdots&a_n\\
		a_0^2&a_1^2&\cdots&a_n^2\\
		\vdots&\vdots&\ddots&\vdots\\
		a_0^n&a_1^n&\cdots&a_n^n
	\end{pmatrix}.
	\]
	Show that \(\det(A)\neq 0\).
	\par\noindent (b)\quad Let \(P_n\) be the space of real polynomials of degree \(\le n\). For \(x\in\mathbb R\) define \(e_x\in P_n^*\) by
	\(e_x(p)=p(x)\). Use (a) to show \(\{e_{a_0},\dots,e_{a_n}\}\) is linearly independent, hence a basis of \(P_n^*\).
	\par\noindent (c)\quad Identify the basis of \(P_n\) to which \((e_{a_0},\dots,e_{a_n})\) is dual.

\end{problem}

\begin{solution}
(a) If \(Ac=0\) for \(c=(c_0,\dots,c_n)^T\), define \(p(t)=\sum_{k=0}^n c_k t^k\). Then \((Ac)_i=p(a_i)\), so \(p(a_i)=0\) for
\(i=0,\dots,n\). A nonzero polynomial of degree \(\le n\) cannot have \(n+1\) distinct roots, hence \(p=0\) and \(c=0\).
Thus \(\ker A=\{0\}\), so \(\det(A)\neq0\).

(b) If \(\sum_{i=0}^n \lambda_i e_{a_i}=0\), evaluate on \(t^k\) for \(k=0,\dots,n\):
\[
0=\sum_{i=0}^n \lambda_i e_{a_i}(t^k)=\sum_{i=0}^n \lambda_i a_i^k,
\]
i.e. \(A\lambda=0\). By (a), \(\lambda=0\). Hence the \(e_{a_i}\) are independent, and since \(\dim P_n^*=n+1\), they form a basis.

(c) The dual basis is the Lagrange basis \(\{\ell_0,\dots,\ell_n\}\) with \(\ell_i(a_j)=\delta_{ij}\), namely
\[
\ell_i(t)=\prod_{\substack{0\le j\le n\\ j\ne i}}\frac{t-a_j}{a_i-a_j}.
\]
\end{solution}

\begin{problem}[CP-I-0107]
\label{prob:cp-i-0107}
\textbf{Characteristic Polynomials of Block-Triangular Matrices}

\textbf{Canonical authored differentiation of the retained triangular-spectrum fragment.}

\hspace*{1.5em}Let
\[
A=
\begin{pmatrix}
B&C\\
0&D
\end{pmatrix}
\]
be block upper triangular, where \(B\) and \(D\) are square. Prove that
\[
\chi_A(t)=\chi_B(t)\chi_D(t).
\]
Deduce that the eigenvalues of \(A\), counted with algebraic multiplicity, are the eigenvalues of
\(B\) together with those of \(D\). Recover the usual triangular-matrix result as a special case.
\end{problem}

\begin{solution}
We have
\[
A-tI=
\begin{pmatrix}
B-tI&C\\
0&D-tI
\end{pmatrix}.
\]
The determinant of a block triangular matrix is the product of the determinants of its diagonal
blocks. Therefore
\[
\chi_A(t)=\det(A-tI)
=\det(B-tI)\det(D-tI)
=\chi_B(t)\chi_D(t).
\]
Thus the roots of \(\chi_A\) are exactly the roots of \(\chi_B\) and \(\chi_D\), with multiplicities
added when the same root occurs in both.

An ordinary upper triangular matrix is obtained by repeatedly taking \(1\times1\) diagonal
blocks. Hence
\[
\chi_A(t)=\prod_i(a_{ii}-t).
\]
\end{solution}

\begin{problem}[CP-I-0108]
\label{prob:cp-i-0108}
\textbf{Change-of-basis matrix of a linear map}

\textbf{Exact statement.}
\hspace*{1.5em}Let $\alpha:\mathbb R^3\to\mathbb R^3$ be given by
\[
\alpha\begin{pmatrix}x_1\\x_2\\x_3\end{pmatrix}
=
\begin{pmatrix}
	2&1&0\\
	0&2&1\\
	0&0&2
\end{pmatrix}
\begin{pmatrix}x_1\\x_2\\x_3\end{pmatrix}.
\]
Find the matrix representing $\alpha$ relative to the basis
\[
b_1=\begin{pmatrix}1\\1\\1\end{pmatrix},\quad
b_2=\begin{pmatrix}1\\1\\0\end{pmatrix},\quad
b_3=\begin{pmatrix}1\\0\\0\end{pmatrix}
\]
for both the domain and the range.
Write down bases for the domain and range with respect to which the matrix of $\alpha$ is the identity.
\end{problem}

\begin{solution}
Compute
\[
\alpha(b_1)=\begin{pmatrix}3\\3\\2\end{pmatrix},\qquad
\alpha(b_2)=\begin{pmatrix}3\\2\\0\end{pmatrix},\qquad
\alpha(b_3)=\begin{pmatrix}2\\0\\0\end{pmatrix}.
\]
Express each in the basis $B=(b_1,b_2,b_3)$.

For $\alpha(b_1)$, since only $b_1$ has third coordinate $1$, we must take $2b_1=(2,2,2)^T$, leaving $(1,1,0)^T=b_2$:
\[
\alpha(b_1)=2b_1+b_2 \ \Rightarrow\ [\alpha(b_1)]_B=\begin{pmatrix}2\\1\\0\end{pmatrix}.
\]
For $\alpha(b_2)$, we have
\[
\alpha(b_2)=2b_2+b_3 \ \Rightarrow\ [\alpha(b_2)]_B=\begin{pmatrix}0\\2\\1\end{pmatrix}.
\]
For $\alpha(b_3)$,
\[
\alpha(b_3)=2b_3 \ \Rightarrow\ [\alpha(b_3)]_B=\begin{pmatrix}0\\0\\2\end{pmatrix}.
\]
Therefore
\[
[\alpha]_B=
\begin{pmatrix}
	2&0&0\\
	1&2&0\\
	0&1&2
\end{pmatrix}.
\]

To make the matrix the identity, take the domain basis $B$ and the range basis
\[
\alpha(B)=(\alpha(b_1),\alpha(b_2),\alpha(b_3))
=\left(
\begin{pmatrix}3\\3\\2\end{pmatrix},
\begin{pmatrix}3\\2\\0\end{pmatrix},
\begin{pmatrix}2\\0\\0\end{pmatrix}
\right).
\]
Then $[\alpha]_{B}^{\alpha(B)}=I_3$.

\bigskip
\end{solution}

\begin{problem}[CP-I-0109]
\label{prob:cp-i-0109}
\textbf{Determinant of a circulant matrix}

\textbf{Exact statement.}
\hspace*{1.5em}Let $f(x)=a_0+a_1x+\cdots +a_n x^n$ with $a_i\in\mathbb C$, and let $C$ be the $(n+1)\times(n+1)$ circulant matrix
\[
C=\begin{pmatrix}
	a_0&a_1&a_2&\cdots&a_n\\
	a_n&a_0&a_1&\cdots&a_{n-1}\\
	a_{n-1}&a_n&a_0&\cdots&a_{n-2}\\
	\vdots&\vdots&\vdots&\ddots&\vdots\\
	a_1&a_2&a_3&\cdots&a_0
\end{pmatrix}.
\]
Show that $\det(C)=\prod_{j=0}^{n} f(\zeta^j)$ where $\zeta=\exp(2\pi i/(n+1))$.
\end{problem}

\begin{solution}
For $j=0,1,\dots,n$, define
\[
v_j=(1,\zeta^j,\zeta^{2j},\dots,\zeta^{nj})^T.
\]
A direct computation using the circulant structure gives
\[
Cv_j=f(\zeta^j)\,v_j,
\]
so $v_j$ is an eigenvector with eigenvalue $f(\zeta^j)$.
The vectors $v_j$ form a Vandermonde/Fourier matrix and are linearly independent, hence these are all eigenvalues.
Therefore
\[
\det(C)=\prod_{j=0}^{n} f(\zeta^j).
\]
\end{solution}

\begin{problem}[CP-I-0110]
\label{prob:cp-i-0110}
\textbf{Normal operators and an orthonormal eigenbasis}

\textbf{Exact statement.}
\hspace*{1.5em}Let $V$ be a finite dimensional complex inner product space, and let $\alpha$ be an endomorphism on $V$.
Assume that $\alpha$ is \emph{normal}, i.e.\ $\alpha\alpha^*=\alpha^*\alpha$.
Show that $\alpha$ and $\alpha^*$ have a common eigenvector $v$, and that the corresponding eigenvalues are complex conjugates.
Show that the subspace $\langle v\rangle^\perp$ is invariant under both $\alpha$ and $\alpha^*$.
Deduce that there is an orthonormal basis of eigenvectors of $\alpha$.

\end{problem}

\begin{solution}
Since $V$ is complex and finite dimensional, $\alpha$ has an eigenvalue $\lambda$ and eigenvector $v\neq 0$:
\[
\alpha v=\lambda v.
\]
Consider $N:=\alpha-\lambda I$, which is normal because $\alpha$ is normal.
For any $x$,
\[
\|Nx\|^2=\langle N^*Nx,x\rangle,\qquad \|N^*x\|^2=\langle NN^*x,x\rangle,
\]
and normality gives $N^*N=NN^*$, hence $\|Nx\|=\|N^*x\|$ for all $x$.
In particular, $Nv=0$ implies $N^*v=0$, i.e.
\[
(\alpha^*-\overline{\lambda}I)v=0 \quad\Rightarrow\quad \alpha^*v=\overline{\lambda}\,v.
\]
Thus $v$ is a common eigenvector of $\alpha$ and $\alpha^*$, with conjugate eigenvalues.

If $w\perp v$, then
\[
\langle \alpha w,v\rangle=\langle w,\alpha^*v\rangle=\langle w,\overline{\lambda}v\rangle=0,
\]
so $\alpha w\perp v$. Similarly,
\[
\langle \alpha^* w,v\rangle=\langle w,\alpha v\rangle=\langle w,\lambda v\rangle=0,
\]
so $\alpha^*w\perp v$. Hence $\langle v\rangle^\perp$ is invariant under both $\alpha$ and $\alpha^*$.

Restrict $\alpha$ to $\langle v\rangle^\perp$. The restriction is again normal.
By induction on $\dim V$, $\alpha|_{\langle v\rangle^\perp}$ has an orthonormal basis of eigenvectors.
Together with $v$ (normalized), this yields an orthonormal basis of $V$ consisting of eigenvectors of $\alpha$.
\end{solution}

\begin{problem}[CP-I-0111]
\label{prob:cp-i-0111}
\textbf{Cayley--Hamilton for $2\times 2$, powers of a specific matrix}

\textbf{Exact statement.}
\hspace*{1.5em}Check, by direct calculation, that the Cayley--Hamilton Theorem holds for a general $2\times 2$ matrix.
Find the characteristic polynomial for
\[
A=\begin{pmatrix}3&4\\-1&-1\end{pmatrix}
\]
and deduce that $A^2=2A-I$. Is $A$ diagonalisable?
Show by induction that
\[
A^n=\alpha_n A+\beta_n I,\qquad n\ge 0,
\]
for real numbers $\alpha_n$ and $\beta_n$.
Solve the recurrences satisfied by $\alpha_n$ and $\beta_n$ and hence find $A^n$ explicitly.

\end{problem}

\begin{solution}
\textbf{(a) General $2\times 2$ case.}
Let
\[
A=\begin{pmatrix}a&b\\c&d\end{pmatrix},\qquad \operatorname{tr}(A)=a+d,\qquad \det(A)=ad-bc.
\]
Then
\[
p_A(t)=\det(tI-A)=t^2-(a+d)t+(ad-bc).
\]
Cayley--Hamilton claims
\[
p_A(A)=A^2-(\operatorname{tr} A)A+(\det A)I=0.
\]
Compute
\[
A^2=\begin{pmatrix}a^2+bc&ab+bd\\ac+dc&bc+d^2\end{pmatrix}.
\]
Then entrywise one checks:
\[
a^2+bc-(a+d)a+(ad-bc)=0,\quad ab+bd-(a+d)b=0,
\]
\[
ac+dc-(a+d)c=0,\quad bc+d^2-(a+d)d+(ad-bc)=0,
\]
so indeed $p_A(A)=0$.

\textbf{(b) The given matrix.}
For
\[
A=\begin{pmatrix}3&4\\-1&-1\end{pmatrix}
\]
we have $\operatorname{tr}(A)=2$ and $\det(A)=1$, hence
\[
p_A(t)=t^2-2t+1=(t-1)^2.
\]
By Cayley--Hamilton,
\[
A^2-2A+I=0 \quad\Longrightarrow\quad A^2=2A-I.
\]
The only eigenvalue is $1$ (algebraic multiplicity $2$). Since
\[
A-I=\begin{pmatrix}2&4\\-1&-2\end{pmatrix}
\]
has rank $1$, the eigenspace has dimension $1$, so $A$ is not diagonalizable.

\textbf{(c) Powers.}
Set $A^n=\alpha_n A+\beta_n I$. From $A^0=I$ and $A^1=A$ we have
\[
(\alpha_0,\beta_0)=(0,1),\qquad (\alpha_1,\beta_1)=(1,0).
\]
Assuming $A^n=\alpha_n A+\beta_n I$, multiply by $A$:
\[
A^{n+1}=\alpha_n A^2+\beta_n A=\alpha_n(2A-I)+\beta_n A=(2\alpha_n+\beta_n)A-\alpha_n I.
\]
Thus
\[
\alpha_{n+1}=2\alpha_n+\beta_n,\qquad \beta_{n+1}=-\alpha_n.
\]
Eliminating $\beta_n$ via $\beta_n=-\alpha_{n-1}$ gives
\[
\alpha_{n+1}=2\alpha_n-\alpha_{n-1},\qquad \alpha_0=0,\ \alpha_1=1.
\]
This yields $\alpha_n=n$, and then $\beta_n=-\alpha_{n-1}=1-n$. Therefore
\[
\boxed{\,A^n=nA+(1-n)I\,,\qquad n\ge 0.}
\]
\end{solution}

\begin{problem}[CP-I-0112]
\label{prob:cp-i-0112}
\textbf{Rank and signature of quadratic forms; diagonalization to $\pm 1$}

\textbf{Exact statement.}
\hspace*{1.5em}Find the rank and signature of the following quadratic forms over $\mathbb R$:
\[
q_1=x^2+y^2+z^2-2xz-2yz,\qquad
q_2=x^2+2y^2-2z^2-4xy-4yz,
\]
\[
q_3=16xy-z^2,\qquad
q_4=2xy+2yz+2zx.
\]
If $A$ is the matrix of $q_1$, find a non-singular matrix $P$ such that $P^TAP$ is diagonal with entries $\pm 1$.
\end{problem}

\begin{solution}
\textbf{(1) } For $q_1$ complete squares:
\[
q_1=(x-z)^2+(y-z)^2-z^2.
\]
Set $u=x-z$, $v=y-z$, $w=z$. Then $q_1=u^2+v^2-w^2$, so $\mathrm{rank}(q_1)=3$ and inertia/signature $(2,1)$.

To produce $P$ with $P^TAP=\mathrm{diag}(1,1,-1)$, write $(u,v,w)^T=T(x,y,z)^T$ with
\[
T=\begin{pmatrix}1&0&-1\\[2pt]0&1&-1\\[2pt]0&0&1\end{pmatrix},\qquad D=\mathrm{diag}(1,1,-1).
\]
Then $q_1=(x,y,z)\,A\,(x,y,z)^T=(u,v,w)\,D\,(u,v,w)^T$, hence $A=T^TDT$.
Taking $P=T^{-1}$ gives $P^TAP=D$:
\[
P=\begin{pmatrix}1&0&1\\[2pt]0&1&1\\[2pt]0&0&1\end{pmatrix},\qquad
P^TAP=\mathrm{diag}(1,1,-1).
\]

\textbf{(2) } For $q_2$:
\[
x^2-4xy+2y^2=(x-2y)^2-2y^2,
\]
so
\[
q_2=(x-2y)^2-\bigl(2y^2+4yz+2z^2\bigr)=(x-2y)^2-2(y+z)^2.
\]
Thus $\mathrm{rank}(q_2)=2$ with inertia $(1,1)$ and nullity $1$.

\textbf{(3) } For $q_3=16xy-z^2$, put $u=x+y$, $v=x-y$. Then $16xy=4(u^2-v^2)$, hence
\[
q_3=4u^2-4v^2-z^2.
\]
Therefore $\mathrm{rank}(q_3)=3$ and inertia $(1,2)$.

\textbf{(4) } For $q_4=2(xy+yz+zx)$, put
\[
u=x+y+z,\quad v=x-y,\quad w=x+y-2z.
\]
A direct substitution yields
\[
q_4=\frac{2}{3}u^2-\frac{1}{2}v^2-\frac{1}{6}w^2,
\]
so $\mathrm{rank}(q_4)=3$ and inertia $(1,2)$.
\end{solution}

\begin{problem}[CP-I-0113]
\label{prob:cp-i-0113}
\textbf{Spectral Theorem for Real Symmetric Matrices}

\textbf{Canonical authored completion of the retained spectral-theorem fragment.}

\hspace*{1.5em}Let \(A\in M_n(\mathbb R)\) satisfy \(A^T=A\).
State the spectral theorem and use it to prove that there is an orthogonal matrix \(P\) such that
\[
P^TAP=D
\]
is diagonal with real entries. Deduce that eigenvectors for distinct eigenvalues are orthogonal.
If \(A\) is invertible, show that
\[
A^{-1}=PD^{-1}P^T.
\]
\end{problem}

\begin{solution}
The real spectral theorem states that every real symmetric matrix admits an orthonormal basis
of eigenvectors and all its eigenvalues are real. Let \(u_1,\dots,u_n\) be such an orthonormal
eigenbasis, with
\[
Au_i=\lambda_i u_i.
\]
Put
\[
P=[u_1\ \cdots\ u_n].
\]
Then \(P^TP=I\), so \(P\) is orthogonal, and
\[
P^TAP=\operatorname{diag}(\lambda_1,\dots,\lambda_n)=D.
\]

For completeness, if \(Av=\lambda v\) and \(Aw=\mu w\), then symmetry gives
\[
\lambda\,v^Tw=(Av)^Tw=v^TAw=\mu\,v^Tw.
\]
Thus
\[
(\lambda-\mu)v^Tw=0,
\]
so \(v\perp w\) whenever \(\lambda\ne\mu\).

Finally,
\[
A=PDP^T.
\]
If \(A\) is invertible, every \(\lambda_i\ne0\), hence
\[
A^{-1}=PD^{-1}P^T.
\]
\end{solution}

\begin{problem}[CP-I-0114]
\label{prob:cp-i-0114}
\textbf{(Morse functions)}

\hspace*{1.5em}\textbf{Problem.}
	Let $X$ be a $k$-manifold and $f:X\to\mathbb{R}$ a smooth function.
	Recall that $x\in X$ is a \emph{critical point} if $df_x$ is not surjective, i.e.\ $df_x=0$.
	A critical point is \emph{non-degenerate} if, in local coordinates near $x$, the Hessian matrix
	$\big(\frac{\partial^2 f}{\partial x_i\partial x_j}\big)$ has non-vanishing determinant.
	If all critical points of $f$ are non-degenerate, $f$ is called a \emph{Morse function}.

		\par\noindent\textbullet\quad Show that the condition $\det\!\big(\frac{\partial^2 f}{\partial x_i\partial x_j}\big)\neq 0$
		is independent of the choice of chart.
		\par\noindent\textbullet\quad Suppose now that $X$ is an open subset of $\mathbb{R}^k$. Given $a\in\mathbb{R}^k$, set
		\[
		f_a(x)=f(x)+\langle x,a\rangle,
		\]
		where $\langle x,a\rangle$ is the usual inner product. Show that $f_a$ is a Morse function for a dense set
		of values of $a$. \emph{(Hint: consider $\nabla f:X\to\mathbb{R}^k$.)}
		\par\noindent\textbullet\quad Show that the determinant function on $M(n)$ is Morse if $n=2$, but not if $n>2$.

	\medskip

	\noindent\textbf{T1. Hessian under change of coordinates at a critical point.}
	Let $f$ be smooth and $p$ a critical point. If $y=y(x)$ is a coordinate change with Jacobian
	$J=\big(\frac{\partial x_i}{\partial y_\alpha}\big)$, then at $p$,
	\[
	\mathrm{Hess}_y(f)=J^T\,\mathrm{Hess}_x(f)\,J,
	\qquad
	\det(\mathrm{Hess}_y(f))=\det(J)^2\,\det(\mathrm{Hess}_x(f)).
	\]

	\medskip
	\noindent\textbf{T2. Regular values and Sard's theorem (used for generic Morse).}
	For a smooth map $F:X\to\mathbb{R}^k$, a value $v$ is \emph{regular} if $dF_x$ is surjective for all $x\in F^{-1}(v)$.
	Sard's theorem implies that the set of regular values of $F$ is dense in $\mathbb{R}^k$.

	\medskip
	\noindent\textbf{T3. Differential of $\det$ and critical matrices.}
	For $A\in M(n)$ and $H\in M(n)$,
	\[
	d(\det)_A(H)=\mathrm{tr}(\operatorname{adj}(A)^T H).
	\]
	Hence $A$ is critical for $\det$ iff $\operatorname{adj}(A)=0$, equivalently $\mathrm{rank}(A)\le n-2$.

	\medskip
\end{problem}

\begin{solution}
\noindent\textbf{(i)}
	Let $(x_1,\dots,x_k)$ and $(y_1,\dots,y_k)$ be two charts near $p\in X$ and write $F(y)=f(x(y))$.
	By the chain rule,
	\[
	\frac{\partial F}{\partial y_\alpha}
	=\sum_i \frac{\partial f}{\partial x_i}\frac{\partial x_i}{\partial y_\alpha}.
	\]
	Differentiating again,
	\[
	\frac{\partial^2 F}{\partial y_\alpha\partial y_\beta}
	=\sum_{i,j}\frac{\partial^2 f}{\partial x_i\partial x_j}
	\frac{\partial x_i}{\partial y_\alpha}\frac{\partial x_j}{\partial y_\beta}
	\;+\;\sum_i \frac{\partial f}{\partial x_i}\frac{\partial^2 x_i}{\partial y_\alpha\partial y_\beta}.
	\]
	At a critical point $p$, $\frac{\partial f}{\partial x_i}(p)=0$, so the second sum vanishes at $p$.
	Thus $\mathrm{Hess}_y(f)=J^T\mathrm{Hess}_x(f)J$ at $p$, and
	\[
	\det(\mathrm{Hess}_y(f))(p)=\det(J(p))^2\det(\mathrm{Hess}_x(f))(p).
	\]
	Since $\det(J(p))\neq 0$, we have $\det(\mathrm{Hess}_x(f))(p)\neq 0$ iff $\det(\mathrm{Hess}_y(f))(p)\neq 0$.

	\medskip
	\noindent\textbf{(ii)}
	Let $X\subset\mathbb{R}^k$ be open and define $f_a(x)=f(x)+\langle x,a\rangle$.
	Then
	\[
	\nabla f_a(x)=\nabla f(x)+a,
	\]
	so $x$ is a critical point of $f_a$ iff $\nabla f(x)=-a$.
	Moreover $\mathrm{Hess}(f_a)=\mathrm{Hess}(f)$ because $\langle x,a\rangle$ is linear.
	But $d(\nabla f)_x$ is represented by $\mathrm{Hess}(f)(x)$, hence $x$ is a non-degenerate critical point of $f_a$
	iff $-a$ is a regular value of $\nabla f$ at $x$.
	By Sard's theorem, the set of regular values of $\nabla f$ is dense in $\mathbb{R}^k$,
	so for a dense set of $a$ the function $f_a$ is Morse.

	\medskip
	\noindent\textbf{(iii)}
	Consider $\det:M(n)\to\mathbb{R}$. By \textbf{T3}, the critical points are exactly the matrices with
	$\mathrm{rank}(A)\le n-2$.

	\smallskip
	\noindent\emph{Case $n=2$.}
	Write $A=\begin{pmatrix}a&b\\c&d\end{pmatrix}$, so $f(a,b,c,d)=ad-bc$.
	The only critical point is $A=0$. The Hessian is constant:
	\[
	\mathrm{Hess}(f)=
	\begin{pmatrix}
		0&0&0&1\\
		0&0&-1&0\\
		0&-1&0&0\\
		1&0&0&0
	\end{pmatrix},
	\qquad
	\det(\mathrm{Hess}(f))=1\neq 0.
	\]
	Hence $\det$ is Morse on $M(2)$.

	\smallskip
	\noindent\emph{Case $n>2$.}
	The critical set contains the smooth stratum $\{A:\mathrm{rank}(A)=n-2\}$, which has positive dimension.
	On this set we have $\det\equiv 0$. Therefore, for any curve $\gamma(t)$ lying in this stratum with
	$\gamma(0)=A$ and $\gamma'(0)=\xi\neq 0$, the function $\det(\gamma(t))$ is identically $0$, so
	$\frac{d^2}{dt^2}\big|_{t=0}\det(\gamma(t))=0$.
	Thus the Hessian quadratic form at $A$ vanishes on a nonzero tangent vector $\xi$, so the Hessian is degenerate.
	Hence $\det$ is not Morse for $n>2$.

	\medskip
\end{solution}

\begin{problem}[CP-I-0117]
\label{prob:cp-i-0117}
\textbf{Eigenvalues/eigenvectors; reality + orthogonality \(\Leftrightarrow\) Hermitian}

\textbf{Exact statement.}
\hspace*{1.5em}Find the eigenvalues and eigenvectors of the matrix
\[
A=\begin{pmatrix}
	1&\alpha&0\\
	\beta&1&0\\
	0&0&1
\end{pmatrix},
\]
where \(\alpha,\beta\) are non-zero complex numbers. Find the conditions for which
(i) the eigenvalues are real and
(ii) eigenvectors for different eigenvalues are orthogonal.
Show that both these conditions hold if and only if \(A\) is hermitian.

\emph{Recall:} the complex inner product on \(\mathbb C^3\) is
\[
z^\dagger w=\overline z_1 w_1+\overline z_2 w_2+\overline z_3 w_3,
\]
and \(z,w\) are orthogonal if \(z^\dagger w=0\).

\end{problem}

\begin{solution}
Write \(A=\mathrm{diag}(B,1)\) with
\[
B=\begin{pmatrix}1&\alpha\\ \beta&1\end{pmatrix}.
\]
Then
\[
\chi_B(\lambda)=\det(B-\lambda I)=(1-\lambda)^2-\alpha\beta,
\]
so with \(s=\sqrt{\alpha\beta}\) we have eigenvalues
\[
\lambda_1=1+s,\qquad \lambda_2=1-s,\qquad \lambda_3=1.
\]
For \(\lambda=1+s\), solving \((B-(1+s)I)v=0\) gives \(-su+\alpha v=0\), so we may take
\[
v_+=\begin{pmatrix}\alpha\\ s\\ 0\end{pmatrix}.
\]
For \(\lambda=1-s\), we may take
\[
v_-=\begin{pmatrix}\alpha\\ -s\\ 0\end{pmatrix}.
\]
For \(\lambda=1\), an eigenvector is \(v_0=e_3=(0,0,1)^T\).

\textbf{(i) Real eigenvalues.}
The eigenvalues are real iff \(s\in\mathbb R\), i.e.\ iff \(\alpha\beta\in(0,\infty)\subset\mathbb R\).

\textbf{(ii) Orthogonality for distinct eigenvalues.}
Clearly \(v_0\perp v_\pm\). Thus we require \(v_+^\dagger v_-=0\):
\[
v_+^\dagger v_-=\overline\alpha\,\alpha+\overline s(-s)=|\alpha|^2-|s|^2.
\]
If (i) holds then \(s\in\mathbb R\) so \(|s|^2=s^2=\alpha\beta\), hence
\[
v_+^\dagger v_-=0 \iff |\alpha|^2=\alpha\beta \iff \beta=\overline\alpha \quad(\alpha\neq 0).
\]
But \(\beta=\overline\alpha\) is exactly the condition that \(A\) is hermitian. Conversely, if \(A\) is hermitian then
\(\alpha\beta=\alpha\overline\alpha=|\alpha|^2>0\) so (i) holds, and eigenvectors for distinct eigenvalues are orthogonal.
\end{solution}

\begin{problem}[CP-I-0118]
\label{prob:cp-i-0118}
\textbf{Spectral Properties of Anti-Hermitian and Unitary Matrices}

\textbf{Canonical authored completion of the retained spectral fragment.}

\hspace*{1.5em}Prove the following.
Prove that if \(A^\dagger=-A\), then every eigenvalue of \(A\) is purely imaginary.

Prove that if \(U^\dagger U=I\), then every eigenvalue of \(U\) has modulus \(1\).

Finally, prove that in either case eigenvectors corresponding to distinct eigenvalues are orthogonal.
\end{problem}

\begin{solution}
Let \(Av=\lambda v\) with \(v\neq0\). If \(A^\dagger=-A\), then
\[
\langle Av,v\rangle
=
\lambda\langle v,v\rangle.
\]
On the other hand,
\[
\langle Av,v\rangle
=
\langle v,A^\dagger v\rangle^{\!*}
=
-\overline{\langle Av,v\rangle},
\]
so \(\langle Av,v\rangle\) is purely imaginary. Since \(\langle v,v\rangle>0\) is real,
\(\lambda\in i\mathbb R\).

If \(Uv=\lambda v\) and \(U\) is unitary, then
\[
\|v\|=\|Uv\|=|\lambda|\,\|v\|,
\]
hence \(|\lambda|=1\).

Finally, anti-Hermitian matrices are normal because
\[
AA^\dagger=A(-A)=-A^2=(-A)A=A^\dagger A,
\]
and unitary matrices are normal because
\[
UU^\dagger=U^\dagger U=I.
\]
For a normal matrix, eigenvectors belonging to distinct eigenvalues are orthogonal. Therefore
the claimed orthogonality holds in both cases.
\end{solution}

\begin{problem}[CP-I-0120]
\label{prob:cp-i-0120}
\textbf{Orthonormal eigenbases and quadric surfaces}

\textbf{Exact statement.}
\hspace*{1.5em}Find all eigenvalues, and an orthonormal set of eigenvectors, of
\[
A=\begin{pmatrix}5&0&\sqrt3\\[2pt]0&3&0\\[2pt]\sqrt3&0&3\end{pmatrix},
\qquad
B=\begin{pmatrix}2&-1&-1\\[2pt]-1&2&-1\\[2pt]-1&-1&2\end{pmatrix}.
\]
Hence sketch the surfaces
\[
5x^2+3y^2+3z^2+2\sqrt3\,xz=1,
\qquad
x^2+y^2+z^2-xy-yz-zx=1.
\]

\end{problem}

\begin{solution}
\textbf{(A)} The eigenvalues are
\[
\lambda=2,\ 3,\ 6.
\]
An orthonormal eigenbasis is
\[
u_{2}=\begin{pmatrix}-\tfrac12\\[2pt]0\\[2pt]\tfrac{\sqrt3}{2}\end{pmatrix},\quad
u_{3}=\begin{pmatrix}0\\[2pt]1\\[2pt]0\end{pmatrix},\quad
u_{6}=\begin{pmatrix}\tfrac{\sqrt3}{2}\\[2pt]0\\[2pt]\tfrac12\end{pmatrix},
\]
corresponding to $\lambda=2,3,6$ respectively.
Let
\[
U=u_2\cdot(x,y,z)=\frac{-x+\sqrt3 z}{2},\quad
V=u_3\cdot(x,y,z)=y,\quad
W=u_6\cdot(x,y,z)=\frac{\sqrt3 x+z}{2}.
\]
Then
\[
(x,y,z)A(x,y,z)^T=2U^2+3V^2+6W^2,
\]
so the surface $5x^2+3y^2+3z^2+2\sqrt3\,xz=1$ is the ellipsoid
\[
2U^2+3V^2+6W^2=1,
\]
with semi-axis lengths $1/\sqrt2$, $1/\sqrt3$, $1/\sqrt6$ (a rotation about the $y$-axis).

\textbf{(B)} The eigenvalues of $B$ are
\[
0\ (\text{mult.\ }1),\qquad 3\ (\text{mult.\ }2).
\]
An orthonormal eigenbasis is
\[
v_0=\frac1{\sqrt3}\begin{pmatrix}1\\1\\1\end{pmatrix},\quad
v_{3,1}=\frac1{\sqrt2}\begin{pmatrix}1\\-1\\0\end{pmatrix},\quad
v_{3,2}=\frac1{\sqrt6}\begin{pmatrix}1\\1\\-2\end{pmatrix}.
\]
If $(\alpha,\beta,\gamma)$ are coordinates in this basis then
\[
(x,y,z)B(x,y,z)^T=3\beta^2+3\gamma^2.
\]
Moreover,
\[
x^2+y^2+z^2-xy-yz-zx=\tfrac12(x,y,z)B(x,y,z)^T,
\]
so the surface $x^2+y^2+z^2-xy-yz-zx=1$ becomes
\[
\frac12(3\beta^2+3\gamma^2)=1
\quad\Longleftrightarrow\quad
\beta^2+\gamma^2=\frac23,
\]
with $\alpha$ free. Hence it is a circular cylinder of radius $\sqrt{2/3}$ whose axis is the line spanned by $(1,1,1)$.
\end{solution}

\section{Euclidean and Hilbert-Space Geometry}

\paragraph{Related material.} Volume I, Chapters \texttt{I/13--I/18}.

\begin{problem}[CP-I-0001]
\label{prob:cp-i-0001}
\textbf{Operator norm: basic and useful properties}

\hspace*{1.5em}\textbf{Statement.}\quad

With Euclidean norms, let $\|A\|_{op}=\sup_{\|x\|=1}\|Ax\|$ for $A\in\mathcal{L}(\mathbb{R}^n;\mathbb{R}^m)$.
Prove:
 \setlength{\itemsep}{6pt}
	\par\noindent\textbullet\quad $\displaystyle \|A\|_{op}=\sup_{x\neq 0}\frac{\|Ax\|}{\|x\|}$.
	\par\noindent\textbullet\quad For $A\in\mathcal{L}(\mathbb{R}^n;\mathbb{R}^m)$, $B\in\mathcal{L}(\mathbb{R}^m;\mathbb{R}^p)$,
	\ $\|B\circ A\|_{op}\le \|B\|_{op}\,\|A\|_{op}$.
	\par\noindent\textbullet\quad If $A\in\mathcal{L}(\mathbb{R}^n;\mathbb{R})$, then $\exists\,a\in\mathbb{R}^n$ with $Ax=\langle a,x\rangle$ and $\|A\|_{op}=\|a\|$.
	\par\noindent\textbullet\quad If $A\in\mathcal{L}(\mathbb{R};\mathbb{R}^m)$, then $\exists\,a\in\mathbb{R}^m$ with $Ax=xa$ and $\|A\|_{op}=\|a\|$.
	\par\noindent\textbullet\quad If $(A_{ij})$ is the matrix of $A$ in the standard bases, then
	\[
	\frac{1}{\sqrt{n}}\Big(\sum_{j=1}^n\sum_{i=1}^m A_{ij}^2\Big)^{1/2}
	\ \le\ \|A\|_{op}\ \le\
	\Big(\sum_{j=1}^n\sum_{i=1}^m A_{ij}^2\Big)^{1/2},
	\]
	and the right inequality is an equality iff $A=0$ or $\operatorname{rank}(A)=1$.

\bigskip\hrule\bigskip

\textbf{Theory \& Solution.}\quad

\textbf{(1) Homogeneity).}\;
$\|Ax\|/\|x\|$ depends only on $x/\|x\|$; taking $x$ on the unit sphere yields the same supremum.

\medskip

\textbf{(2) Submultiplicativity).}\;
For $\|x\|=1$, $\|B(Ax)\|\le \|B\|_{op}\|Ax\|\le \|B\|_{op}\|A\|_{op}$.
Taking the supremum gives $\|B\circ A\|_{op}\le \|B\|_{op}\|A\|_{op}$.

\medskip

\textbf{(3) Linear functionals as inner products).}\;
Finite-dimensional Riesz: $\exists a$ with $Ax=\langle a,x\rangle$ and by Cauchy–Schwarz
$\|Ax\|\le \|a\|\,\|x\|$, hence $\|A\|_{op}=\|a\|$.

\medskip

\textbf{(4) Maps $\mathbb{R}\to\mathbb{R}^m$).}\;
Let $a:=A(1)$. Then $Ax=A(x\cdot 1)=xa$ and
$\|Ax\|=\|x a\|=|x|\,\|a\|$, so $\|A\|_{op}=\|a\|$.

\medskip

\textbf{(5) Frobenius vs.\ operator norm; bounds and equality).}\;
Let $\|A\|_{F} := \big(\sum_{i,j} A_{ij}^2\big)^{1/2}$. For any unit $x$,
\[
\|Ax\|^2=\sum_{i=1}^{m}\Big(\sum_{j=1}^{n} A_{ij}x_j\Big)^2
\ \le\ \sum_{i=1}^{m}\Big(\sum_{j=1}^{n} A_{ij}^2\Big)\Big(\sum_{j=1}^{n}x_j^2\Big)
\ =\ \|A\|_F^2,
\]
so $\|A\|_{op}\le \|A\|_F$.
For the lower bound, let $e_1,\dots,e_n$ be the standard orthonormal basis of $\mathbb{R}^n$. Since $\|Ae_j\|\le \|A\|_{op}$ for every $j$,
\[
\|A\|_F^2=\sum_{j=1}^{n}\|Ae_j\|^2\le n\,\|A\|_{op}^2.
\]
Therefore $\|A\|_{op}\ge \|A\|_F/\sqrt{n}$.
Equality $\|A\|_{op}=\|A\|_F$ holds iff $A$ has exactly one nonzero singular value,
i.e.\ $A=0$ or $\operatorname{rank}(A)=1$. \hfill$\square$

\bigskip\hrule\bigskip

\textbf{Examples.}\quad

 \setlength{\itemsep}{6pt}
	\par\noindent\textbullet\quad Rank–one matrix $A=uv^\top$: $\|A\|_{op}=\|u\|\,\|v\|=\|A\|_F$ (right equality case).
	\par\noindent\textbullet\quad Identity $A=I_n$: $\|A\|_{op}=1$, $\|A\|_F=\sqrt{n}$; the bounds read
	$1\le 1 \le \sqrt{n}$.

\bigskip\hrule\bigskip
\end{problem}

\begin{problem}[CP-I-0006]
\label{prob:cp-i-0006}
\textbf{Householder reflections and ``two reflections = rotation''}

\hspace*{1.5em}Reflection in the plane through the origin with unit normal $n$ is
\[
H(n)=I-2nn^T,\qquad H(n)^2=I,\qquad \det H(n)=-1.
\]
The composition of reflections in two planes through the origin is a rotation about their line of intersection; the rotation angle equals twice the angle between the planes (equivalently twice the angle between their normals).
Moreover, $-H(n)$ fixes $n$ and flips every vector orthogonal to $n$, hence it is a rotation by $\pi$ about axis $n$.
\end{problem}

\begin{problem}[CP-I-0021]
\label{prob:cp-i-0021}
\par\noindent\textbullet\quad Show that the inner product has the following properties for all $x,y,z\in\mathbb{R}^n$ and $a\in\mathbb{R}$:
		\[
		\langle x,y\rangle=\langle y,x\rangle,\qquad
		\langle x+y,z\rangle=\langle x,z\rangle+\langle y,z\rangle,\qquad
		\langle ax,y\rangle=a\langle x,y\rangle.
		\]
		\textit{Solution.} These follow from the Euclidean inner product $\langle x,y\rangle=\sum_{i=1}^n x_i y_i$.

		\par\noindent\textbullet\quad For $t\in\mathbb{R}$ and $x,y\in\mathbb{R}^n$ show
		\[
		\|x+ty\|^2=\|x\|^2+2t\langle x,y\rangle+t^2\|y\|^2\ge 0.
		\]
		\textit{Solution.} Expand $\langle x+ty,x+ty\rangle$.

		\par\noindent\textbullet\quad Deduce Cauchy--Schwarz: $|\langle x,y\rangle|\le \|x\|\,\|y\|$.\\
		\textit{Solution.} View $\|x+ty\|^2$ as a quadratic in $t$ and use non-positivity of its discriminant.
		Equality iff $x,y$ are linearly dependent.

		\par\noindent\textbullet\quad Deduce the triangle inequality $\|x+y\|\le \|x\|+\|y\|$.

		\par\noindent\textbullet\quad Show the reverse triangle inequality $|\|x\|-\|y\||\le \|x-y\|$.

		\par\noindent\textbullet\quad Suppose $x=(x^1,\dots,x^n)^t$.

			\par\noindent\textbullet\quad $\displaystyle \max_{k}|x^k|\le \|x\|$.
			\par\noindent\textbullet\quad $\displaystyle \|x\|\le \sqrt{n}\,\max_{k}|x^k|$.

		\textit{Solution.} Immediate from $\|x\|^2=\sum (x^i)^2$.
\end{problem}

\begin{problem}[CP-I-0031]
\label{prob:cp-i-0031}
\textbf{Apollonius circle; a coaxal family and orthogonal intersections}

\textbf{Exact statement.}
\hspace*{1.5em}(a) Let \(z,a,b\in\mathbb C\) (\(a\neq b\)) correspond to points \(P,A,B\) in the Argand diagram.
Let \(C_\lambda\) be the locus of \(P\) defined by
\[
\frac{|PA|}{|PB|}=\lambda,
\]
where \(\lambda\) is a fixed positive real constant. Show that \(C_\lambda\) is a circle if \(\lambda\neq 1\),
and find its centre and radius. What happens if \(\lambda=1\)?

(b) For the case \(a=-b=p\) with \(p\in\mathbb R\), and for each fixed \(\mu\in\mathbb R\), show that
\[
S_\mu=\left\{z\in\mathbb C:\ |z-i\mu|=\sqrt{p^2+\mu^2}\right\}
\]
is a circle passing through \(A\) and \(B\) with its centre on the perpendicular bisector of \(AB\).
Show that the circles \(C_\lambda\) and \(S_\mu\) intersect orthogonally for all \(\lambda,\mu\).

\end{problem}

\begin{solution}
\textbf{(a)} The condition \(|PA|/|PB|=\lambda\) is
\[
|z-a|=\lambda |z-b|
\quad\Longleftrightarrow\quad
|z-a|^2=\lambda^2|z-b|^2.
\]
Expand:
\[
|z|^2- z\bar a-\bar z a+|a|^2=\lambda^2\bigl(|z|^2-z\bar b-\bar z b+|b|^2\bigr).
\]
So
\[
(1-\lambda^2)|z|^2- z(\bar a-\lambda^2\bar b)-\bar z(a-\lambda^2 b)+(|a|^2-\lambda^2|b|^2)=0.
\]
If \(\lambda\neq 1\), divide by \(1-\lambda^2\) and complete the square:
\[
\left|\,z-\frac{a-\lambda^2 b}{\,1-\lambda^2\,}\right|^2
=\frac{\lambda^2|a-b|^2}{(1-\lambda^2)^2}.
\]
Hence \(C_\lambda\) is a circle with centre
\[
c_\lambda=\frac{a-\lambda^2 b}{1-\lambda^2}
\]
and radius
\[
r_\lambda=\frac{\lambda|a-b|}{|1-\lambda^2|}.
\]
If \(\lambda=1\), the condition is \(|z-a|=|z-b|\): the locus is the perpendicular bisector of \(AB\).

\medskip
\noindent\textbf{(b)} Here \(a=p\), \(b=-p\) with \(p\in\mathbb R\). The set
\[
S_\mu:\ |z-i\mu|=\sqrt{p^2+\mu^2}
\]
is a circle with centre \(i\mu\), which lies on the imaginary axis (the perpendicular bisector of \(AB\)).
It passes through \(A=p\) and \(B=-p\) since
\[
|p-i\mu|=\sqrt{p^2+\mu^2},\qquad |-p-i\mu|=\sqrt{p^2+\mu^2}.
\]
From (a),
\[
c_\lambda=\frac{p-\lambda^2(-p)}{1-\lambda^2}=p\frac{1+\lambda^2}{1-\lambda^2}\in\mathbb R,\qquad
r_\lambda=\frac{\lambda|2p|}{|1-\lambda^2|}.
\]
Two circles meet orthogonally iff the squared distance between centres equals the sum of squared radii:
\[
|c_\lambda-i\mu|^2=r_\lambda^2+(p^2+\mu^2).
\]
But
\[
|c_\lambda-i\mu|^2=c_\lambda^2+\mu^2,
\]
so it suffices to show \(c_\lambda^2=r_\lambda^2+p^2\). Indeed,
\[
c_\lambda^2=p^2\frac{(1+\lambda^2)^2}{(1-\lambda^2)^2},\qquad
r_\lambda^2=\frac{4\lambda^2p^2}{(1-\lambda^2)^2},
\]
hence
\[
c_\lambda^2-r_\lambda^2
=\frac{p^2\bigl((1+\lambda^2)^2-4\lambda^2\bigr)}{(1-\lambda^2)^2}
=\frac{p^2(1-\lambda^2)^2}{(1-\lambda^2)^2}
=p^2.
\]
Therefore \(C_\lambda\) and \(S_\mu\) intersect orthogonally for all \(\lambda\neq 1\).
If \(\lambda=1\), then \(C_1\) is the imaginary axis, which passes through the centre \(i\mu\) of \(S_\mu\),
so it meets \(S_\mu\) orthogonally as well.
\end{solution}

\begin{problem}[CP-I-0041]
\label{prob:cp-i-0041}
\textbf{An $L^2[-1,1]$ inner product on $P_n$ and Legendre-type eigenvectors}

\textbf{Exact statement.}
\hspace*{1.5em}Let $P_n$ be the $(n+1)$-dimensional space of real polynomials of degree $\le n$. Define
\[
(f,g)=\int_{-1}^{1} f(t)g(t)\,dt.
\]
Show that $(\ ,\ )$ is an inner product on $P_n$ and that the endomorphism $\alpha:P_n\to P_n$ defined by
\[
\alpha(f)(t)=(1-t^2)f''(t)-2t f'(t)
\]
is self-adjoint.
If $f$ is an eigenvector of $\alpha$ of degree $k$, what is the corresponding eigenvalue?
Why must $\alpha$ have precisely one monic eigenvector of degree $k$ for each $0\le k\le n$?

Let $s_k\in P_n$ be defined by $s_k(t)=\dfrac{d^k}{dt^k}(1-t^2)^k$. Prove:

	\par\noindent (i)\quad For $i\ne j$, $(s_i,s_j)=0$.
	\par\noindent (ii)\quad $s_0,\dots,s_n$ forms a basis for $P_n$.
	\par\noindent (iii)\quad For all $1\le k\le n$, $s_k$ spans the orthogonal complement of $P_{k-1}$ in $P_k$.
	\par\noindent (iv)\quad $s_k$ is an eigenvector of $\alpha$.

What is the relation between the $s_k$ and the result of applying Gram--Schmidt to the sequence $1,t,t^2,t^3,\dots$?
Explain why that is the case.
\end{problem}

\begin{solution}
\textbf{Inner product.}
Bilinearity and symmetry are clear. Moreover,
\[
(f,f)=\int_{-1}^{1} f(t)^2\,dt\ge 0,
\]
and $(f,f)=0$ implies $f(t)=0$ for all $t\in[-1,1]$, hence $f\equiv 0$ as a polynomial.
So $(\ ,\ )$ is an inner product on $P_n$.

\textbf{Self-adjointness of $\alpha$.}
Note that
\[
\alpha(f)=((1-t^2)f')'.
\]
Then integration by parts yields
\[
(\alpha f,g)=\int_{-1}^{1} ((1-t^2)f')'g\,dt
=\Big[(1-t^2)f'g\Big]_{-1}^{1}-\int_{-1}^{1}(1-t^2)f'g'\,dt.
\]
Since $1-t^2=0$ at $t=\pm 1$, the boundary term vanishes, so
\[
(\alpha f,g)=-\int_{-1}^{1}(1-t^2)f'g'\,dt.
\]
The right-hand side is symmetric in $f,g$, hence $(\alpha f,g)=(f,\alpha g)$ and $\alpha$ is self-adjoint.

\textbf{Eigenvalue for degree $k$.}
Let $f(t)=t^k+\text{(lower degree)}$ with $k\ge 1$. Then
\[
f'(t)=k t^{k-1}+\cdots,\qquad f''(t)=k(k-1)t^{k-2}+\cdots,
\]
so
\[
(1-t^2)f''(t)=k(k-1)t^{k-2}-k(k-1)t^k+\cdots,\qquad -2t f'(t)=-2k t^k+\cdots.
\]
Hence the leading term of $\alpha(f)$ is
\[
-(k(k-1)+2k)t^k+\cdots=-k(k+1)t^k+\cdots.
\]
If $\alpha(f)=\lambda f$ and $\deg f=k$, comparing leading coefficients gives
\[
\lambda=-k(k+1).
\]

\textbf{One monic eigenvector in each degree.}
In the ordered basis $(1,t,t^2,\dots,t^n)$, $\alpha$ maps $P_k$ to itself and is upper triangular, with diagonal entries
\[
\lambda_k=-k(k+1)\qquad (k=0,1,\dots,n),
\]
which are all distinct. Therefore each eigenspace is $1$-dimensional, and in each degree $k$ there is a unique eigenvector with leading coefficient $1$ (monic).

\textbf{Now define } $s_k(t)=\frac{d^k}{dt^k}(1-t^2)^k$.

\textbf{Degree.} Since $(1-t^2)^k$ has leading term $(-1)^k t^{2k}$, differentiating $k$ times gives
\[
s_k(t) = \frac{d^k}{dt^k}\bigl((-1)^k t^{2k} + \cdots\bigr)
= (-1)^k \frac{(2k)!}{k!}\, t^k + \cdots
\]
so $\deg s_k=k$ and $s_k\neq 0$.

\textbf{(i) Orthogonality.}
Let $p$ be any polynomial with $\deg p<k$. Then repeated integration by parts gives
\[
\int_{-1}^{1} p(t)\,s_k(t)\,dt
=\int_{-1}^{1} p(t)\,\frac{d^k}{dt^k}(1-t^2)^k\,dt
=(-1)^k\int_{-1}^{1} p^{(k)}(t)\,(1-t^2)^k\,dt=0,
\]
because $p^{(k)}\equiv 0$. The boundary terms vanish since $(1-t^2)^k$ has a zero of order $k$ at $t=\pm 1$.
Taking $p=s_i$ with $i<k$ gives $(s_i,s_k)=0$, hence $(s_i,s_j)=0$ for $i\ne j$.

\textbf{(ii) Basis.}
The polynomials $s_0,\dots,s_n$ have distinct degrees $0,1,\dots,n$ with nonzero leading coefficients, so they are linearly independent.
Since $\dim P_n=n+1$, they form a basis.

\textbf{(iii) Orthogonal complement in } $P_k$.
From the argument in (i), $s_k\perp P_{k-1}$ and $s_k\in P_k$.
But $\dim P_k-\dim P_{k-1}=1$, so the orthogonal complement of $P_{k-1}$ in $P_k$ is $1$-dimensional; hence it is $\mathrm{span}(s_k)$.

\textbf{(iv) $s_k$ is an eigenvector.}
Since $\alpha(P_{k-1})\subseteq P_{k-1}$ and $\alpha$ is self-adjoint, the orthogonal complement $(P_{k-1})^\perp\cap P_k$ is invariant under $\alpha$:
if $w\perp P_{k-1}$ and $u\in P_{k-1}$ then
\[
(\alpha w,u)=(w,\alpha u)=0.
\]
By (iii), $(P_{k-1})^\perp\cap P_k=\mathrm{span}(s_k)$, so $\alpha(s_k)=\mu_k s_k$ for some scalar $\mu_k$.
Comparing leading terms as above forces $\mu_k=-k(k+1)$, so
\[
\alpha(s_k)=-k(k+1)s_k.
\]

\textbf{Gram--Schmidt relation.}
Applying Gram--Schmidt to $1,t,t^2,\dots$ in the inner product $(f,g)=\int_{-1}^{1}fg$
produces orthogonal polynomials $p_k\in P_k$ such that $p_k\perp P_{k-1}$.
By (iii) this orthogonal complement is $1$-dimensional and spanned by $s_k$.
Therefore each Gram--Schmidt polynomial $p_k$ is a nonzero scalar multiple of $s_k$
(these are the Legendre polynomials up to normalization).
\end{solution}

\begin{problem}[CP-I-0047]
\label{prob:cp-i-0047}
\textbf{Spherical unit vectors form an orthonormal right-handed frame}

\textbf{Exact statement.}
\hspace*{1.5em}The vectors \(\mathbf e_r,\mathbf e_\theta,\mathbf e_\phi\) are defined in terms of the standard basis vectors
\(\mathbf i,\mathbf j,\mathbf k\) by
\[
\mathbf e_r=\cos\phi\,\sin\theta\,\mathbf i+\sin\phi\,\sin\theta\,\mathbf j+\cos\theta\,\mathbf k,
\]
\[
\mathbf e_\theta=\cos\phi\,\cos\theta\,\mathbf i+\sin\phi\,\cos\theta\,\mathbf j-\sin\theta\,\mathbf k,
\]
\[
\mathbf e_\phi=-\sin\phi\,\mathbf i+\cos\phi\,\mathbf j,
\]
where \(\theta\) and \(\phi\) are real numbers. Show, as efficiently as possible, that
\(\mathbf e_r,\mathbf e_\theta,\mathbf e_\phi\) are an orthonormal right-handed set.

\end{problem}

\begin{solution}
In components,
\begin{center}
\resizebox{\linewidth}{!}{$\displaystyle
\mathbf e_r=(\cos\phi\sin\theta,\ \sin\phi\sin\theta,\ \cos\theta),\quad
\mathbf e_\theta=(\cos\phi\cos\theta,\ \sin\phi\cos\theta,\ -\sin\theta),\quad
\mathbf e_\phi=(-\sin\phi,\ \cos\phi,\ 0).
$}
\end{center}

\medskip
\noindent\textbf{Unit lengths.}
\[
|\mathbf e_r|^2=\sin^2\theta(\cos^2\phi+\sin^2\phi)+\cos^2\theta=1,\qquad
|\mathbf e_\theta|^2=\cos^2\theta(\cos^2\phi+\sin^2\phi)+\sin^2\theta=1,
\]
\[
|\mathbf e_\phi|^2=\sin^2\phi+\cos^2\phi=1.
\]

\medskip
\noindent\textbf{Orthogonality.}
\[
\mathbf e_r\cdot\mathbf e_\theta
=\sin\theta\cos\theta(\cos^2\phi+\sin^2\phi)-\sin\theta\cos\theta=0,
\]
\[
\mathbf e_r\cdot\mathbf e_\phi
=\cos\phi\sin\theta(-\sin\phi)+\sin\phi\sin\theta(\cos\phi)=0,
\]
\[
\mathbf e_\theta\cdot\mathbf e_\phi
=\cos\phi\cos\theta(-\sin\phi)+\sin\phi\cos\theta(\cos\phi)=0.
\]
Hence the set is orthonormal.

\medskip
\noindent\textbf{Right-handedness.}
\[
\mathbf e_\theta\times \mathbf e_\phi
=
\begin{vmatrix}
	\mathbf i&\mathbf j&\mathbf k\\
	\cos\phi\cos\theta&\sin\phi\cos\theta&-\sin\theta\\
	-\sin\phi&\cos\phi&0
\end{vmatrix}
=(\sin\theta\cos\phi,\ \sin\theta\sin\phi,\ \cos\theta)=\mathbf e_r.
\]
Therefore \(\mathbf e_\theta\times\mathbf e_\phi=\mathbf e_r\), so
\((\mathbf e_r,\mathbf e_\theta,\mathbf e_\phi)\) is a right-handed orthonormal set.
\end{solution}

\begin{problem}[CP-I-0068]
\label{prob:cp-i-0068}
\textbf{Sharp Comparison of the Euclidean and Maximum Norms}

\textbf{Canonical authored differentiation of a duplicated imported exercise.}

\hspace*{1.5em}For \(x=(x_1,\dots,x_n)\in\mathbb{R}^n\), define
\[
\|x\|_2=\left(\sum_{k=1}^{n}x_k^2\right)^{1/2},
\qquad
\|x\|_\infty=\max_{1\le k\le n}|x_k|.
\]
Prove the sharp inequalities
\[
\|x\|_\infty\le \|x\|_2\le \sqrt n\,\|x\|_\infty.
\]
Determine all equality cases and conclude that \(\|\cdot\|_2\) and \(\|\cdot\|_\infty\) induce the same
topology on \(\mathbb{R}^n\).
\end{problem}

\begin{solution}
Choose \(j\) with \(|x_j|=\|x\|_\infty\). Then
\[
\|x\|_2^2=\sum_{k=1}^{n}x_k^2\ge x_j^2=\|x\|_\infty^2,
\]
hence
\[
\|x\|_\infty\le \|x\|_2.
\]
Equality holds precisely when every coordinate except possibly one is zero.

For the other direction,
\[
x_k^2\le \|x\|_\infty^2
\]
for every \(k\), so
\[
\|x\|_2^2=\sum_{k=1}^{n}x_k^2
\le n\|x\|_\infty^2.
\]
Therefore
\[
\|x\|_2\le \sqrt n\,\|x\|_\infty.
\]
Equality holds precisely when
\[
|x_1|=\cdots=|x_n|=\|x\|_\infty.
\]

The constants are sharp: \(e_1\) gives equality in the first inequality, while
\((1,\dots,1)\) gives equality in the second. Since
\[
\|x\|_\infty\le \|x\|_2\le \sqrt n\,\|x\|_\infty,
\]
the two norms are equivalent and therefore induce the same topology on \(\mathbb{R}^n\).
\end{solution}

\begin{problem}[CP-I-0073]
\label{prob:cp-i-0073}
\textbf{Positive definite operators: square roots and polar decomposition}

\textbf{Exact statement.}
\hspace*{1.5em}An endomorphism $\alpha$ of a finite dimensional inner product space $V$ is \emph{positive definite} if it is self-adjoint and satisfies
$\langle \alpha(x),x\rangle>0$ for all non-zero $x\in V$.

	\par\noindent (i)\quad Prove that a positive definite endomorphism has a unique positive definite square root.
	\par\noindent (ii)\quad Let $\alpha$ be an invertible endomorphism of $V$ and $\alpha^*$ its adjoint. By considering $\alpha^*\alpha$, show that $\alpha$
	can be factored as $\beta\gamma$ with $\beta$ unitary and $\gamma$ positive definite.

\end{problem}

\begin{solution}
\textbf{(i)} Since $\alpha$ is self-adjoint, the spectral theorem gives an orthonormal basis $(v_i)$ with
\[
\alpha v_i=\lambda_i v_i,\qquad \lambda_i\in\mathbb R.
\]
Positive definiteness implies $\lambda_i>0$ for all $i$ (otherwise $\langle \alpha v_i,v_i\rangle=\lambda_i\|v_i\|^2\le 0$).
Define $\sqrt{\alpha}$ by
\[
\sqrt{\alpha}\,v_i=\sqrt{\lambda_i}\,v_i.
\]
Then $(\sqrt{\alpha})^2=\alpha$, and $\sqrt{\alpha}$ is self-adjoint and positive definite.

For uniqueness, let $T$ be positive definite with $T^2=\alpha$.
If $\alpha v=\lambda v$, then $T^2v=\lambda v$, so $Tv=\pm\sqrt{\lambda}\,v$ on the eigenspace.
Positivity forces $Tv=\sqrt{\lambda}\,v$, hence $T=\sqrt{\alpha}$.

\textbf{(ii)} The operator $\alpha^*\alpha$ is self-adjoint and for $x\neq 0$,
\[
\langle (\alpha^*\alpha)x,x\rangle=\langle \alpha x,\alpha x\rangle>0,
\]
so $\alpha^*\alpha$ is positive definite. Let
\[
\gamma:=(\alpha^*\alpha)^{1/2}.
\]
Define
\[
\beta:=\alpha\,\gamma^{-1}.
\]
Then
\[
\beta^*\beta=\gamma^{-1}\alpha^*\alpha\gamma^{-1}=\gamma^{-1}\gamma^2\gamma^{-1}=I,
\]
so $\beta$ is unitary, and $\alpha=\beta\gamma$ with $\gamma$ positive definite (polar decomposition).
\end{solution}

\begin{problem}[CP-I-0106]
\label{prob:cp-i-0106}
\textbf{A positive definite quadratic form; orthonormal basis; Gram--Schmidt}

\textbf{Exact statement.}
\hspace*{1.5em}Show that the quadratic form $2(x^2+y^2+z^2+xy+yz+zx)$ is positive definite.
Write down an orthonormal basis for the corresponding inner product on $\mathbb R^3$.
Compute the basis of $\mathbb R^3$ obtained by applying the Gram--Schmidt process to the standard basis with respect to this inner product.

\end{problem}

\begin{solution}
Let
\[
q(x,y,z)=2(x^2+y^2+z^2+xy+yz+zx)=\begin{pmatrix}x&y&z\end{pmatrix}A\begin{pmatrix}x\\y\\z\end{pmatrix},
\]
where
\[
A=\begin{pmatrix}2&1&1\\1&2&1\\1&1&2\end{pmatrix}=I+J
\quad\text{with }J=\begin{pmatrix}1&1&1\\1&1&1\\1&1&1\end{pmatrix}.
\]
Since $J$ has eigenvalues $3,0,0$, the eigenvalues of $A=I+J$ are $4,1,1$, hence $A$ is positive definite and so is $q$.

The associated inner product is
\[
\langle u,v\rangle := u^TAv.
\]
An orthonormal basis (with respect to $\langle\cdot,\cdot\rangle$) is obtained from eigenvectors:
\[
e_1=\frac{(1,1,1)}{2\sqrt3},\qquad
e_2=\frac{(1,-1,0)}{\sqrt2},\qquad
e_3=\frac{(1,1,-2)}{\sqrt6},
\]
because $Ae_1=4e_1$ and $Ae_2=e_2$, $Ae_3=e_3$, and these vectors are mutually orthogonal.

Now apply Gram--Schmidt to the standard basis $a_1=(1,0,0)$, $a_2=(0,1,0)$, $a_3=(0,0,1)$ with respect to $\langle\cdot,\cdot\rangle$.

First,
\[
\|a_1\|^2=\langle a_1,a_1\rangle=2,\qquad
u_1=\frac{a_1}{\|a_1\|}=\frac1{\sqrt2}(1,0,0).
\]
Next,
\[
\langle a_2,u_1\rangle = a_2^TAu_1=\frac1{\sqrt2},
\quad
v_2=a_2-\langle a_2,u_1\rangle u_1=(0,1,0)-\frac12(1,0,0)=\Bigl(-\frac12,1,0\Bigr).
\]
Compute $\|v_2\|^2=\langle v_2,v_2\rangle=\frac32$, hence
\[
u_2=\frac{v_2}{\|v_2\|}=\sqrt{\frac23}\Bigl(-\frac12,1,0\Bigr).
\]
Finally,
\[
\langle a_3,u_1\rangle=\frac1{\sqrt2},\qquad
\langle a_3,u_2\rangle=\frac1{\sqrt6},
\]
so
\[
v_3=a_3-\langle a_3,u_1\rangle u_1-\langle a_3,u_2\rangle u_2
=(0,0,1)-\Bigl(\frac12,0,0\Bigr)-\Bigl(-\frac16,\frac13,0\Bigr)=\Bigl(-\frac13,-\frac13,1\Bigr).
\]
Then $\|v_3\|^2=\langle v_3,v_3\rangle=\frac43$, so
\[
u_3=\frac{v_3}{\|v_3\|}=\frac{\sqrt3}{2}\Bigl(-\frac13,-\frac13,1\Bigr).
\]
Thus $(u_1,u_2,u_3)$ is the Gram--Schmidt orthonormal basis for $\langle\cdot,\cdot\rangle$.
\end{solution}

\begin{problem}[CP-I-0115]
\label{prob:cp-i-0115}
\textbf{Orthogonal Matrices and Completing an Orthonormal Frame}

\textbf{Canonical authored completion of the retained orthogonal-matrix fragment.}

\hspace*{1.5em}Prove that for a real square matrix \(Q\), the following are equivalent:
\[
Q^TQ=I,\qquad Q^{-1}=Q^T,\qquad
\text{the columns of }Q\text{ form an orthonormal basis}.
\]
In \(\mathbb R^3\), suppose \(v_1,v_2\) are orthonormal. Show that the only unit vectors
orthogonal to both are
\[
v_3=\pm(v_1\times v_2).
\]
Explain how the sign determines the orientation.
\end{problem}

\begin{solution}
The \((i,j)\)-entry of \(Q^TQ\) is the dot product of the \(i\)-th and \(j\)-th columns of \(Q\).
Thus
\[
Q^TQ=I
\]
holds exactly when those columns are orthonormal. For a square matrix this also implies
invertibility, and multiplying \(Q^TQ=I\) by \(Q^{-1}\) gives
\[
Q^{-1}=Q^T.
\]
Conversely, \(Q^{-1}=Q^T\) immediately yields \(Q^TQ=I\).

In \(\mathbb R^3\), \(v_1\times v_2\) is perpendicular to both \(v_1\) and \(v_2\), and
\[
\|v_1\times v_2\|=\|v_1\|\|v_2\|\sin\frac{\pi}{2}=1.
\]
The common orthogonal complement of \(\operatorname{span}\{v_1,v_2\}\) is one-dimensional,
so its only unit vectors are
\[
\pm(v_1\times v_2).
\]
The \(+\) choice gives a right-handed frame and determinant \(+1\); the \(-\) choice gives a
left-handed frame and determinant \(-1\).
\end{solution}

\begin{problem}[CP-I-0116]
\label{prob:cp-i-0116}
\textbf{Quadratic Forms and Principal Axes}

\textbf{Canonical authored completion of the retained principal-axes fragment.}

\hspace*{1.5em}Let
\[
q(x)=x^TAx,\qquad A=A^T\in M_n(\mathbb R).
\]
Prove that there is an orthogonal change of coordinates \(x=Py\) for which
\[
q(x)=\lambda_1y_1^2+\cdots+\lambda_ny_n^2.
\]
Explain why this removes all cross-terms and identify the principal axes. For \(n=2\), relate
this change of coordinates to an ordinary rotation of the plane.
\end{problem}

\begin{solution}
By the spectral theorem, there is an orthogonal matrix \(P\) whose columns are an orthonormal
eigenbasis of \(A\), so
\[
P^TAP=D=\operatorname{diag}(\lambda_1,\dots,\lambda_n).
\]
Set \(x=Py\). Then
\[
q(x)=x^TAx
=y^TP^TAPy
=y^TDy
=\sum_{i=1}^{n}\lambda_i y_i^2.
\]
Since \(D\) is diagonal, no terms \(y_iy_j\) with \(i\ne j\) occur. The coordinate directions in
the \(y\)-system correspond to the eigenvectors of \(A\); these are the principal axes.

When \(n=2\), an orientation-preserving orthogonal matrix has the form
\[
P=
\begin{pmatrix}
\cos\theta&-\sin\theta\\
\sin\theta&\cos\theta
\end{pmatrix},
\]
so principal-axis reduction is precisely a rotation of coordinates.
\end{solution}
